\begin{filecontents*}[overwrite]{\jobname.ist}
headings_flag 1
numhead_positive "数字"
symhead_positive "符号"
heading_prefix "{\\bfseries\\hfil "
heading_suffix "\\hfil}\\nopagebreak\\par\\vskip 5pt "
\end{filecontents*}



\usepackage{amsmath,amssymb,amsthm}
\usepackage{mathtools}
\usepackage{mathrsfs}          %
\usepackage{tikz}              %
\usepackage{tikz-cd}           %
\usepackage{etoolbox}          %
\usepackage[a4paper,margin=2.8cm]{geometry}
\usepackage{enumitem}
\setlist{itemsep=2pt,topsep=4pt}
\emergencystretch=1.5em %
\vfuzz=3pt              %
\hfuzz=4pt              %
\hbadness=2500           %
\xeCJKsetup{PunctStyle=quanjiao}%
\usepackage{array}             %
\usepackage{multicol}          %


\newif\ifzhnames
\zhnamestrue
\newcommand{\trn}[2]{\ifzhnames #2\else #1\fi}

\newcommand{\trGrothendieck}{\trn{Grothendieck}{格罗滕迪克}}%
\newcommand{\trSerre}{\trn{Serre}{塞尔}}%
\newcommand{\trZariski}{\trn{Zariski}{扎里斯基}}%
\newcommand{\trNoether}{\trn{Noether}{诺特}}%
\newcommand{\trArtin}{\trn{Artin}{阿廷}}%
\newcommand{\trDedekind}{\trn{Dedekind}{戴德金}}%
\newcommand{\trKrull}{\trn{Krull}{克鲁尔}}%
\newcommand{\trHilbert}{\trn{Hilbert}{希尔伯特}}%
\newcommand{\trRiemann}{\trn{Riemann}{黎曼}}%
\newcommand{\trRoch}{\trn{Roch}{罗赫}}%
\newcommand{\trChevalley}{\trn{Chevalley}{谢瓦莱}}%
\newcommand{\trCartier}{\trn{Cartier}{卡蒂耶}}%
\newcommand{\trWeil}{\trn{Weil}{韦伊}}%
\newcommand{\trBertini}{\trn{Bertini}{贝尔蒂尼}}%
\newcommand{\trMumford}{\trn{Mumford}{芒福德}}%
\newcommand{\trDeligne}{\trn{Deligne}{德利涅}}%
\newcommand{\trSegre}{\trn{Segre}{塞格雷}}%
\newcommand{\trCastelnuovo}{\trn{Castelnuovo}{卡斯泰尔诺沃}}%
\newcommand{\trChow}{\trn{Chow}{周炜良}}%
\newcommand{\trNakai}{\trn{Nakai}{中井}}%
\newcommand{\trMoishezon}{\trn{Moishezon}{莫伊舍宗}}%
\newcommand{\trKleiman}{\trn{Kleiman}{克莱曼}}%
\newcommand{\trKlein}{\trn{Klein}{克莱因}}%
\newcommand{\trLeibniz}{\trn{Leibniz}{莱布尼茨}}%
\newcommand{\trDwork}{\trn{Dwork}{德沃克}}%
\newcommand{\trWitt}{\trn{Witt}{维特}}%
\newcommand{\trMobius}{\trn{M\"obius}{莫比乌斯}}%
\newcommand{\trAtiyah}{\trn{Atiyah}{阿蒂亚}}%
\newcommand{\trHirzebruch}{\trn{Hirzebruch}{希策布鲁赫}}%
\newcommand{\trHodge}{\trn{Hodge}{霍奇}}%
\newcommand{\trKahler}{\trn{K\"ahler}{凯勒}}%
\newcommand{\trPicard}{\trn{Picard}{皮卡}}%
\newcommand{\trLefschetz}{\trn{Lefschetz}{莱夫谢茨}}%
\newcommand{\trHironaka}{\trn{Hironaka}{广中平祐}}%
\newcommand{\trNakayama}{\trn{Nakayama}{中山}}%
\newcommand{\trJordan}{\trn{Jordan}{若尔当}}%
\newcommand{\trHolder}{\trn{H\"older}{赫尔德}}%
\newcommand{\trCech}{\trn{\v{C}ech}{切赫}}%
\newcommand{\trStein}{\trn{Stein}{斯坦因}}%
\newcommand{\trTate}{\trn{Tate}{泰特}}%
\newcommand{\trRussell}{\trn{Russell}{罗素}}%
\newcommand{\trVakil}{\trn{Vakil}{瓦基尔}}%
\newcommand{\trConrad}{\trn{Conrad}{康拉德}}%
\newcommand{\trHartshorne}{\trn{Hartshorne}{哈茨霍恩}}%
\newcommand{\trEisenbud}{\trn{Eisenbud}{艾森巴德}}%
\newcommand{\trHarris}{\trn{Harris}{哈里斯}}%
\newcommand{\trKedlaya}{\trn{Kedlaya}{凯德莱亚}}%
\newcommand{\trPoonen}{\trn{Poonen}{普宁}}%
\newcommand{\trKnutson}{\trn{Knutson}{克努森}}%
\newcommand{\trTao}{\trn{Tao}{陶哲轩}}%
\newcommand{\trKisin}{\trn{Kisin}{基辛}}%
\newcommand{\trMcLarty}{\trn{McLarty}{麦克拉蒂}}%
\newcommand{\trMurfet}{\trn{Murfet}{穆费特}}%
\newcommand{\trStay}{\trn{Stay}{斯泰}}%
\newcommand{\trKovacs}{\trn{Kov\'{a}cs}{科瓦奇}}%
\newcommand{\trGillooly}{\trn{Gillooly}{吉洛伊}}%
\newcommand{\trHassett}{\trn{Hassett}{哈塞特}}%
\newcommand{\trDeJong}{\trn{de~Jong}{德容}}%
\newcommand{\trHokusai}{\trn{Hokusai}{葛饰北斋}}%
\newcommand{\trZorn}{\trn{Zorn}{佐恩}}%
\newcommand{\trYoneda}{\trn{Yoneda}{米田}}%
\newcommand{\trNietzsche}{\trn{Nietzsche}{尼采}}%
\newcommand{\trGodel}{\trn{G\"odel}{哥德尔}}%
\newcommand{\trBernays}{\trn{Bernays}{贝尔奈斯}}%

\newcommand{\trhTheorem}{定理}%
\newcommand{\trhProposition}{命题}%
\newcommand{\trhLemma}{引理}%
\newcommand{\trhCorollary}{推论}%
\newcommand{\trCorollary}{推论}%
\newcommand{\trhConjecture}{猜想}%
\newcommand{\trhDefinition}{定义}%
\newcommand{\trhExample}{例}%
\newcommand{\trhExercise}{练习}%
\newcommand{\trhNotation}{记号}%
\newcommand{\trhConvention}{约定}%
\newcommand{\trhRemark}{注记}%
\newcommand{\trhWarning}{警告}%
\newcommand{\trhCaveat}{警示}%

\newcommand{\trCategory}{范畴}%
\let\trCategories\trCategory%
\newcommand{\trObject}{对象}%
\let\trObjects\trObject%
\newcommand{\trMorphism}{态射}%
\let\trMorphisms\trMorphism%
\newcommand{\trArrow}{箭头}%
\newcommand{\trMap}{映射}%
\let\trMaps\trMap%
\newcommand{\trSource}{源}%
\newcommand{\trTarget}{目标}%
\newcommand{\trIdMorphism}{恒等态射}%
\newcommand{\trIdFunctor}{恒等函子}%
\newcommand{\trIsomorphism}{同构}%
\newcommand{\trIsomorphic}{同构}%
\newcommand{\trCanonical}{典范}%
\newcommand{\trCanonically}{典范地}%
\newcommand{\trComposition}{复合}%
\newcommand{\trFunctor}{函子}%
\newcommand{\trCovariantFunctor}{协变函子}%
\newcommand{\trContravariantFunctor}{反变函子}%
\newcommand{\trFaithful}{忠实}%
\newcommand{\trFull}{满}%
\newcommand{\trFullyFaithful}{完全忠实}%
\newcommand{\trForgetfulFunctor}{遗忘函子}%
\newcommand{\trSubcategory}{子范畴}%
\newcommand{\trAbelianCategory}{阿贝尔范畴}%
\newcommand{\trAbelianGroups}{阿贝尔群}%
\newcommand{\trGroupHom}{群同态}%
\newcommand{\trGroupoid}{群胚}%
\newcommand{\trMonoid}{幺半群}%
\newcommand{\trAutomorphism}{自同构}%
\newcommand{\trAutomorphismGroup}{自同构群}%
\newcommand{\trHomomorphism}{同态}%
\newcommand{\trKernel}{核}%
\newcommand{\trCokernel}{余核}%
\newcommand{\trUniversalProperty}{泛性质}%
\newcommand{\trLimit}{极限}%
\newcommand{\trColimit}{余极限}%
\newcommand{\trAdjoint}{伴随}%
\newcommand{\trRepresentableFunctor}{可表函子}%
\newcommand{\trSpectralSequence}{谱序列}%
\newcommand{\trConcreteCategory}{具体范畴}%
\newcommand{\trPoset}{偏序集}%
\newcommand{\trFundamentalGroup}{基本群}%
\newcommand{\trHomology}{同调}%

\newcommand{\trSet}{集合}%
\let\trSetsWord\trSet%
\newcommand{\trTopologicalSpace}{拓扑空间}%
\newcommand{\trHomeomorphism}{同胚}%
\newcommand{\trField}{域}%
\newcommand{\trRing}{环}%
\let\trRingsWord\trRing%
\newcommand{\trModule}{模}%
\let\trModules\trModule%
\newcommand{\trVectorSpace}{向量空间}%
\newcommand{\trGroup}{群}%
\let\trGroups\trGroup%
\newcommand{\trOpenSet}{开集}%
\newcommand{\trClosedSet}{闭集}%
\newcommand{\trOpenSubset}{开子集}%
\newcommand{\trClosedSubset}{闭子集}%
\newcommand{\trCompact}{紧致}%
\newcommand{\trCompactness}{紧致性}%

\newcommand{\trPresheaf}{预层}%
\newcommand{\trSheaf}{层}%
\let\trSheaves\trSheaf%
\newcommand{\trSubsheaf}{子层}%
\newcommand{\trQuotientSheaf}{商层}%
\newcommand{\trStalk}{茎}%
\newcommand{\trSheafification}{层化}%
\newcommand{\trScheme}{概形}%
\let\trSchemes\trScheme%
\newcommand{\trAffineScheme}{仿射概形}%
\newcommand{\trProjectiveScheme}{射影概形}%
\newcommand{\trQCSheaf}{拟凝聚层}%
\newcommand{\trQC}{拟凝聚}%
\newcommand{\trCoherentSheaf}{凝聚层}%
\newcommand{\trSection}{截面}%
\newcommand{\trRestrictionMap}{限制映射}%
\newcommand{\trGerm}{芽}%
\newcommand{\trEspaceEtale}{平展空间}%
\newcommand{\trIdentityAxiom}{恒等公理}%
\newcommand{\trGluingAxiom}{粘合公理}%
\newcommand{\trGluability}{可粘性}%
\newcommand{\trStructureSheaf}{结构层}%
\newcommand{\trSkyscraperSheaf}{摩天楼层}%
\newcommand{\trConstantSheaf}{常层}%
\newcommand{\trConstantPresheaf}{常预层}%
\newcommand{\trLocallyConstant}{局部常值}%
\newcommand{\trPushforward}{前推}%
\newcommand{\trInverseImage}{逆像}%
\newcommand{\trCompatible}{相容}%
\newcommand{\trCompatibleGerm}{相容芽}%
\newcommand{\trStalkLocal}{茎局部}%
\newcommand{\trExponentialSequence}{指数正合列}%
\newcommand{\trIndicatorSheaf}{指示层}%
\newcommand{\trSubpresheaf}{子预层}%
\newcommand{\trExactness}{正合性}%
\newcommand{\trLeftExactness}{左正合性}%
\newcommand{\trSheafOnBase}{基上的层}%
\newcommand{\trPresheafOnBase}{基上的预层}%
\newcommand{\trBasicOpenSet}{基本开集}%
\newcommand{\trCocycleCondition}{上闭链条件}%
\newcommand{\trContractible}{可缩}%
\newcommand{\trOpenBall}{开球}%
\newcommand{\trSupport}{支撑}%
\newcommand{\trDirectImage}{直像}%
\newcommand{\trPushPullMap}{推拉映射}%
\newcommand{\trExtensionByZero}{零扩张}%
\newcommand{\trAdjointness}{伴随性}%
\newcommand{\trInclusionMap}{包含映射}%
\newcommand{\trLocus}{轨迹}%
\newcommand{\trAbstractNonsense}{抽象废话}%
\newcommand{\trTopology}{拓扑}%
\newcommand{\trAlgebraic}{代数}%
\newcommand{\trAlgebraicFunction}{代数函数}%
\newcommand{\trGeometricSpace}{几何空间}%
\newcommand{\trHausdorff}{\trn{Hausdorff}{豪斯多夫}}%
\newcommand{\trSecondCountable}{第二可数}%
\newcommand{\trSeparatedness}{分离性}%
\newcommand{\trUnramified}{非分歧}%
\newcommand{\trSubmersion}{淹没}%
\newcommand{\trIrreducible}{不可约}%
\newcommand{\trGalois}{\trn{Galois}{伽罗瓦}}%
\newcommand{\trNullstellensatz}{零点定理}%
\newcommand{\trDistinguishedOpen}{特异开集}%
\newcommand{\trAffineLine}{仿射直线}%
\newcommand{\trLittlewood}{\trn{Littlewood}{利特尔伍德}}%
\newcommand{\trPicasso}{\trn{Picasso}{毕加索}}%
\newcommand{\trFitzGerald}{\trn{FitzGerald}{菲茨杰拉德}}%
\newcommand{\trCarroll}{\trn{Carroll}{卡罗尔}}%
\newcommand{\trSpectrum}{谱}%
\newcommand{\trEuclideanDomain}{欧几里得整环}%
\newcommand{\trNilradical}{幂零根}%
\newcommand{\trHypersurface}{超曲面}%
\newcommand{\trCoordinate}{坐标}%
\newcommand{\trConjugate}{共轭}%
\newcommand{\trResidueRing}{剩余环}%
\newcommand{\trFieldExtension}{域扩张}%
\newcommand{\trFiniteFieldExtension}{有限扩张}%
\newcommand{\trDualNumbers}{对偶数}%
\newcommand{\trDivisionAlgorithm}{带余除法}%
\newcommand{\trRationalFunction}{有理函数}%
\newcommand{\trRegularFunction}{正则函数}%
\newcommand{\trPole}{极点}%
\newcommand{\trSophieGermain}{\trn{Sophie Germain}{苏菲·热尔曼}}%
\newcommand{\trEuclid}{\trn{Euclid}{欧几里得}}%
\newcommand{\trVanishingSet}{消失集}%
\newcommand{\trZeroSet}{零点集}%
\newcommand{\trRadical}{根}%
\newcommand{\trRadicalIdeal}{根理想}%
\newcommand{\trClosure}{闭包}%
\newcommand{\trZariskiClosure}{\trZariski\trClosure}%
\newcommand{\trCoarse}{粗}%
\newcommand{\trCofiniteTopology}{余有限拓扑}%
\newcommand{\trSubspaceTopology}{子空间拓扑}%
\newcommand{\trDense}{稠密}%
\newcommand{\trForm}{型}%
\newcommand{\trComplement}{补集}%
\newcommand{\trMinimalPrimeIdeal}{极小素理想}%
\newcommand{\trMinimalPrime}{极小素}%
\newcommand{\trIrreducibleComponent}{不可约分量}%
\newcommand{\trSpecialization}{特化}%
\newcommand{\trGenerization}{泛化}%
\newcommand{\trGeneralPoint}{泛点}%
\newcommand{\trConnectedComponent}{连通分量}%
\newcommand{\trDescendingChainCondition}{降链条件}%
\newcommand{\trAscendingChainCondition}{升链条件}%
\newcommand{\trIdempotent}{幂等元}%
\newcommand{\trCone}{锥}%
\newcommand{\trUltrafilter}{超滤子}%
\newcommand{\trQuasicompactness}{拟紧性}%
\newcommand{\trVeronese}{\trn{Veronese}{韦罗内塞}}%
\newcommand{\trDistinguishedBase}{特异基底}%
\newcommand{\trEmbeddedPoint}{嵌入点}%
\newcommand{\trPartitionOfUnity}{单位分解}%
\newcommand{\trAffineCommunicationLemma}{仿射通信引理}%
\newcommand{\trQuadricSurface}{二次曲面}%
\newcommand{\trFuzz}{绒毛}%
\newcommand{\trAssociatedPoint}{相伴点}%
\newcommand{\trClosedSubscheme}{闭子概形}%
\newcommand{\trMultiplicity}{重数}%
\newcommand{\trPoincare}{\trn{Poincaré}{庞加莱}}%
\newcommand{\trOpenSubscheme}{开子概形}%
\newcommand{\trRank}{秩}%
\newcommand{\trCharacteristic}{特征}%
\newcommand{\trRealValued}{实值}%
\newcommand{\trHartogs}{\trn{Hartogs}{哈托格斯}}%
\newcommand{\trWeyl}{\trn{Weyl}{外尔}}%
\newcommand{\trChineseRemainderTheorem}{中国剩余定理}%
\newcommand{\trPrincipalIdeal}{主理想}%
\newcommand{\trQuotientTopology}{商拓扑}%
\newcommand{\trHyperplane}{超平面}%
\newcommand{\trRuling}{直纹}%
\newcommand{\trProjectiveLine}{射影直线}%
\newcommand{\trGradedIdeal}{分次理想}%
\newcommand{\trGrading}{分次}%
\newcommand{\trBaseRing}{基环}%
\newcommand{\trIrrelevantIdeal}{无关理想}%
\newcommand{\trSymmetricAlgebra}{对称代数}%
\newcommand{\trQuasiprojective}{拟射影}%
\newcommand{\trVanishingScheme}{消失概形}%
\newcommand{\trProjectivization}{射影化}%
\newcommand{\trLineAtInfinity}{无穷远线}%
\newcommand{\trPlaneAtInfinity}{无穷远平面}%
\newcommand{\trConic}{二次曲线}%
\newcommand{\trProperness}{固有性}%
\newcommand{\trReducedness}{既约性}%
\newcommand{\trIntegrality}{整性}%
\newcommand{\trAffineLocal}{仿射局部}%
\newcommand{\trSubscheme}{子概形}%
\newcommand{\trNormality}{正规性}%
\newcommand{\trDegree}{次数}%
\newcommand{\trQuasiprojectivity}{拟射影性}%
\newcommand{\trIntegrallyClosed}{整闭}%
\newcommand{\trIntegralClosedness}{整闭性}%
\newcommand{\trIntegralClosure}{整闭包}%
\newcommand{\trFactorial}{唯一分解}%
\newcommand{\trFactoriality}{唯一分解性}%
\newcommand{\trIdealOfDenominators}{分母理想}%
\newcommand{\trQuadraticForm}{二次型}%
\newcommand{\trEllipticCurve}{椭圆曲线}%
\newcommand{\trClassGroup}{类群}%
\newcommand{\trSquareFree}{无平方因子}%
\newcommand{\trSeparable}{可分}%
\newcommand{\trHeight}{高度}%
\newcommand{\trFullRank}{满秩}%
\newcommand{\trMaximalRank}{极大秩}%
\newcommand{\trRegularity}{正则性}%
\newcommand{\trGauss}{\trn{Gauss}{高斯}}%
\newcommand{\trNagata}{\trn{Nagata}{永田}}%
\newcommand{\trPolynomialRing}{多项式环}%
\newcommand{\trDiagonalize}{对角化}%
\newcommand{\trFaithfullyFlat}{忠实平坦}%
\newcommand{\trDescent}{下降}%
\newcommand{\trTorsion}{挠}%
\newcommand{\trTorsionSubmodule}{挠子模}%
\newcommand{\trDistinguishedTopology}{特异拓扑}%
\newcommand{\trFreeSheaf}{自由层}%
\newcommand{\trComplexAnalyticGeometry}{复解析几何}%
\newcommand{\trDistinguishedAffineBase}{特异仿射基底}%
\newcommand{\trEtale}{平展}%
\newcommand{\trSite}{景}%
\newcommand{\trTopos}{拓扑斯}%
\newcommand{\trNonArchimedean}{非阿基米德}%
\newcommand{\trTotallyDisconnected}{全不连通}%
\newcommand{\trVoevodsky}{\trn{Voevodsky}{沃沃斯基}}%
\newcommand{\trBaseField}{基域}%
\newcommand{\trOka}{\trn{Oka}{冈洁}}%
\newcommand{\trCyclicModule}{循环模}%
\newcommand{\trCompositionSeries}{合成列}%
\newcommand{\trSimple}{单}%
\newcommand{\trLength}{长度}%
\newcommand{\trFiniteLength}{有限长度}%
\newcommand{\trFiltration}{滤过}%
\newcommand{\trNormalSubgroup}{正规子群}%
\newcommand{\trSecondIsomorphismTheorem}{第二同构定理}%
\newcommand{\trDescendingInduction}{递降归纳}%
\newcommand{\trIsomorphismClass}{同构类}%
\newcommand{\trRefine}{加细}%
\newcommand{\trAnnihilate}{零化}%
\newcommand{\trValuationRing}{赋值环}%
\newcommand{\trRestrictedPowerSeries}{受限幂级数}%
\newcommand{\trAlgebraicClosure}{代数闭包}%
\newcommand{\trAssociatedPrimeIdeal}{相伴素理想}%
\newcommand{\trEmbeddedPrime}{嵌入素}%
\newcommand{\trAnnihilatorIdeal}{零化理想}%
\newcommand{\trPrimary}{准素}%
\newcommand{\trPrimaryDecomposition}{准素分解}%
\newcommand{\trTotalFractionRing}{全分式环}%
\newcommand{\trTotalQuotientRing}{全商环}%
\newcommand{\trDomainOfDefinition}{定义域}%
\newcommand{\trIndeterminacyLocus}{不定性轨迹}%
\newcommand{\trSchemeTheoreticallyDense}{概形论稠密}%
\newcommand{\trCompleteIntersection}{完全相交}%
\newcommand{\trRegularSequence}{正则序列}%
\newcommand{\trLasker}{\trn{Lasker}{拉斯克}}%
\newcommand{\trLocallyRingedSpace}{局部赋环空间}%
\newcommand{\trStructureMorphism}{结构态射}%
\newcommand{\trValuedPoint}{值点}%
\newcommand{\trAnalytification}{解析化}%
\newcommand{\trPrevariety}{预簇}%
\newcommand{\trRationalPoint}{有理点}%
\newcommand{\trTrivialBundle}{平凡丛}%
\newcommand{\trTangentVector}{切向量}%
\newcommand{\trGradedRingMap}{分次环映射}%
\newcommand{\trVeroneseSubring}{\trVeronese{}子环}%
\newcommand{\trVeroneseEmbedding}{\trVeronese{}嵌入}%
\newcommand{\trMonomial}{单项式}%
\newcommand{\trRegrading}{重新分次}%
\newcommand{\trRationalMap}{有理映射}%
\newcommand{\trDominant}{支配}%
\newcommand{\trBirational}{双有理}%
\newcommand{\trRational}{有理}%
\newcommand{\trFermat}{\trn{Fermat}{费马}}%
\newcommand{\trCremona}{\trn{Cremona}{克雷莫纳}}%
\newcommand{\trGroupObject}{群对象}%
\newcommand{\trMultiplicativeGroup}{乘法群}%
\newcommand{\trRingScheme}{环概形}%
\newcommand{\trGroupSchemeAction}{群概形作用}%
\newcommand{\trGeneralLinearGroup}{一般线性群}%
\newcommand{\trDeterminant}{行列式}%
\newcommand{\trHopf}{\trn{Hopf}{霍普夫}}%
\newcommand{\trHopfAlgebra}{\trHopf{}代数}%
\newcommand{\trComultiplication}{余乘法}%
\newcommand{\trFlagVariety}{旗簇}%
\newcommand{\trPartialFlagVariety}{部分旗簇}%
\newcommand{\trSubspace}{子空间}%
\newcommand{\trStratification}{分层}%
\newcommand{\trNormalForm}{范式}%
\newcommand{\trRowReducedEchelonForm}{行约简阶梯形}%
\newcommand{\trSchubertCell}{舒伯特胞腔}%
\newcommand{\trSchubert}{\trn{Schubert}{舒伯特}}%
\newcommand{\trPlucker}{\trn{Plücker}{普吕克}}%
\newcommand{\trPluckerEmbedding}{\trPlucker{}嵌入}%
\newcommand{\trIntegralExtension}{整扩张}%
\newcommand{\trCancellationTheorem}{消去定理}%
\newcommand{\trJacobson}{\trn{Jacobson}{雅各布森}}%
\newcommand{\trJacobsonRadical}{\trJacobson{}根}%
\newcommand{\trMonicPolynomial}{首一多项式}%
\newcommand{\trAdjugateMatrix}{伴随矩阵}%
\newcommand{\trSubalgebra}{子代数}%
\newcommand{\trClosed}{闭}%
\newcommand{\trConstructible}{可构造}%
\newcommand{\trFundamentalTheoremOfEliminationTheory}{消元理论基本定理}%
\newcommand{\trResultant}{结式}%
\newcommand{\trEliminationOfQuantifiers}{量词消去}%
\newcommand{\trQuasifinite}{拟有限}%
\newcommand{\trTrueDegree}{真次数}%
\newcommand{\trFibralLocus}{纤维轨迹}%
\newcommand{\trRealClosedField}{实闭域}%
\newcommand{\trTarski}{\trn{Tarski}{塔斯基}}%
\newcommand{\trFinite}{有限}%
\newcommand{\trEndomorphism}{自同态}%
\newcommand{\trZariskiMainTheorem}{\trZariski{}主定理}%
\newcommand{\trBranchedCover}{分歧覆盖}%
\newcommand{\trNode}{结点}%
\newcommand{\trTautologicalMap}{重言映射}%
\newcommand{\trFiniteToOne}{有限对一}%
\newcommand{\trStandingHypothesis}{常设假设}%
\newcommand{\trProjectiveSystem}{投射系统}%
\newcommand{\trOpen}{开}%
\newcommand{\trIdealSheaf}{理想层}%
\newcommand{\trSchemeTheoretic}{概形论}%
\newcommand{\trEnrichment}{充实}%
\newcommand{\trTwistedCubic}{扭曲三次曲线}%
\newcommand{\trRationalNormalCurve}{有理正规曲线}%
\newcommand{\trLinearSpace}{线性空间}%
\newcommand{\trWeightedProjectiveSpace}{加权射影空间}%
\newcommand{\trBezout}{\trn{Bézout}{贝祖}}%
\newcommand{\trMinor}{子式}%
\newcommand{\trProjectiveCompletion}{射影完备化}%
\newcommand{\trEffectiveCartierDivisor}{有效\trCartier{}除子}%
\newcommand{\trRegularEmbedding}{正则嵌入}%
\newcommand{\trLocallyPrincipal}{局部主}%
\newcommand{\trSlicing}{切片}%
\newcommand{\trReduction}{既约化}%
\newcommand{\trTaylor}{\trn{Taylor}{泰勒}}%
\newcommand{\trSchemeTheoreticImage}{概形论像}%
\newcommand{\trZariskiSheaf}{\trZariski{}层}%
\newcommand{\trSubfunctor}{子函子}%
\newcommand{\trTorsor}{主齐性空间}%
\newcommand{\trProfiniteTopology}{投射有限拓扑}%
\newcommand{\trExtension}{扩张}%
\newcommand{\trGeometricPoint}{几何点}%
\newcommand{\trGeometricFiber}{几何纤维}%
\newcommand{\trGeometricallyConnected}{几何连通}%
\newcommand{\trGeometricallyIrreducible}{几何不可约}%
\newcommand{\trGeometricallyIntegral}{几何整}%
\newcommand{\trGeometricallyReduced}{几何既约}%
\newcommand{\trUniversallyInjective}{泛单射}%
\newcommand{\trUniversallyClosed}{泛闭}%
\newcommand{\trUniversallyOpen}{泛开}%
\newcommand{\trRadicial}{根性}%
\newcommand{\trPurelyInseparable}{纯不可分}%
\newcommand{\trSeparablyClosed}{可分闭}%
\newcommand{\trPerfectClosure}{完全闭包}%
\newcommand{\trPerfectField}{完全域}%
\newcommand{\trSeparablyGenerated}{可分生成}%
\newcommand{\trTranscendenceBase}{超越基}%
\newcommand{\trSeparatingTranscendenceBase}{分离超越基}%
\newcommand{\trInseparableDegree}{不可分次数}%
\newcommand{\trMinimalPolynomial}{极小多项式}%
\newcommand{\trLocallyClosedEmbedding}{局部闭嵌入}%
\newcommand{\trReasonable}{合理}%
\newcommand{\trFamily}{族}%
\newcommand{\trBaseScheme}{基概形}%
\newcommand{\trSchemeTheoreticPreimage}{概形论原像}%
\newcommand{\trSchemeTheoreticFiber}{概形论纤维}%
\newcommand{\trGenericFiber}{一般纤维}%
\newcommand{\trGeneralFiber}{泛纤维}%
\newcommand{\trGenericallyFinite}{一般有限}%
\newcommand{\trExceptionalDivisor}{例外除子}%
\newcommand{\trBaseChangeDiagram}{基变换图}%
\newcommand{\trChangeOfBase}{换基}%
\newcommand{\trParabola}{抛物线}%
\newcommand{\trStandardOpenSet}{标准开集}%
\newcommand{\trNumberField}{数域}%
\newcommand{\trRingOfIntegers}{整数环}%
\newcommand{\trCusp}{尖点}%
\newcommand{\trTacnode}{双切点}%
\newcommand{\trCoordinateFree}{无坐标}%
\newcommand{\trRealPoint}{实点}%
\newcommand{\trCharacter}{特征标}%
\newcommand{\trNormallyClosed}{正规闭}%
\newcommand{\trPinchedPlane}{捏合平面}%
\newcommand{\trGraphMorphism}{图像态射}%
\newcommand{\trSubvariety}{子簇}%
\newcommand{\trCancellationProperty}{消去性质}%
\newcommand{\trDiagonalSubscheme}{对角子概形}%
\newcommand{\trAffineDiagonal}{仿射对角}%
\newcommand{\trIrreducibility}{不可约性}%
\newcommand{\trFundamentalPoint}{基本点}%
\newcommand{\trGroupVariety}{群簇}%
\newcommand{\trAlgebraicGroup}{代数群}%
\newcommand{\trAbelianVariety}{\trAbelian{}簇}%
\newcommand{\trLie}{\trn{Lie}{李}}%
\newcommand{\trComplexTorus}{复环面}%
\newcommand{\trLattice}{格}%
\newcommand{\trSimplyConnected}{单连通}%
\newcommand{\trUniversalCover}{泛覆盖}%
\newcommand{\trTranslation}{平移}%
\newcommand{\trIdentityPoint}{单位点}%
\newcommand{\trComplete}{完备}%
\newcommand{\trInversionMorphism}{求逆态射}%
\newcommand{\trAbel}{\trn{Abel}{阿贝尔}}%
\newcommand{\trClosedMap}{闭映射}%
\newcommand{\trCoh}{凝聚}%
\newcommand{\trZariskiTopology}{\trZariski{}拓扑}%
\newcommand{\trGenericPoint}{一般点}%
\newcommand{\trClosedPoint}{闭点}%
\newcommand{\trNilpotent}{幂零}%
\newcommand{\trReduced}{既约}%
\newcommand{\trIntegral}{整}%
\newcommand{\trNormal}{正规}%
\newcommand{\trRegular}{正则}%
\newcommand{\trSmooth}{光滑}%
\newcommand{\trDimension}{维数}%
\newcommand{\trCodimension}{余维数}%
\newcommand{\trPureDimensional}{纯维}%
\newcommand{\trPureDimension}{纯维数}%
\newcommand{\trEquidimensional}{等维}%
\newcommand{\trNFold}{$n$ 维体}%
\newcommand{\trNoetherNormalization}{\trNoether{}正规化}%
\newcommand{\trAlgebraicallyIndependent}{代数无关}%
\newcommand{\trAlgebraicallyDependent}{代数相关}%
\newcommand{\trCompositum}{合成域}%
\newcommand{\trPurelyTranscendental}{纯超越}%
\newcommand{\trCatenary}{链式}%
\newcommand{\trSemicontinuity}{半连续性}%
\newcommand{\trUpperSemicontinuous}{上半连续}%
\newcommand{\trLowerSemicontinuous}{下半连续}%
\newcommand{\trIncidenceCorrespondence}{关联对应}%
\newcommand{\trIncidenceVariety}{关联簇}%
\newcommand{\trNormalClosure}{正规闭包}%
\newcommand{\trJacobianMatrix}{雅可比矩阵}%
\newcommand{\trVeryGeneral}{非常泛}%
\newcommand{\trPrincipalIdealTheorem}{主理想定理}%
\newcommand{\trSystemOfParameters}{参数系}%
\newcommand{\trSymbolicPower}{符号幂}%
\newcommand{\trUniruled}{单有理}%
\newcommand{\trPrimeElement}{素元}%
\newcommand{\trCotangentVector}{余切向量}%
\newcommand{\trDerivation}{导子}%
\newcommand{\trTangentLine}{切线}%
\newcommand{\trCorank}{余秩}%
\newcommand{\trAssociatedGraded}{相伴分次}%
\newcommand{\trJacobianCriterion}{雅可比判别法}%
\newcommand{\trNonsingular}{非奇异}%
\newcommand{\trSingular}{奇异}%
\newcommand{\trSingularPoint}{奇点}%
\newcommand{\trMixedCharacteristic}{混合特征}%
\newcommand{\trEisenstein}{\trn{Eisenstein}{艾森斯坦}}%
\newcommand{\trRees}{\trn{Rees}{里斯}}%
\newcommand{\trPuiseux}{\trn{Puiseux}{皮瑟}}%
\newcommand{\trTangentPlane}{切平面}%
\newcommand{\trProjection}{投射}%
\newcommand{\trDualVariety}{对偶簇}%
\newcommand{\trDualProjectiveSpace}{对偶射影空间}%
\newcommand{\trUniversalHyperplane}{万有超平面}%
\newcommand{\trUniformizer}{单值化元}%
\newcommand{\trValuation}{赋值}%
\newcommand{\trZeros}{零点}%
\newcommand{\trDedekindDomain}{\trDedekind{}整环}%
\newcommand{\trTransitionFunction}{转移函数}%
\newcommand{\trTrivialization}{平凡化}%
\newcommand{\trTrivial}{平凡}%
\newcommand{\trLocallyFree}{局部自由}%
\newcommand{\trFractionalIdeal}{分式理想}%
\newcommand{\trTensorAlgebra}{张量代数}%
\newcommand{\trExteriorAlgebra}{外代数}%
\newcommand{\trPerfectPairing}{完美配对}%
\newcommand{\trReflexiveSheaf}{自反层}%
\newcommand{\trTraceMap}{迹映射}%
\newcommand{\trSaturation}{饱和化}%
\newcommand{\trBasePoint}{基点}%
\newcommand{\trBaseLocus}{基点轨迹}%
\newcommand{\trBasePointFree}{无基点}%
\newcommand{\trLinearSystem}{线性系}%
\newcommand{\trDegenerate}{退化}%
\newcommand{\trBoxProduct}{盒积}%
\newcommand{\trSecant}{割线}%
\newcommand{\trPrincipal}{主}%
\newcommand{\trExcision}{切除}%
\newcommand{\trQuillen}{\trn{Quillen}{奎伦}}%
\newcommand{\trSuslin}{\trn{Suslin}{苏斯林}}%
\newcommand{\trVaserstein}{\trn{Vaserstein}{瓦瑟斯坦}}%
\newcommand{\trProjectiveTransformation}{射影变换}%
\newcommand{\trFractionalLinearTransformation}{分式线性变换}%
\newcommand{\trStabilizer}{稳定子群}%
\newcommand{\trNormalBundle}{正规丛}%
\newcommand{\trSaturated}{饱和}%
\newcommand{\trGloballyGenerated}{整体生成的}%
\newcommand{\trAmpleness}{丰富性}%
\newcommand{\trSyzygy}{合冲}%
\newcommand{\trGlobalGeneration}{整体生成}%
\newcommand{\trKodaira}{\trn{Kodaira}{小平}}%
\newcommand{\trSemiample}{半丰富}%
\newcommand{\trVeryAmpleness}{极丰富性}%
\newcommand{\trHermitianMetric}{厄米特度量}%
\newcommand{\trSubidealSheaf}{子理想层}%
\newcommand{\trPrincipalModule}{主模}%
\newcommand{\trTautologicalQuotientBundle}{重言商丛}%
\newcommand{\trTautologicalSubbundle}{重言子丛}%
\newcommand{\trUniversalExactSequence}{泛正合列}%
\newcommand{\trGrassmannianBundle}{格拉斯曼丛}%
\newcommand{\trQuadricHypersurface}{二次超曲面}%
\newcommand{\trPureTensor}{纯张量}%
\newcommand{\trSheafOfAlgebras}{代数层}%
\newcommand{\trSheafOfGradedAlgebras}{分次代数层}%
\newcommand{\trTotalSpace}{全空间}%
\newcommand{\trTautologicalBundle}{重言丛}%
\newcommand{\trHirzebruchSurface}{希策布鲁赫曲面}%
\newcommand{\trQuasiaffine}{拟仿射}%
\newcommand{\trQuasiisomorphism}{拟同构}%
\newcommand{\trDerivedCategory}{导出范畴}%
\newcommand{\trDualizingSheaf}{对偶化层}%
\newcommand{\trMultidegree}{多重次数}%
\newcommand{\trReducedHomology}{约化同调}%
\newcommand{\trSimplex}{单形}%
\newcommand{\trLaurentPolynomial}{Laurent 多项式}%
\newcommand{\trEulerCharacteristic}{欧拉示性数}%
\newcommand{\trArithmeticGenus}{算术亏格}%
\newcommand{\trNef}{数值有效}%
\newcommand{\trNumericallyTrivial}{数值平凡}%
\newcommand{\trNumericallyEquivalent}{数值等价}%
\newcommand{\trChern}{\trn{Chern}{陈}}%
\newcommand{\trWeber}{\trn{Weber}{韦伯}}%
\newcommand{\trAffineness}{仿射性}%
\newcommand{\trDesideratum}{愿望}%
\newcommand{\trAlgebraicVolumeForm}{代数体积形式}%
\newcommand{\trZeroSheaf}{零层}%
\newcommand{\trHyperplaneClass}{超平面类}%
\newcommand{\trTopologicalGenus}{拓扑亏格}%
\newcommand{\trHigherPushforward}{高阶前推}%
\newcommand{\trHigherDirectImage}{高阶直像}%
\newcommand{\trGeometricallyRegular}{几何正则}%
\newcommand{\trHyperelliptic}{超椭圆}%
\newcommand{\trBranchPoint}{分支点}%
\newcommand{\trRamification}{分歧}%
\newcommand{\trDoubleCover}{二重覆盖}%
\newcommand{\trInvolution}{对合}%
\newcommand{\trhClaim}{论断}%
\newcommand{\trPlaneQuartic}{平面四次曲线}%
\newcommand{\trTrigonal}{三角}%
\newcommand{\trCrossRatio}{交比}%
\newcommand{\trFlex}{拐}%
\newcommand{\trWeierstrassNormalForm}{魏尔斯特拉斯范式}%
\newcommand{\trModularForm}{模形式}%
\newcommand{\trModularCurve}{模曲线}%
\newcommand{\trDoublePlane}{二重平面}%
\newcommand{\trGroupLaw}{群律}%
\newcommand{\trFaltings}{\trn{Faltings}{法尔廷斯}}%
\newcommand{\trMordell}{\trn{Mordell}{莫德尔}}%
\newcommand{\trTaniyama}{\trn{Taniyama}{谷山}}%
\newcommand{\trShimura}{\trn{Shimura}{志村}}%
\newcommand{\trHurwitz}{\trn{Hurwitz}{赫尔维茨}}%
\newcommand{\trWeierstrass}{\trn{Weierstrass}{魏尔斯特拉斯}}%
\newcommand{\trLuroth}{\trn{L\"uroth}{吕罗特}}%
\newcommand{\trPappus}{\trn{Pappus}{帕普斯}}%
\newcommand{\trPascal}{\trn{Pascal}{帕斯卡}}%
\newcommand{\trChasles}{\trn{Chasles}{沙勒}}%
\newcommand{\trCayleyBacharach}{\trn{Cayley--Bacharach}{凯莱--巴赫拉赫}}%
\newcommand{\trViete}{\trn{Vi\`ete}{韦达}}%
\newcommand{\trPoncelet}{\trn{Poncelet}{彭赛列}}%
\newcommand{\trDelPezzo}{\trn{del Pezzo}{德尔·佩佐}}%
\newcommand{\trIntersectionProduct}{相交积}%
\newcommand{\trIntersectionNumber}{相交数}%
\newcommand{\trIntersection}{相交}%
\newcommand{\trFibration}{纤维化}%
\newcommand{\trProperTransform}{真变换}%
\newcommand{\trSelfIntersection}{自相交}%
\newcommand{\trIntersectionForm}{相交型}%
\newcommand{\trAdjunctionFormula}{伴随公式}%
\newcommand{\trSignature}{符号差}%
\newcommand{\trMovingLemma}{移动引理}%
\newcommand{\trConiveauFiltration}{余水平滤过}%
\newcommand{\trEffective}{有效}%
\newcommand{\trBorel}{\trn{Borel}{博雷尔}}%
\newcommand{\trMoore}{\trn{Moore}{穆尔}}%
\newcommand{\trCotangentSheaf}{余切层}%
\newcommand{\trTangentSheaf}{切层}%
\newcommand{\trConormalSheaf}{余正规模}%
\newcommand{\trConormal}{余正规}%
\newcommand{\trGeometricGenus}{几何亏格}%
\newcommand{\trHomogeneousSpace}{齐性空间}%
\newcommand{\trRamificationLocus}{分歧轨迹}%
\newcommand{\trBranchLocus}{分支轨迹}%
\newcommand{\trRamificationDivisor}{分歧除子}%
\newcommand{\trBranchDivisor}{分支除子}%
\newcommand{\trRamificationIndex}{分歧指数}%
\newcommand{\trTame}{驯}%
\newcommand{\trWild}{野}%
\newcommand{\trUnbranched}{无分支}%
\newcommand{\trGenericallySeparable}{一般可分}%
\newcommand{\trDifferentIdeal}{差积理想}%
\newcommand{\trDiscriminantIdeal}{判别式理想}%
\newcommand{\trDiscriminant}{判别式}%
\newcommand{\trSmallResolution}{小解消}%
\newcommand{\trTubularNeighborhood}{管状邻域}%
\newcommand{\trRaynaud}{\trn{Raynaud}{雷诺}}%
\newcommand{\trKummer}{\trn{Kummer}{库默尔}}%
\newcommand{\trEnriques}{\trn{Enriques}{恩里克斯}}%
\newcommand{\trSchreier}{\trn{Schreier}{施赖埃尔}}%
\newcommand{\trAndre}{\trn{Andr\'e}{安德烈}}%
\newcommand{\trFiber}{纤维}%
\newcommand{\trFiberProduct}{纤维积}%
\newcommand{\trBaseChange}{基变换}%
\newcommand{\trProper}{固有}%
\newcommand{\trSeparated}{分离}%
\newcommand{\trFiniteType}{有限型}%
\newcommand{\trLocally}{局部}%
\newcommand{\trNeighborhood}{邻域}%
\newcommand{\trCover}{覆盖}%
\newcommand{\trOpenCover}{开覆盖}%
\newcommand{\trDiscreteTopology}{离散拓扑}%
\newcommand{\trClassicalTopology}{经典拓扑}%
\newcommand{\trTopologicalGroup}{拓扑群}%
\newcommand{\trBounded}{有界}%
\newcommand{\trContinuous}{连续}%
\newcommand{\trOpenImmersion}{开浸入}%
\newcommand{\trClosedImmersion}{闭嵌入}%
\newcommand{\trImmersion}{浸入}%
\newcommand{\trOpenEmbedding}{开嵌入}%
\newcommand{\trClosedEmbedding}{闭嵌入}%
\newcommand{\trAmple}{丰富}%
\newcommand{\trVeryAmple}{极丰富}%
\newcommand{\trDivisor}{除子}%
\newcommand{\trCartierDivisor}{\trCartier{}除子}%
\newcommand{\trWeilDivisor}{\trWeil{}除子}%
\newcommand{\trEffectiveDivisor}{有效除子}%
\newcommand{\trInvertibleSheaf}{可逆层}%
\newcommand{\trLineBundle}{线丛}%
\newcommand{\trVectorBundle}{向量丛}%
\newcommand{\trTangentBundle}{切丛}%
\newcommand{\trLocallyFreeSheaf}{局部自由层}%
\newcommand{\trGlobalSections}{整体截面}%
\newcommand{\trCohomology}{上同调}%
\newcommand{\trCechCohomology}{\trCech{}上同调}%
\newcommand{\trDerivedFunctor}{导出函子}%
\newcommand{\trFlat}{平坦}%
\newcommand{\trTorsionFree}{无挠}%
\newcommand{\trNormalization}{正规化}%
\newcommand{\trBlowup}{胀开}%
\newcommand{\trModuliSpace}{模空间}%
\newcommand{\trGrassmannian}{格拉斯曼流形}%
\newcommand{\trAlgClosed}{代数闭}%
\newcommand{\trAlgClosedField}{代数闭域}%
\newcommand{\trVariety}{簇}%
\newcommand{\trCurve}{曲线}%
\newcommand{\trSurface}{曲面}%
\newcommand{\trPicardGroup}{\trPicard{}群}%
\let\trNoetherian\trNoether%
\let\trArtinian\trArtin%
\newcommand{\trLocalRing}{局部环}%
\newcommand{\trLocalization}{局部化}%
\newcommand{\trPrimeIdeal}{素理想}%
\newcommand{\trMaximalIdeal}{极大理想}%
\newcommand{\trIdeal}{理想}%
\newcommand{\trCotangentSpace}{余切空间}%
\newcommand{\trResidueField}{剩余域}%
\newcommand{\trFunctionField}{函数域}%
\newcommand{\trDVR}{离散赋值环}%
\newcommand{\trTrDeg}{超越次数}%
\newcommand{\trKrullDimension}{\trKrull{}维数}%
\newcommand{\trDepth}{深度}%
\newcommand{\trCM}{\trCohen{}--\trMacaulay}%
\newcommand{\trDifferential}{微分}%
\newcommand{\trKahlerDifferentials}{\trKahler{}微分}%
\newcommand{\trCotangentBundle}{余切丛}%
\newcommand{\trTangentSpace}{切空间}%
\newcommand{\trIntersectionTheory}{相交理论}%
\newcommand{\trGradedRing}{分次环}%
\newcommand{\trGradedModule}{分次模}%
\newcommand{\trHomogeneous}{齐次}%
\newcommand{\trGroupScheme}{群概形}%
\newcommand{\trGenus}{亏格}%
\newcommand{\trRiemannRoch}{\trRiemann{}--\trRoch}%
\newcommand{\trSerreDuality}{\trSerre{}对偶}%
\newcommand{\trDuality}{对偶}%

\newcommand{\trLocallySmallCategory}{局部小范畴}%
\newcommand{\trEqualizer}{等子}%
\newcommand{\trClass}{类}%
\newcommand{\trFunctoriality}{函子性}%
\newcommand{\trPath}{道路}%
\newcommand{\trHomotopy}{同伦}%
\newcommand{\trConnected}{连通}%
\newcommand{\trLinearTransformation}{线性变换}%
\newcommand{\trInvertible}{可逆}%
\newcommand{\trSetTheoretic}{集合论}%
\newcommand{\trParadox}{悖论}%

\newcommand{\trAlgebraicGeometry}{代数几何}%
\newcommand{\trSchemeTheory}{\trScheme{}理论}%
\newcommand{\trCanon}{正典}%
\newcommand{\trAffine}{仿射}%
\newcommand{\trComplex}{复}%
\newcommand{\trReal}{实}%
\newcommand{\trManifold}{流形}%
\newcommand{\trSmoothness}{光滑性}%
\newcommand{\trDifferentialGeometry}{微分几何}%
\newcommand{\trComplexGeometry}{\trComplex{}几何}%
\newcommand{\trAnalysis}{分析}%
\newcommand{\trNumberTheory}{数论}%
\newcommand{\trAlgebraicTopology}{代数拓扑}%
\newcommand{\trHomologicalAlgebra}{同调代数}%
\newcommand{\trCoordinatePatch}{坐标片}%
\newcommand{\trRingedSpace}{赋环空间}%
\newcommand{\trComplexAnalyticSpace}{复解析空间}%
\newcommand{\trAlgebraicSpace}{代数空间}%
\newcommand{\trOrbifold}{轨形}%
\newcommand{\trStack}{叠}%
\newcommand{\trRigidAnalyticSpace}{刚性解析空间}%
\newcommand{\trFormalScheme}{形式概形}%
\newcommand{\trYonedaLemma}{\trYoneda{}引理}%
\newcommand{\trProduct}{积}%
\newcommand{\trBase}{基底}%
\newcommand{\trVarietyTheory}{簇论}%
\newcommand{\trRiemannSurface}{\trRiemann{}曲面}%
\newcommand{\trCombinatorialAlgGeom}{组合代数几何}%
\newcommand{\trComputationalAlgGeom}{计算代数几何}%
\newcommand{\trHolomorphic}{全纯}%
\newcommand{\trAnalytic}{解析}%
\newcommand{\trDifferentiable}{可微}%
\newcommand{\trUniverse}{宇宙}%
\newcommand{\trQuasicompact}{拟紧}%
\newcommand{\trQuasiSeparated}{拟分离}%
\newcommand{\trQuasiseparatedness}{拟分离性}%
\newcommand{\trValuativeCriteria}{赋值判别法}%
\newcommand{\trCommutative}{交换}%
\newcommand{\trUnit}{单位元}%
\newcommand{\trZeroRing}{零环}%
\newcommand{\trIntegralDomain}{整环}%
\newcommand{\trFinalObject}{终对象}%
\newcommand{\trInitialObject}{初始对象}%
\newcommand{\trAxiomOfChoice}{选择公理}%
\newcommand{\trZornsLemma}{\trZorn{}引理}%
\newcommand{\trProperIdeal}{真理想}%
\newcommand{\trInjective}{内射}%
\newcommand{\trProjectiveObj}{投射}%
\newcommand{\trProjective}{射影}%
\newcommand{\trProjectiveSpace}{射影空间}%
\newcommand{\trPID}{主理想整环}%
\newcommand{\trUFD}{唯一分解整环}%
\newcommand{\trTensorProduct}{张量积}%
\newcommand{\trExactSequence}{正合列}%
\newcommand{\trDirectSum}{直和}%
\newcommand{\trBijection}{双射}%
\let\trBijections\trBijection%
\newcommand{\trInjectiveMap}{单射}%
\let\trInjection\trInjectiveMap%
\newcommand{\trSurjectiveMap}{满射}%
\newcommand{\trPullback}{拉回}%
\newcommand{\trFullSubcategory}{满子范畴}%
\newcommand{\trOppositeCategory}{反范畴}%
\newcommand{\trNaturalTransformation}{自然变换}%
\newcommand{\trNaturalIsomorphism}{自然同构}%
\newcommand{\trEquivalenceOfCategories}{范畴的等价}%
\newcommand{\trFunctorOfPoints}{点函子}%
\newcommand{\trUnderlyingSet}{底集合}%
\newcommand{\trContinuousMap}{连续映射}%
\newcommand{\trFiniteDimensional}{有限维}%
\newcommand{\trDual}{对偶}%
\newcommand{\trSmallCategory}{小范畴}%
\newcommand{\trCommutativeDiagram}{交换图}%
\newcommand{\trZeroObject}{零对象}%
\newcommand{\trZerodivisor}{零因子}%
\newcommand{\trNonZerodivisor}{非零因子}%
\newcommand{\trFinitelyGenerated}{有限生成}%
\newcommand{\trFractionField}{分式域}%
\newcommand{\trMultiplicativeSubset}{乘性子集}%
\newcommand{\trFreeModule}{自由模}%
\newcommand{\trSubmodule}{子模}%
\newcommand{\trBilinear}{双线性}%
\newcommand{\trRightExact}{右正合}%
\newcommand{\trLeftExact}{左正合}%
\newcommand{\trRingMorphism}{环映射}%
\newcommand{\trAlgebra}{代数}%
\let\trAlgebras\trAlgebra%
\newcommand{\trSubring}{子环}%
\newcommand{\trWellDefined}{良定}%
\newcommand{\trEquivalenceRelation}{等价关系}%
\newcommand{\trEquivalenceClass}{等价类}%
\let\trEquivalenceClasses\trEquivalenceClass%
\newcommand{\trBinaryOperation}{二元运算}%
\newcommand{\trFinitelyPresented}{有限表现}%
\newcommand{\trQuotient}{商}%
\newcommand{\trDiagonalMorphism}{对角态射}%
\newcommand{\trCartesianDiagram}{笛卡尔方图}%
\newcommand{\trCoproduct}{余积}%
\newcommand{\trFiberedCoproduct}{纤维余积}%
\newcommand{\trDisjointUnion}{不交并}%
\newcommand{\trMonomorphism}{单态射}%
\newcommand{\trEpimorphism}{满态射}%
\newcommand{\trDivisibleGroup}{可除群}%
\newcommand{\trRepresentable}{可表}%
\newcommand{\trFunctorCategory}{函子范畴}%
\newcommand{\trYonedaEmbedding}{\trYoneda{}嵌入}%
\newcommand{\trUniversalObject}{泛对象}%
\newcommand{\trYoda}{\trn{Yoda}{尤达}}%
\let\trNonZerodivisors\trNonZerodivisor%
\newcommand{\trDiagram}{图}%
\newcommand{\trIndexCategory}{指标范畴}%
\newcommand{\trIndexSet}{指标集}%
\newcommand{\trInverseLimit}{逆极限}%
\newcommand{\trProjectiveLimit}{投射极限}%
\newcommand{\trDirectLimit}{正向极限}%
\newcommand{\trInductiveLimit}{归纳极限}%
\newcommand{\trInjectiveLimit}{内射极限}%
\newcommand{\trFiltered}{滤过}%
\newcommand{\trDirected}{有向}%
\newcommand{\trCompletion}{完备化}%
\newcommand{\trFormalPowerSeries}{形式幂级数}%
\newcommand{\trConvergentPowerSeries}{收敛幂级数}%
\newcommand{\trPadic}{$p$~进}%
\newcommand{\trPadicIntegers}{\trPadic{}整数}%
\newcommand{\trPadics}{$p$~进数}%
\newcommand{\trComathematician}{余数学家}%
\newcommand{\trCoauthor}{余作者}%
\newcommand{\trAdjointPair}{伴随对}%
\newcommand{\trLeftAdjoint}{左伴随}%
\newcommand{\trRightAdjoint}{右伴随}%
\newcommand{\trNatural}{自然}%
\newcommand{\trAdjunctionUnit}{单位}%
\newcommand{\trAdjunctionCounit}{余单位}%
\newcommand{\trRepresentation}{表示}%
\newcommand{\trRepresentationTheory}{表示论}%
\newcommand{\trFiniteGroup}{有限群}%
\newcommand{\trSubgroup}{子群}%
\newcommand{\trInduction}{诱导}%
\newcommand{\trRestriction}{限制}%
\newcommand{\trFrobenius}{\trn{Frobenius}{弗罗贝尼乌斯}}%
\newcommand{\trFrobeniusReciprocity}{\trFrobenius{}互反律}%
\newcommand{\trSuspensionFunctor}{纬垂函子}%
\newcommand{\trLoopSpace}{环路空间}%
\newcommand{\trSemigroup}{半群}%
\newcommand{\trAbelian}{阿贝尔}%
\newcommand{\trAbelianSemigroup}{\trAbelian\trSemigroup}%
\newcommand{\trGroupification}{群化}%
\newcommand{\trInverseElement}{逆元}%
\newcommand{\trExtensionOfScalars}{标量扩张}%
\newcommand{\trRestrictionOfScalars}{标量限制}%
\newcommand{\trForgetful}{遗忘}%
\newcommand{\trAdditive}{加性}%
\newcommand{\trAdditiveCategory}{\trAdditive{}范畴}%
\newcommand{\trAdditiveFunctor}{\trAdditive{}函子}%
\newcommand{\trExact}{正合}%
\newcommand{\trZeroMorphism}{零态射}%
\newcommand{\trSubobject}{子对象}%
\newcommand{\trQuotientObject}{商对象}%
\newcommand{\trImage}{像}%
\newcommand{\trComplexN}{复形}%
\let\trComplexes\trComplexN%
\newcommand{\trHomotopic}{同伦的}%
\newcommand{\trConnectingHomomorphism}{连接同态}%
\newcommand{\trCycle}{闭链}%
\let\trCycles\trCycle%
\newcommand{\trBoundary}{边缘链}%
\let\trBoundaries\trBoundary%
\newcommand{\trLongExactSequence}{长正合列}%
\newcommand{\trShortExactSequence}{短正合列}%
\newcommand{\trFlatness}{平坦性}%
\newcommand{\trCommutativeMonoid}{\trCommutative{}幺半群}%
\newcommand{\trCodiagonal}{余对角}%
\newcommand{\trGenerator}{生成元}%
\let\trGenerators\trGenerator%
\newcommand{\trRelation}{关系}%
\newcommand{\trBanach}{\trn{Banach}{巴拿赫}}%
\newcommand{\trBanachSpace}{\trBanach{}空间}%
\newcommand{\trFreyd}{\trn{Freyd}{弗雷德}}%
\newcommand{\trMitchell}{\trn{Mitchell}{米切尔}}%
\newcommand{\trFreydMitchellEmbeddingTheorem}{\trFreyd{}--\trMitchell{}嵌入定理}%
\newcommand{\trMacLane}{\trn{Mac~Lane}{麦克莱恩}}%
\newcommand{\trDieudonne}{\trn{Dieudonn\'e}{迪厄多内}}%
\newcommand{\trInverseSystem}{逆向系统}%
\newcommand{\trCofiltered}{余滤过}%
\newcommand{\trTransitionMap}{过渡映射}%
\newcommand{\trStalkification}{茎化}%
\newcommand{\trLeray}{\trn{Leray}{勒雷}}%
\newcommand{\trDoubleComplex}{双复形}%
\newcommand{\trTotalComplex}{全复形}%
\newcommand{\trAntidiagonal}{反对角线}%
\newcommand{\trHypercohomology}{超上同调}%
\newcommand{\trPage}{页}%
\newcommand{\trFirstQuadrant}{第一象限}%
\newcommand{\trAnticommute}{反交换}%
\newcommand{\trOrientation}{定向}%
\newcommand{\trSubquotient}{子商}%
\newcommand{\trConverge}{收敛}%
\newcommand{\trAbut}{抵达}%
\newcommand{\trCollapse}{塌缩}%
\newcommand{\trMappingCone}{映射锥}%
\newcommand{\trSpectralFunctor}{谱函子}%
\newcommand{\trTangent}{切}%
\newcommand{\trDesingularization}{去奇点}%
\newcommand{\trResolutionOfSingularities}{奇点解消}%
\newcommand{\trTotalTransform}{全变换}%
\newcommand{\trCenter}{中心}%
\newcommand{\trFiberedCube}{纤维立方体}%
\newcommand{\trMinimalModelProgram}{极小模型纲领}%
\newcommand{\trBlowupClosureLemma}{胀开闭包引理}%
\newcommand{\trRemainderTheorem}{余式定理}%
\newcommand{\trGradedAlgebra}{分次代数}%
\newcommand{\trDynkinDiagram}{邓金图}%
\newcommand{\trAlterationTheorem}{变更定理}%
\newcommand{\trWeakFactorizationTheorem}{弱分解定理}%
\newcommand{\trAtiyahFlop}{\trAtiyah{}翻转}%
\newcommand{\trEllipticFibration}{椭圆纤维化}%
\newcommand{\trIntersectionHomology}{相交同调}%
\newcommand{\trDeformationToTheNormalCone}{到正规锥的形变}%
\newcommand{\trAnalyticallyIsomorphic}{解析同构}%
\newcommand{\trInvariantTheory}{不变量论}%
\newcommand{\trSingularLocus}{奇异轨迹}%
\newcommand{\trDeltaFunctor}{$\delta$-函子}%
\newcommand{\trEffaceable}{可消去}%
\newcommand{\trAcyclic}{零调}%
\newcommand{\trCartan}{\trn{Cartan}{嘉当}}%
\newcommand{\trEilenberg}{\trn{Eilenberg}{艾伦伯格}}%
\newcommand{\trFlasque}{松软}%
\newcommand{\trAlexandrov}{\trn{Alexandrov}{亚历山德罗夫}}%
\newcommand{\trDeformation}{形变}%
\newcommand{\trDeformationTheory}{形变理论}%
\newcommand{\trOpenMap}{开映射}%
\newcommand{\trHilbertScheme}{\trHilbert{}概形}%
\newcommand{\trHilbertPolynomial}{\trHilbert{}多项式}%
\newcommand{\trTolstoy}{\trn{Tolstoy}{托尔斯泰}}%
\newcommand{\trHalfExact}{半正合}%
\newcommand{\trCayley}{\trn{Cayley}{凯莱}}%
\newcommand{\trSalmon}{\trn{Salmon}{萨蒙}}%
\newcommand{\trDonagi}{\trn{Donagi}{多纳吉}}%
\newcommand{\trSmith}{\trn{Smith}{史密斯}}%
\newcommand{\trTsen}{\trn{Tsen}{曾炯}}%
\newcommand{\trGrauert}{\trn{Grauert}{格劳尔特}}%
\newcommand{\trMiracleFlat}{奇迹平坦}%
\newcommand{\trCohen}{\trn{Cohen}{科恩}}%
\newcommand{\trMacaulay}{\trn{Macaulay}{麦考利}}%
\newcommand{\trHensel}{\trn{Hensel}{亨泽尔}}%
\newcommand{\trTriplePoint}{三重点}%
\newcommand{\trSplitNode}{分裂结点}%
\newcommand{\trNonSplitNode}{非分裂结点}%
\newcommand{\trThickening}{加厚}%
\newcommand{\trDeltaInvariant}{$\delta$-不变量}%



\newif\ifrkfont \rkfontfalse
\IfFontExistsTF{SimSun-ExtB}{%
  \setCJKfamilyfont{rk}{SimSun-ExtB}\rkfonttrue}{%
  \IfFontExistsTF{Source Han Serif CN}{%
    \setCJKfamilyfont{rk}{Source Han Serif CN}\rkfonttrue}{}}
\newcommand{\rfchar}[1]{{\ifrkfont\CJKfamily{rk}\fi #1}}

\newcounter{secthm}[section]
\renewcommand{\thesecthm}{\thesection.\arabic{secthm}}
\newcounter{secex}[section]
\renewcommand{\thesecex}{\thesection.\Alph{secex}}
\newcommand{\vkthmsync}[1]{\setcounter{secthm}{\numexpr(#1)-1\relax}}
\newcommand{\vkexsync}[1]{\setcounter{secex}{\numexpr(`#1)-65\relax}}

\newtheoremstyle{vakplain}{\topsep}{\topsep}{\itshape}{}
  {\bfseries}{}{.5em}{#1~#2\ifstrempty{#3}{}{（#3）}}
\newtheoremstyle{vakdefinition}{\topsep}{\topsep}{}{}
  {\bfseries}{}{.5em}{#1~#2\ifstrempty{#3}{}{（#3）}}
\newtheoremstyle{vakremark}{\topsep}{\topsep}{\upshape}{}
  {\itshape}{}{.5em}{#1~#2\ifstrempty{#3}{}{（#3）}}

\theoremstyle{vakplain}
\newtheorem{theorem}    [secthm]{\trhTheorem}
\newtheorem{proposition}[secthm]{\trhProposition}
\newtheorem{lemma}      [secthm]{\trhLemma}
\newtheorem{corollary}  [secthm]{\trhCorollary}
\newtheorem{conjecture} [secthm]{\trhConjecture}
\theoremstyle{vakdefinition}
\newtheorem{definition} [secthm]{\trhDefinition}
\newtheorem{example}    [secthm]{\trhExample}
\newtheorem{notation}   [secthm]{\trhNotation}
\newtheorem{convention} [secthm]{\trhConvention}
\newtheorem{exercise}   [secex] {\trhExercise}
\theoremstyle{vakremark}
\newtheorem{remark}  [secthm]{\trhRemark}
\newtheorem{warning} [secthm]{\trhWarning}
\newtheorem{caveat}  [secthm]{\trhCaveat}

\newcommand{\hard}{$\star$}
\newcommand{\extrahard}{$\star\star$}

\counterwithin{figure}{chapter}
\counterwithin{table}{chapter}

\newcommand{\epigraph}[2]{%
  \begin{center}\begin{minipage}{0.86\textwidth}
    {\kaishu #1}\par\vspace{2pt}
    \hfill{\kaishu\small #2}\par
  \end{minipage}\end{center}\vspace{6pt}
}

\newcommand{\OO}{\mathcal{O}}
\newcommand{\Spec}{\operatorname{Spec}}
\newcommand{\Proj}{\operatorname{Proj}}
\newcommand{\Mor}{\operatorname{Mor}}
\newcommand{\Hom}{\operatorname{Hom}}
\newcommand{\colim}{\operatorname{colim}}
\newcommand{\Ind}{\operatorname{Ind}}   %
\newcommand{\Res}{\operatorname{Res}}   %
\newcommand{\coker}{\operatorname{coker}} %
\newcommand{\Supp}{\operatorname{Supp}}   %
\newcommand{\Sym}{\operatorname{Sym}}     %
\newcommand{\Gal}{\operatorname{Gal}}     %
\newcommand{\vkbracedown}[5]{%
  ({#1},{#3+#5}) .. controls ({#1},{#3}) and ({(#1+#2)/2-0.4},{#3-#4}) ..
  ({(#1+#2)/2-0.09},{#3-#4+0.09}) -- ({(#1+#2)/2},{#3-#4}) --
  ({(#1+#2)/2+0.09},{#3-#4+0.09}) .. controls ({(#1+#2)/2+0.4},{#3-#4}) and
  ({#2},{#3}) .. ({#2},{#3+#5})}
\newcommand{\vkbraceup}[5]{%
  ({#1},{#3-#5}) .. controls ({#1},{#3}) and ({(#1+#2)/2-0.4},{#3+#4}) ..
  ({(#1+#2)/2-0.09},{#3+#4-0.09}) -- ({(#1+#2)/2},{#3+#4}) --
  ({(#1+#2)/2+0.09},{#3+#4-0.09}) .. controls ({(#1+#2)/2+0.4},{#3+#4}) and
  ({#2},{#3}) .. ({#2},{#3-#5})}
\newcommand{\im}{\operatorname{im}}       %
\newcommand{\obj}{\operatorname{obj}}
\newcommand{\Aut}{\operatorname{Aut}}
\newcommand{\id}{\mathrm{id}}
\newcommand{\scr}[1]{\mathscr{#1}}
\newcommand{\catSets}{\mathbf{Sets}}
\newcommand{\catVec}{\mathbf{Vec}}
\newcommand{\catTop}{\mathbf{Top}}
\newcommand{\catAb}{\mathbf{Ab}}
\newcommand{\catGrp}{\mathbf{Grp}}
\newcommand{\catRings}{\mathbf{Rings}}
\newcommand{\catSetsX}[1]{\textsc{Sets}_{#1}}
\newcommand{\catAbX}[1]{\textsc{Ab}_{#1}}
\newcommand{\catModOX}{\textsc{Mod}_{\OO_X}}
\newcommand{\catSetsPreX}[1]{\textsc{Sets}^{\mathrm{pre}}_{#1}}
\newcommand{\catAbPreX}[1]{\textsc{Ab}^{\mathrm{pre}}_{#1}}
\newcommand{\catQCo}{\mathbf{QCoh}}              %
\newcommand{\catCoh}{\mathbf{Coh}}               %
\newcommand{\catGroups}{\mathit{Groups}}         %
\newcommand{\catMod}[1]{\mathbf{Mod}_{#1}}      %
\newcommand{\fdVec}[1]{\mathbf{f.d.Vec}_{#1}}   %
\newcommand{\Ho}{\mathbf{Ho}}                   %
\newcommand{\opp}{\mathrm{opp}}                 %
\newcommand{\Eto}[2]{{}_{\to}E^{#1}_{#2}}
\newcommand{\Eup}[2]{{}_{\uparrow}E^{#1}_{#2}}

\usepackage{xcolor}            %
\usepackage{hyperref}
\hypersetup{
  colorlinks=true,
  linkcolor=blue!55!black,
  citecolor=teal,
  urlcolor=violet,
  bookmarksnumbered=true,
  pdflang={zh-CN},
  pdftitle={上升之海：代数几何基础（中文译本）},
  pdfauthor={Ravi Vakil（原著）},
}

\usepackage[capitalise,noabbrev]{cleveref}
\crefname{theorem}    {\trhTheorem}    {\trhTheorem}
\crefname{proposition}{\trhProposition}{\trhProposition}
\crefname{lemma}      {\trhLemma}      {\trhLemma}
\crefname{corollary}  {\trhCorollary}  {\trhCorollary}
\crefname{conjecture} {\trhConjecture} {\trhConjecture}
\crefname{definition} {\trhDefinition} {\trhDefinition}
\crefname{example}    {\trhExample}    {\trhExample}
\crefname{notation}   {\trhNotation}   {\trhNotation}
\crefname{convention} {\trhConvention} {\trhConvention}
\crefname{exercise}   {\trhExercise}   {\trhExercise}
\crefname{remark}     {\trhRemark}     {\trhRemark}
\crefname{warning}    {\trhWarning}    {\trhWarning}
\crefname{caveat}     {\trhCaveat}     {\trhCaveat}
\crefname{equation}   {}{}  %
\crefformat{equation}{#2(#1)#3}
\crefname{secthm}{\trhTheorem}{\trhTheorem}
\crefname{secex}{\trhExercise}{\trhExercise}
\crefname{figure}{图}{图}%
\AtBeginEnvironment{theorem}    {\crefalias{secthm}{theorem}}
\AtBeginEnvironment{proposition}{\crefalias{secthm}{proposition}}
\AtBeginEnvironment{lemma}      {\crefalias{secthm}{lemma}}
\AtBeginEnvironment{corollary}  {\crefalias{secthm}{corollary}}
\AtBeginEnvironment{conjecture} {\crefalias{secthm}{conjecture}}
\AtBeginEnvironment{definition} {\crefalias{secthm}{definition}}
\AtBeginEnvironment{example}    {\crefalias{secthm}{example}}
\AtBeginEnvironment{notation}   {\crefalias{secthm}{notation}}
\AtBeginEnvironment{convention} {\crefalias{secthm}{convention}}
\AtBeginEnvironment{remark}     {\crefalias{secthm}{remark}}
\AtBeginEnvironment{warning}    {\crefalias{secthm}{warning}}
\AtBeginEnvironment{caveat}     {\crefalias{secthm}{caveat}}
\AtBeginEnvironment{exercise}   {\crefalias{secex}{exercise}}

\newcommand{\vkglossrule}{\par\noindent\rule{\linewidth}{0.4pt}\par}
\newcommand{\vkglosshead}[1]{\par\textbf{#1}\par}
\newcommand{\vkgloss}[2]{\par\hangindent1.5em\hangafter1\noindent #1\ \emph{#2}}
\newcommand{\vkdict}[2]{\par\hangindent1.5em\hangafter1\noindent \emph{#1}\ #2}
\newcommand{\vkidxentry}[2]{\par\hangindent2em\hangafter1\noindent #1\quad #2}

\usepackage{imakeidx}
\makeindex[title={译本术语索引},intoc,noautomatic]

\usepackage{multicol}
\newcommand{\BookGlossEntry}[2]{%
  \par\noindent\hangindent=1em\hangafter=1 #1\quad #2\par}
\newcommand{\BookSourceIndexEntry}[2]{%
  \par\begingroup\leftskip=#1em\relax\hangindent=1em\relax
  \noindent #2\par\endgroup}
\makeatletter
\newif\ifBookIndexing
\BookIndexingfalse
\newcommand{\BookBoldPage}[1]{\textbf{\hyperpage{#1}}}
\DeclareRobustCommand{\BookRecordIndex}[1]{%
  \ifBookIndexing
    \expandafter\index\expandafter{\csname BookIndex@#1\endcsname|hyperpage}%
  \fi}
\DeclareRobustCommand{\BookIndexDefinition}[1]{%
  \ifBookIndexing
    \expandafter\index\expandafter{\csname BookIndex@#1\endcsname|BookBoldPage}%
  \fi}
\DeclareRobustCommand{\BookTerm}[2]{%
  \BookRecordIndex{#1}%
  \begingroup\let\BookRecordIndex\@gobble #2\endgroup}
\AddToHook{cmd/@outputpage/before}{\let\BookRecordIndex\@gobble}
\pdfstringdefDisableCommands{%
  \def\BookTerm#1#2{#2}%
  \let\BookRecordIndex\@gobble
  \let\BookIndexDefinition\@gobble}
\expandafter\def\csname BookIndex@T0001\endcsname{tuilun@推论（corollary）}
\expandafter\let\csname BookOriginal@trCorollary\endcsname\trCorollary
\renewcommand{\trCorollary}{\BookTerm{T0001}{\csname BookOriginal@trCorollary\endcsname}}
\expandafter\def\csname BookIndex@T0002\endcsname{fanchou@范畴（category）}
\expandafter\let\csname BookOriginal@trCategory\endcsname\trCategory
\renewcommand{\trCategory}{\BookTerm{T0002}{\csname BookOriginal@trCategory\endcsname}}
\expandafter\def\csname BookIndex@T0003\endcsname{duixiang@对象（object）}
\expandafter\let\csname BookOriginal@trObject\endcsname\trObject
\renewcommand{\trObject}{\BookTerm{T0003}{\csname BookOriginal@trObject\endcsname}}
\expandafter\def\csname BookIndex@T0004\endcsname{taishe@态射（morphism）}
\expandafter\let\csname BookOriginal@trMorphism\endcsname\trMorphism
\renewcommand{\trMorphism}{\BookTerm{T0004}{\csname BookOriginal@trMorphism\endcsname}}
\expandafter\def\csname BookIndex@T0005\endcsname{jiantou@箭头（arrow）}
\expandafter\let\csname BookOriginal@trArrow\endcsname\trArrow
\renewcommand{\trArrow}{\BookTerm{T0005}{\csname BookOriginal@trArrow\endcsname}}
\expandafter\def\csname BookIndex@T0006\endcsname{yingshe@映射（map）}
\expandafter\let\csname BookOriginal@trMap\endcsname\trMap
\renewcommand{\trMap}{\BookTerm{T0006}{\csname BookOriginal@trMap\endcsname}}
\expandafter\def\csname BookIndex@T0007\endcsname{yuan@源（source）}
\expandafter\let\csname BookOriginal@trSource\endcsname\trSource
\renewcommand{\trSource}{\BookTerm{T0007}{\csname BookOriginal@trSource\endcsname}}
\expandafter\def\csname BookIndex@T0008\endcsname{mubiao@目标（target）}
\expandafter\let\csname BookOriginal@trTarget\endcsname\trTarget
\renewcommand{\trTarget}{\BookTerm{T0008}{\csname BookOriginal@trTarget\endcsname}}
\expandafter\def\csname BookIndex@T0009\endcsname{hengdengtaishe@恒等态射（identity morphism）}
\expandafter\let\csname BookOriginal@trIdMorphism\endcsname\trIdMorphism
\renewcommand{\trIdMorphism}{\BookTerm{T0009}{\csname BookOriginal@trIdMorphism\endcsname}}
\expandafter\def\csname BookIndex@T0010\endcsname{hengdenghanzi@恒等函子（identity functor）}
\expandafter\let\csname BookOriginal@trIdFunctor\endcsname\trIdFunctor
\renewcommand{\trIdFunctor}{\BookTerm{T0010}{\csname BookOriginal@trIdFunctor\endcsname}}
\expandafter\def\csname BookIndex@T0011\endcsname{tonggou@同构（isomorphism /\allowbreak{} isomorphic）}
\expandafter\let\csname BookOriginal@trIsomorphism\endcsname\trIsomorphism
\renewcommand{\trIsomorphism}{\BookTerm{T0011}{\csname BookOriginal@trIsomorphism\endcsname}}
\expandafter\def\csname BookIndex@T0012\endcsname{tonggou@同构（isomorphism /\allowbreak{} isomorphic）}
\expandafter\let\csname BookOriginal@trIsomorphic\endcsname\trIsomorphic
\renewcommand{\trIsomorphic}{\BookTerm{T0012}{\csname BookOriginal@trIsomorphic\endcsname}}
\expandafter\def\csname BookIndex@T0013\endcsname{dianfan@典范（canonical）}
\expandafter\let\csname BookOriginal@trCanonical\endcsname\trCanonical
\renewcommand{\trCanonical}{\BookTerm{T0013}{\csname BookOriginal@trCanonical\endcsname}}
\expandafter\def\csname BookIndex@T0014\endcsname{dianfande@典范地（canonically）}
\expandafter\let\csname BookOriginal@trCanonically\endcsname\trCanonically
\renewcommand{\trCanonically}{\BookTerm{T0014}{\csname BookOriginal@trCanonically\endcsname}}
\expandafter\def\csname BookIndex@T0015\endcsname{fuhe@复合（composition）}
\expandafter\let\csname BookOriginal@trComposition\endcsname\trComposition
\renewcommand{\trComposition}{\BookTerm{T0015}{\csname BookOriginal@trComposition\endcsname}}
\expandafter\def\csname BookIndex@T0016\endcsname{hanzi@函子（functor）}
\expandafter\let\csname BookOriginal@trFunctor\endcsname\trFunctor
\renewcommand{\trFunctor}{\BookTerm{T0016}{\csname BookOriginal@trFunctor\endcsname}}
\expandafter\def\csname BookIndex@T0017\endcsname{xiebianhanzi@协变函子（covariant functor）}
\expandafter\let\csname BookOriginal@trCovariantFunctor\endcsname\trCovariantFunctor
\renewcommand{\trCovariantFunctor}{\BookTerm{T0017}{\csname BookOriginal@trCovariantFunctor\endcsname}}
\expandafter\def\csname BookIndex@T0018\endcsname{fanbianhanzi@反变函子（contravariant functor）}
\expandafter\let\csname BookOriginal@trContravariantFunctor\endcsname\trContravariantFunctor
\renewcommand{\trContravariantFunctor}{\BookTerm{T0018}{\csname BookOriginal@trContravariantFunctor\endcsname}}
\expandafter\def\csname BookIndex@T0019\endcsname{zhongshi@忠实（faithful）}
\expandafter\let\csname BookOriginal@trFaithful\endcsname\trFaithful
\renewcommand{\trFaithful}{\BookTerm{T0019}{\csname BookOriginal@trFaithful\endcsname}}
\expandafter\def\csname BookIndex@T0020\endcsname{man@满（full）}
\expandafter\let\csname BookOriginal@trFull\endcsname\trFull
\renewcommand{\trFull}{\BookTerm{T0020}{\csname BookOriginal@trFull\endcsname}}
\expandafter\def\csname BookIndex@T0021\endcsname{wanquanzhongshi@完全忠实（fully faithful）}
\expandafter\let\csname BookOriginal@trFullyFaithful\endcsname\trFullyFaithful
\renewcommand{\trFullyFaithful}{\BookTerm{T0021}{\csname BookOriginal@trFullyFaithful\endcsname}}
\expandafter\def\csname BookIndex@T0022\endcsname{yiwanghanzi@遗忘函子（forgetful functor）}
\expandafter\let\csname BookOriginal@trForgetfulFunctor\endcsname\trForgetfulFunctor
\renewcommand{\trForgetfulFunctor}{\BookTerm{T0022}{\csname BookOriginal@trForgetfulFunctor\endcsname}}
\expandafter\def\csname BookIndex@T0023\endcsname{zifanchou@子范畴（subcategory）}
\expandafter\let\csname BookOriginal@trSubcategory\endcsname\trSubcategory
\renewcommand{\trSubcategory}{\BookTerm{T0023}{\csname BookOriginal@trSubcategory\endcsname}}
\expandafter\def\csname BookIndex@T0024\endcsname{abeierfanchou@阿贝尔范畴（abelian category）}
\expandafter\let\csname BookOriginal@trAbelianCategory\endcsname\trAbelianCategory
\renewcommand{\trAbelianCategory}{\BookTerm{T0024}{\csname BookOriginal@trAbelianCategory\endcsname}}
\expandafter\def\csname BookIndex@T0025\endcsname{abeierqun@阿贝尔群（abelian groups）}
\expandafter\let\csname BookOriginal@trAbelianGroups\endcsname\trAbelianGroups
\renewcommand{\trAbelianGroups}{\BookTerm{T0025}{\csname BookOriginal@trAbelianGroups\endcsname}}
\expandafter\def\csname BookIndex@T0026\endcsname{quntongtai@群同态（group homomorphism）}
\expandafter\let\csname BookOriginal@trGroupHom\endcsname\trGroupHom
\renewcommand{\trGroupHom}{\BookTerm{T0026}{\csname BookOriginal@trGroupHom\endcsname}}
\expandafter\def\csname BookIndex@T0027\endcsname{qunpei@群胚（groupoid）}
\expandafter\let\csname BookOriginal@trGroupoid\endcsname\trGroupoid
\renewcommand{\trGroupoid}{\BookTerm{T0027}{\csname BookOriginal@trGroupoid\endcsname}}
\expandafter\def\csname BookIndex@T0028\endcsname{yaobanqun@幺半群（monoid）}
\expandafter\let\csname BookOriginal@trMonoid\endcsname\trMonoid
\renewcommand{\trMonoid}{\BookTerm{T0028}{\csname BookOriginal@trMonoid\endcsname}}
\expandafter\def\csname BookIndex@T0029\endcsname{zitonggou@自同构（automorphism）}
\expandafter\let\csname BookOriginal@trAutomorphism\endcsname\trAutomorphism
\renewcommand{\trAutomorphism}{\BookTerm{T0029}{\csname BookOriginal@trAutomorphism\endcsname}}
\expandafter\def\csname BookIndex@T0030\endcsname{zitonggouqun@自同构群（automorphism group）}
\expandafter\let\csname BookOriginal@trAutomorphismGroup\endcsname\trAutomorphismGroup
\renewcommand{\trAutomorphismGroup}{\BookTerm{T0030}{\csname BookOriginal@trAutomorphismGroup\endcsname}}
\expandafter\def\csname BookIndex@T0031\endcsname{tongtai@同态（homomorphism）}
\expandafter\let\csname BookOriginal@trHomomorphism\endcsname\trHomomorphism
\renewcommand{\trHomomorphism}{\BookTerm{T0031}{\csname BookOriginal@trHomomorphism\endcsname}}
\expandafter\def\csname BookIndex@T0032\endcsname{he@核（kernel）}
\expandafter\let\csname BookOriginal@trKernel\endcsname\trKernel
\renewcommand{\trKernel}{\BookTerm{T0032}{\csname BookOriginal@trKernel\endcsname}}
\expandafter\def\csname BookIndex@T0033\endcsname{yuhe@余核（cokernel）}
\expandafter\let\csname BookOriginal@trCokernel\endcsname\trCokernel
\renewcommand{\trCokernel}{\BookTerm{T0033}{\csname BookOriginal@trCokernel\endcsname}}
\expandafter\def\csname BookIndex@T0034\endcsname{fanxingzhi@泛性质（universal property）}
\expandafter\let\csname BookOriginal@trUniversalProperty\endcsname\trUniversalProperty
\renewcommand{\trUniversalProperty}{\BookTerm{T0034}{\csname BookOriginal@trUniversalProperty\endcsname}}
\expandafter\def\csname BookIndex@T0035\endcsname{jixian@极限（limit）}
\expandafter\let\csname BookOriginal@trLimit\endcsname\trLimit
\renewcommand{\trLimit}{\BookTerm{T0035}{\csname BookOriginal@trLimit\endcsname}}
\expandafter\def\csname BookIndex@T0036\endcsname{yujixian@余极限（colimit）}
\expandafter\let\csname BookOriginal@trColimit\endcsname\trColimit
\renewcommand{\trColimit}{\BookTerm{T0036}{\csname BookOriginal@trColimit\endcsname}}
\expandafter\def\csname BookIndex@T0037\endcsname{bansui@伴随（adjoint）}
\expandafter\let\csname BookOriginal@trAdjoint\endcsname\trAdjoint
\renewcommand{\trAdjoint}{\BookTerm{T0037}{\csname BookOriginal@trAdjoint\endcsname}}
\expandafter\def\csname BookIndex@T0038\endcsname{kebiaohanzi@可表函子（representable functor）}
\expandafter\let\csname BookOriginal@trRepresentableFunctor\endcsname\trRepresentableFunctor
\renewcommand{\trRepresentableFunctor}{\BookTerm{T0038}{\csname BookOriginal@trRepresentableFunctor\endcsname}}
\expandafter\def\csname BookIndex@T0039\endcsname{puxulie@谱序列（spectral sequence）}
\expandafter\let\csname BookOriginal@trSpectralSequence\endcsname\trSpectralSequence
\renewcommand{\trSpectralSequence}{\BookTerm{T0039}{\csname BookOriginal@trSpectralSequence\endcsname}}
\expandafter\def\csname BookIndex@T0040\endcsname{jutifanchou@具体范畴（concrete category）}
\expandafter\let\csname BookOriginal@trConcreteCategory\endcsname\trConcreteCategory
\renewcommand{\trConcreteCategory}{\BookTerm{T0040}{\csname BookOriginal@trConcreteCategory\endcsname}}
\expandafter\def\csname BookIndex@T0041\endcsname{pianxuji@偏序集（partially ordered set）}
\expandafter\let\csname BookOriginal@trPoset\endcsname\trPoset
\renewcommand{\trPoset}{\BookTerm{T0041}{\csname BookOriginal@trPoset\endcsname}}
\expandafter\def\csname BookIndex@T0042\endcsname{jibenqun@基本群（fundamental group）}
\expandafter\let\csname BookOriginal@trFundamentalGroup\endcsname\trFundamentalGroup
\renewcommand{\trFundamentalGroup}{\BookTerm{T0042}{\csname BookOriginal@trFundamentalGroup\endcsname}}
\expandafter\def\csname BookIndex@T0043\endcsname{tongdiao@同调（homology）}
\expandafter\let\csname BookOriginal@trHomology\endcsname\trHomology
\renewcommand{\trHomology}{\BookTerm{T0043}{\csname BookOriginal@trHomology\endcsname}}
\expandafter\def\csname BookIndex@T0044\endcsname{jihe@集合（set）}
\expandafter\let\csname BookOriginal@trSet\endcsname\trSet
\renewcommand{\trSet}{\BookTerm{T0044}{\csname BookOriginal@trSet\endcsname}}
\expandafter\def\csname BookIndex@T0045\endcsname{tapukongjian@拓扑空间（topological space）}
\expandafter\let\csname BookOriginal@trTopologicalSpace\endcsname\trTopologicalSpace
\renewcommand{\trTopologicalSpace}{\BookTerm{T0045}{\csname BookOriginal@trTopologicalSpace\endcsname}}
\expandafter\def\csname BookIndex@T0046\endcsname{tongpei@同胚（homeomorphism）}
\expandafter\let\csname BookOriginal@trHomeomorphism\endcsname\trHomeomorphism
\renewcommand{\trHomeomorphism}{\BookTerm{T0046}{\csname BookOriginal@trHomeomorphism\endcsname}}
\expandafter\def\csname BookIndex@T0047\endcsname{yu@域（field）}
\expandafter\let\csname BookOriginal@trField\endcsname\trField
\renewcommand{\trField}{\BookTerm{T0047}{\csname BookOriginal@trField\endcsname}}
\expandafter\def\csname BookIndex@T0048\endcsname{huan@环（ring）}
\expandafter\let\csname BookOriginal@trRing\endcsname\trRing
\renewcommand{\trRing}{\BookTerm{T0048}{\csname BookOriginal@trRing\endcsname}}
\expandafter\def\csname BookIndex@T0049\endcsname{mo@模（module）}
\expandafter\let\csname BookOriginal@trModule\endcsname\trModule
\renewcommand{\trModule}{\BookTerm{T0049}{\csname BookOriginal@trModule\endcsname}}
\expandafter\def\csname BookIndex@T0050\endcsname{xiangliangkongjian@向量空间（vector space）}
\expandafter\let\csname BookOriginal@trVectorSpace\endcsname\trVectorSpace
\renewcommand{\trVectorSpace}{\BookTerm{T0050}{\csname BookOriginal@trVectorSpace\endcsname}}
\expandafter\def\csname BookIndex@T0051\endcsname{qun@群（group）}
\expandafter\let\csname BookOriginal@trGroup\endcsname\trGroup
\renewcommand{\trGroup}{\BookTerm{T0051}{\csname BookOriginal@trGroup\endcsname}}
\expandafter\def\csname BookIndex@T0052\endcsname{kaiji@开集（open set）}
\expandafter\let\csname BookOriginal@trOpenSet\endcsname\trOpenSet
\renewcommand{\trOpenSet}{\BookTerm{T0052}{\csname BookOriginal@trOpenSet\endcsname}}
\expandafter\def\csname BookIndex@T0053\endcsname{biji@闭集（closed set）}
\expandafter\let\csname BookOriginal@trClosedSet\endcsname\trClosedSet
\renewcommand{\trClosedSet}{\BookTerm{T0053}{\csname BookOriginal@trClosedSet\endcsname}}
\expandafter\def\csname BookIndex@T0054\endcsname{kaiziji@开子集（open subset）}
\expandafter\let\csname BookOriginal@trOpenSubset\endcsname\trOpenSubset
\renewcommand{\trOpenSubset}{\BookTerm{T0054}{\csname BookOriginal@trOpenSubset\endcsname}}
\expandafter\def\csname BookIndex@T0055\endcsname{biziji@闭子集（closed subset）}
\expandafter\let\csname BookOriginal@trClosedSubset\endcsname\trClosedSubset
\renewcommand{\trClosedSubset}{\BookTerm{T0055}{\csname BookOriginal@trClosedSubset\endcsname}}
\expandafter\def\csname BookIndex@T0056\endcsname{jinzhi@紧致（compact）}
\expandafter\let\csname BookOriginal@trCompact\endcsname\trCompact
\renewcommand{\trCompact}{\BookTerm{T0056}{\csname BookOriginal@trCompact\endcsname}}
\expandafter\def\csname BookIndex@T0057\endcsname{jinzhixing@紧致性（compactness）}
\expandafter\let\csname BookOriginal@trCompactness\endcsname\trCompactness
\renewcommand{\trCompactness}{\BookTerm{T0057}{\csname BookOriginal@trCompactness\endcsname}}
\expandafter\def\csname BookIndex@T0058\endcsname{yuceng@预层（presheaf）}
\expandafter\let\csname BookOriginal@trPresheaf\endcsname\trPresheaf
\renewcommand{\trPresheaf}{\BookTerm{T0058}{\csname BookOriginal@trPresheaf\endcsname}}
\expandafter\def\csname BookIndex@T0059\endcsname{ceng@层（sheaf）}
\expandafter\let\csname BookOriginal@trSheaf\endcsname\trSheaf
\renewcommand{\trSheaf}{\BookTerm{T0059}{\csname BookOriginal@trSheaf\endcsname}}
\expandafter\def\csname BookIndex@T0060\endcsname{ziceng@子层（subsheaf）}
\expandafter\let\csname BookOriginal@trSubsheaf\endcsname\trSubsheaf
\renewcommand{\trSubsheaf}{\BookTerm{T0060}{\csname BookOriginal@trSubsheaf\endcsname}}
\expandafter\def\csname BookIndex@T0061\endcsname{shangceng@商层（quotient sheaf）}
\expandafter\let\csname BookOriginal@trQuotientSheaf\endcsname\trQuotientSheaf
\renewcommand{\trQuotientSheaf}{\BookTerm{T0061}{\csname BookOriginal@trQuotientSheaf\endcsname}}
\expandafter\def\csname BookIndex@T0062\endcsname{jing@茎（stalk）}
\expandafter\let\csname BookOriginal@trStalk\endcsname\trStalk
\renewcommand{\trStalk}{\BookTerm{T0062}{\csname BookOriginal@trStalk\endcsname}}
\expandafter\def\csname BookIndex@T0063\endcsname{cenghua@层化（sheafification）}
\expandafter\let\csname BookOriginal@trSheafification\endcsname\trSheafification
\renewcommand{\trSheafification}{\BookTerm{T0063}{\csname BookOriginal@trSheafification\endcsname}}
\expandafter\def\csname BookIndex@T0064\endcsname{gaixing@概形（scheme）}
\expandafter\let\csname BookOriginal@trScheme\endcsname\trScheme
\renewcommand{\trScheme}{\BookTerm{T0064}{\csname BookOriginal@trScheme\endcsname}}
\expandafter\def\csname BookIndex@T0065\endcsname{fangshegaixing@仿射概形（affine scheme）}
\expandafter\let\csname BookOriginal@trAffineScheme\endcsname\trAffineScheme
\renewcommand{\trAffineScheme}{\BookTerm{T0065}{\csname BookOriginal@trAffineScheme\endcsname}}
\expandafter\def\csname BookIndex@T0066\endcsname{sheyinggaixing@射影概形（projective scheme）}
\expandafter\let\csname BookOriginal@trProjectiveScheme\endcsname\trProjectiveScheme
\renewcommand{\trProjectiveScheme}{\BookTerm{T0066}{\csname BookOriginal@trProjectiveScheme\endcsname}}
\expandafter\def\csname BookIndex@T0067\endcsname{niningjuceng@拟凝聚层（quasicoherent sheaf）}
\expandafter\let\csname BookOriginal@trQCSheaf\endcsname\trQCSheaf
\renewcommand{\trQCSheaf}{\BookTerm{T0067}{\csname BookOriginal@trQCSheaf\endcsname}}
\expandafter\def\csname BookIndex@T0068\endcsname{niningju@拟凝聚（quasicoherent）}
\expandafter\let\csname BookOriginal@trQC\endcsname\trQC
\renewcommand{\trQC}{\BookTerm{T0068}{\csname BookOriginal@trQC\endcsname}}
\expandafter\def\csname BookIndex@T0069\endcsname{ningjuceng@凝聚层（coherent sheaf）}
\expandafter\let\csname BookOriginal@trCoherentSheaf\endcsname\trCoherentSheaf
\renewcommand{\trCoherentSheaf}{\BookTerm{T0069}{\csname BookOriginal@trCoherentSheaf\endcsname}}
\expandafter\def\csname BookIndex@T0070\endcsname{jiemian@截面（section）}
\expandafter\let\csname BookOriginal@trSection\endcsname\trSection
\renewcommand{\trSection}{\BookTerm{T0070}{\csname BookOriginal@trSection\endcsname}}
\expandafter\def\csname BookIndex@T0071\endcsname{xianzhiyingshe@限制映射（restriction map）}
\expandafter\let\csname BookOriginal@trRestrictionMap\endcsname\trRestrictionMap
\renewcommand{\trRestrictionMap}{\BookTerm{T0071}{\csname BookOriginal@trRestrictionMap\endcsname}}
\expandafter\def\csname BookIndex@T0072\endcsname{ya@芽（germ）}
\expandafter\let\csname BookOriginal@trGerm\endcsname\trGerm
\renewcommand{\trGerm}{\BookTerm{T0072}{\csname BookOriginal@trGerm\endcsname}}
\expandafter\def\csname BookIndex@T0073\endcsname{pingzhankongjian@平展空间（espace etale）}
\expandafter\let\csname BookOriginal@trEspaceEtale\endcsname\trEspaceEtale
\renewcommand{\trEspaceEtale}{\BookTerm{T0073}{\csname BookOriginal@trEspaceEtale\endcsname}}
\expandafter\def\csname BookIndex@T0074\endcsname{hengdenggongli@恒等公理（identity axiom）}
\expandafter\let\csname BookOriginal@trIdentityAxiom\endcsname\trIdentityAxiom
\renewcommand{\trIdentityAxiom}{\BookTerm{T0074}{\csname BookOriginal@trIdentityAxiom\endcsname}}
\expandafter\def\csname BookIndex@T0075\endcsname{zhanhegongli@粘合公理（gluing axiom）}
\expandafter\let\csname BookOriginal@trGluingAxiom\endcsname\trGluingAxiom
\renewcommand{\trGluingAxiom}{\BookTerm{T0075}{\csname BookOriginal@trGluingAxiom\endcsname}}
\expandafter\def\csname BookIndex@T0076\endcsname{kezhanxing@可粘性（gluability）}
\expandafter\let\csname BookOriginal@trGluability\endcsname\trGluability
\renewcommand{\trGluability}{\BookTerm{T0076}{\csname BookOriginal@trGluability\endcsname}}
\expandafter\def\csname BookIndex@T0077\endcsname{jiegouceng@结构层（structure sheaf）}
\expandafter\let\csname BookOriginal@trStructureSheaf\endcsname\trStructureSheaf
\renewcommand{\trStructureSheaf}{\BookTerm{T0077}{\csname BookOriginal@trStructureSheaf\endcsname}}
\expandafter\def\csname BookIndex@T0078\endcsname{motianlouceng@摩天楼层（skyscraper sheaf）}
\expandafter\let\csname BookOriginal@trSkyscraperSheaf\endcsname\trSkyscraperSheaf
\renewcommand{\trSkyscraperSheaf}{\BookTerm{T0078}{\csname BookOriginal@trSkyscraperSheaf\endcsname}}
\expandafter\def\csname BookIndex@T0079\endcsname{changceng@常层（constant sheaf）}
\expandafter\let\csname BookOriginal@trConstantSheaf\endcsname\trConstantSheaf
\renewcommand{\trConstantSheaf}{\BookTerm{T0079}{\csname BookOriginal@trConstantSheaf\endcsname}}
\expandafter\def\csname BookIndex@T0080\endcsname{changyuceng@常预层（constant presheaf）}
\expandafter\let\csname BookOriginal@trConstantPresheaf\endcsname\trConstantPresheaf
\renewcommand{\trConstantPresheaf}{\BookTerm{T0080}{\csname BookOriginal@trConstantPresheaf\endcsname}}
\expandafter\def\csname BookIndex@T0081\endcsname{jubuchangzhi@局部常值（locally constant）}
\expandafter\let\csname BookOriginal@trLocallyConstant\endcsname\trLocallyConstant
\renewcommand{\trLocallyConstant}{\BookTerm{T0081}{\csname BookOriginal@trLocallyConstant\endcsname}}
\expandafter\def\csname BookIndex@T0082\endcsname{qiantui@前推（pushforward /\allowbreak{} direct image）}
\expandafter\let\csname BookOriginal@trPushforward\endcsname\trPushforward
\renewcommand{\trPushforward}{\BookTerm{T0082}{\csname BookOriginal@trPushforward\endcsname}}
\expandafter\def\csname BookIndex@T0083\endcsname{nixiang@逆像（inverse image）}
\expandafter\let\csname BookOriginal@trInverseImage\endcsname\trInverseImage
\renewcommand{\trInverseImage}{\BookTerm{T0083}{\csname BookOriginal@trInverseImage\endcsname}}
\expandafter\def\csname BookIndex@T0084\endcsname{xiangrong@相容（compatible）}
\expandafter\let\csname BookOriginal@trCompatible\endcsname\trCompatible
\renewcommand{\trCompatible}{\BookTerm{T0084}{\csname BookOriginal@trCompatible\endcsname}}
\expandafter\def\csname BookIndex@T0085\endcsname{xiangrongya@相容芽（compatible germs）}
\expandafter\let\csname BookOriginal@trCompatibleGerm\endcsname\trCompatibleGerm
\renewcommand{\trCompatibleGerm}{\BookTerm{T0085}{\csname BookOriginal@trCompatibleGerm\endcsname}}
\expandafter\def\csname BookIndex@T0086\endcsname{jingjubu@茎局部（stalk-local）}
\expandafter\let\csname BookOriginal@trStalkLocal\endcsname\trStalkLocal
\renewcommand{\trStalkLocal}{\BookTerm{T0086}{\csname BookOriginal@trStalkLocal\endcsname}}
\expandafter\def\csname BookIndex@T0087\endcsname{zhishuzhenghelie@指数正合列（exponential exact sequence）}
\expandafter\let\csname BookOriginal@trExponentialSequence\endcsname\trExponentialSequence
\renewcommand{\trExponentialSequence}{\BookTerm{T0087}{\csname BookOriginal@trExponentialSequence\endcsname}}
\expandafter\def\csname BookIndex@T0088\endcsname{zhishiceng@指示层（indicator sheaf）}
\expandafter\let\csname BookOriginal@trIndicatorSheaf\endcsname\trIndicatorSheaf
\renewcommand{\trIndicatorSheaf}{\BookTerm{T0088}{\csname BookOriginal@trIndicatorSheaf\endcsname}}
\expandafter\def\csname BookIndex@T0089\endcsname{ziyuceng@子预层（subpresheaf）}
\expandafter\let\csname BookOriginal@trSubpresheaf\endcsname\trSubpresheaf
\renewcommand{\trSubpresheaf}{\BookTerm{T0089}{\csname BookOriginal@trSubpresheaf\endcsname}}
\expandafter\def\csname BookIndex@T0090\endcsname{zhenghexing@正合性（exactness）}
\expandafter\let\csname BookOriginal@trExactness\endcsname\trExactness
\renewcommand{\trExactness}{\BookTerm{T0090}{\csname BookOriginal@trExactness\endcsname}}
\expandafter\def\csname BookIndex@T0091\endcsname{zuozhenghexing@左正合性（left exactness）}
\expandafter\let\csname BookOriginal@trLeftExactness\endcsname\trLeftExactness
\renewcommand{\trLeftExactness}{\BookTerm{T0091}{\csname BookOriginal@trLeftExactness\endcsname}}
\expandafter\def\csname BookIndex@T0092\endcsname{jishangdeceng@基上的层（sheaf on a base）}
\expandafter\let\csname BookOriginal@trSheafOnBase\endcsname\trSheafOnBase
\renewcommand{\trSheafOnBase}{\BookTerm{T0092}{\csname BookOriginal@trSheafOnBase\endcsname}}
\expandafter\def\csname BookIndex@T0093\endcsname{jishangdeyuceng@基上的预层（presheaf on a base）}
\expandafter\let\csname BookOriginal@trPresheafOnBase\endcsname\trPresheafOnBase
\renewcommand{\trPresheafOnBase}{\BookTerm{T0093}{\csname BookOriginal@trPresheafOnBase\endcsname}}
\expandafter\def\csname BookIndex@T0094\endcsname{jibenkaiji@基本开集（basic open set）}
\expandafter\let\csname BookOriginal@trBasicOpenSet\endcsname\trBasicOpenSet
\renewcommand{\trBasicOpenSet}{\BookTerm{T0094}{\csname BookOriginal@trBasicOpenSet\endcsname}}
\expandafter\def\csname BookIndex@T0095\endcsname{shangbiliantiaojian@上闭链条件（cocycle condition）}
\expandafter\let\csname BookOriginal@trCocycleCondition\endcsname\trCocycleCondition
\renewcommand{\trCocycleCondition}{\BookTerm{T0095}{\csname BookOriginal@trCocycleCondition\endcsname}}
\expandafter\def\csname BookIndex@T0096\endcsname{kesuo@可缩（contractible）}
\expandafter\let\csname BookOriginal@trContractible\endcsname\trContractible
\renewcommand{\trContractible}{\BookTerm{T0096}{\csname BookOriginal@trContractible\endcsname}}
\expandafter\def\csname BookIndex@T0097\endcsname{kaiqiu@开球（open ball）}
\expandafter\let\csname BookOriginal@trOpenBall\endcsname\trOpenBall
\renewcommand{\trOpenBall}{\BookTerm{T0097}{\csname BookOriginal@trOpenBall\endcsname}}
\expandafter\def\csname BookIndex@T0098\endcsname{zhicheng@支撑（support）}
\expandafter\let\csname BookOriginal@trSupport\endcsname\trSupport
\renewcommand{\trSupport}{\BookTerm{T0098}{\csname BookOriginal@trSupport\endcsname}}
\expandafter\def\csname BookIndex@T0099\endcsname{zhixiang@直像（direct image）}
\expandafter\let\csname BookOriginal@trDirectImage\endcsname\trDirectImage
\renewcommand{\trDirectImage}{\BookTerm{T0099}{\csname BookOriginal@trDirectImage\endcsname}}
\expandafter\def\csname BookIndex@T0100\endcsname{tuilayingshe@推拉映射（push-pull map）}
\expandafter\let\csname BookOriginal@trPushPullMap\endcsname\trPushPullMap
\renewcommand{\trPushPullMap}{\BookTerm{T0100}{\csname BookOriginal@trPushPullMap\endcsname}}
\expandafter\def\csname BookIndex@T0101\endcsname{lingkuozhang@零扩张（extension by zero）}
\expandafter\let\csname BookOriginal@trExtensionByZero\endcsname\trExtensionByZero
\renewcommand{\trExtensionByZero}{\BookTerm{T0101}{\csname BookOriginal@trExtensionByZero\endcsname}}
\expandafter\def\csname BookIndex@T0102\endcsname{bansuixing@伴随性（adjointness）}
\expandafter\let\csname BookOriginal@trAdjointness\endcsname\trAdjointness
\renewcommand{\trAdjointness}{\BookTerm{T0102}{\csname BookOriginal@trAdjointness\endcsname}}
\expandafter\def\csname BookIndex@T0103\endcsname{baohanyingshe@包含映射（inclusion）}
\expandafter\let\csname BookOriginal@trInclusionMap\endcsname\trInclusionMap
\renewcommand{\trInclusionMap}{\BookTerm{T0103}{\csname BookOriginal@trInclusionMap\endcsname}}
\expandafter\def\csname BookIndex@T0104\endcsname{guiji@轨迹（locus）}
\expandafter\let\csname BookOriginal@trLocus\endcsname\trLocus
\renewcommand{\trLocus}{\BookTerm{T0104}{\csname BookOriginal@trLocus\endcsname}}
\expandafter\def\csname BookIndex@T0105\endcsname{chouxiangfeihua@抽象废话（abstract nonsense）}
\expandafter\let\csname BookOriginal@trAbstractNonsense\endcsname\trAbstractNonsense
\renewcommand{\trAbstractNonsense}{\BookTerm{T0105}{\csname BookOriginal@trAbstractNonsense\endcsname}}
\expandafter\def\csname BookIndex@T0106\endcsname{tapu@拓扑（topology）}
\expandafter\let\csname BookOriginal@trTopology\endcsname\trTopology
\renewcommand{\trTopology}{\BookTerm{T0106}{\csname BookOriginal@trTopology\endcsname}}
\expandafter\def\csname BookIndex@T0107\endcsname{daishu@代数（algebraic /\allowbreak{} algebra）}
\expandafter\let\csname BookOriginal@trAlgebraic\endcsname\trAlgebraic
\renewcommand{\trAlgebraic}{\BookTerm{T0107}{\csname BookOriginal@trAlgebraic\endcsname}}
\expandafter\def\csname BookIndex@T0108\endcsname{daishuhanshu@代数函数（algebraic function）}
\expandafter\let\csname BookOriginal@trAlgebraicFunction\endcsname\trAlgebraicFunction
\renewcommand{\trAlgebraicFunction}{\BookTerm{T0108}{\csname BookOriginal@trAlgebraicFunction\endcsname}}
\expandafter\def\csname BookIndex@T0109\endcsname{jihekongjian@几何空间（geometric space）}
\expandafter\let\csname BookOriginal@trGeometricSpace\endcsname\trGeometricSpace
\renewcommand{\trGeometricSpace}{\BookTerm{T0109}{\csname BookOriginal@trGeometricSpace\endcsname}}
\expandafter\def\csname BookIndex@T0110\endcsname{dierkeshu@第二可数（second-countable）}
\expandafter\let\csname BookOriginal@trSecondCountable\endcsname\trSecondCountable
\renewcommand{\trSecondCountable}{\BookTerm{T0110}{\csname BookOriginal@trSecondCountable\endcsname}}
\expandafter\def\csname BookIndex@T0111\endcsname{fenlixing@分离性（separatedness）}
\expandafter\let\csname BookOriginal@trSeparatedness\endcsname\trSeparatedness
\renewcommand{\trSeparatedness}{\BookTerm{T0111}{\csname BookOriginal@trSeparatedness\endcsname}}
\expandafter\def\csname BookIndex@T0112\endcsname{feifenqi@非分歧（unramified）}
\expandafter\let\csname BookOriginal@trUnramified\endcsname\trUnramified
\renewcommand{\trUnramified}{\BookTerm{T0112}{\csname BookOriginal@trUnramified\endcsname}}
\expandafter\def\csname BookIndex@T0113\endcsname{yanmei@淹没（submersion）}
\expandafter\let\csname BookOriginal@trSubmersion\endcsname\trSubmersion
\renewcommand{\trSubmersion}{\BookTerm{T0113}{\csname BookOriginal@trSubmersion\endcsname}}
\expandafter\def\csname BookIndex@T0114\endcsname{bukeyue@不可约（irreducible）}
\expandafter\let\csname BookOriginal@trIrreducible\endcsname\trIrreducible
\renewcommand{\trIrreducible}{\BookTerm{T0114}{\csname BookOriginal@trIrreducible\endcsname}}
\expandafter\def\csname BookIndex@T0115\endcsname{lingdiandingli@零点定理（Nullstellensatz）}
\expandafter\let\csname BookOriginal@trNullstellensatz\endcsname\trNullstellensatz
\renewcommand{\trNullstellensatz}{\BookTerm{T0115}{\csname BookOriginal@trNullstellensatz\endcsname}}
\expandafter\def\csname BookIndex@T0116\endcsname{teyikaiji@特异开集（distinguished open set）}
\expandafter\let\csname BookOriginal@trDistinguishedOpen\endcsname\trDistinguishedOpen
\renewcommand{\trDistinguishedOpen}{\BookTerm{T0116}{\csname BookOriginal@trDistinguishedOpen\endcsname}}
\expandafter\def\csname BookIndex@T0117\endcsname{fangshezhixian@仿射直线（affine line）}
\expandafter\let\csname BookOriginal@trAffineLine\endcsname\trAffineLine
\renewcommand{\trAffineLine}{\BookTerm{T0117}{\csname BookOriginal@trAffineLine\endcsname}}
\expandafter\def\csname BookIndex@T0118\endcsname{pu@谱（spectrum）}
\expandafter\let\csname BookOriginal@trSpectrum\endcsname\trSpectrum
\renewcommand{\trSpectrum}{\BookTerm{T0118}{\csname BookOriginal@trSpectrum\endcsname}}
\expandafter\def\csname BookIndex@T0119\endcsname{oujilidezhenghuan@欧几里得整环（Euclidean domain）}
\expandafter\let\csname BookOriginal@trEuclideanDomain\endcsname\trEuclideanDomain
\renewcommand{\trEuclideanDomain}{\BookTerm{T0119}{\csname BookOriginal@trEuclideanDomain\endcsname}}
\expandafter\def\csname BookIndex@T0120\endcsname{milinggen@幂零根（nilradical）}
\expandafter\let\csname BookOriginal@trNilradical\endcsname\trNilradical
\renewcommand{\trNilradical}{\BookTerm{T0120}{\csname BookOriginal@trNilradical\endcsname}}
\expandafter\def\csname BookIndex@T0121\endcsname{chaoqumian@超曲面（hypersurface）}
\expandafter\let\csname BookOriginal@trHypersurface\endcsname\trHypersurface
\renewcommand{\trHypersurface}{\BookTerm{T0121}{\csname BookOriginal@trHypersurface\endcsname}}
\expandafter\def\csname BookIndex@T0122\endcsname{zuobiao@坐标（coordinate）}
\expandafter\let\csname BookOriginal@trCoordinate\endcsname\trCoordinate
\renewcommand{\trCoordinate}{\BookTerm{T0122}{\csname BookOriginal@trCoordinate\endcsname}}
\expandafter\def\csname BookIndex@T0123\endcsname{gonge@共轭（conjugate）}
\expandafter\let\csname BookOriginal@trConjugate\endcsname\trConjugate
\renewcommand{\trConjugate}{\BookTerm{T0123}{\csname BookOriginal@trConjugate\endcsname}}
\expandafter\def\csname BookIndex@T0124\endcsname{shengyuhuan@剩余环（residue ring）}
\expandafter\let\csname BookOriginal@trResidueRing\endcsname\trResidueRing
\renewcommand{\trResidueRing}{\BookTerm{T0124}{\csname BookOriginal@trResidueRing\endcsname}}
\expandafter\def\csname BookIndex@T0125\endcsname{yukuozhang@域扩张（field extension）}
\expandafter\let\csname BookOriginal@trFieldExtension\endcsname\trFieldExtension
\renewcommand{\trFieldExtension}{\BookTerm{T0125}{\csname BookOriginal@trFieldExtension\endcsname}}
\expandafter\def\csname BookIndex@T0126\endcsname{youxiankuozhang@有限扩张（finite field extension）}
\expandafter\let\csname BookOriginal@trFiniteFieldExtension\endcsname\trFiniteFieldExtension
\renewcommand{\trFiniteFieldExtension}{\BookTerm{T0126}{\csname BookOriginal@trFiniteFieldExtension\endcsname}}
\expandafter\def\csname BookIndex@T0127\endcsname{duioushu@对偶数（dual numbers）}
\expandafter\let\csname BookOriginal@trDualNumbers\endcsname\trDualNumbers
\renewcommand{\trDualNumbers}{\BookTerm{T0127}{\csname BookOriginal@trDualNumbers\endcsname}}
\expandafter\def\csname BookIndex@T0128\endcsname{daiyuchufa@带余除法（division algorithm）}
\expandafter\let\csname BookOriginal@trDivisionAlgorithm\endcsname\trDivisionAlgorithm
\renewcommand{\trDivisionAlgorithm}{\BookTerm{T0128}{\csname BookOriginal@trDivisionAlgorithm\endcsname}}
\expandafter\def\csname BookIndex@T0129\endcsname{youlihanshu@有理函数（rational function）}
\expandafter\let\csname BookOriginal@trRationalFunction\endcsname\trRationalFunction
\renewcommand{\trRationalFunction}{\BookTerm{T0129}{\csname BookOriginal@trRationalFunction\endcsname}}
\expandafter\def\csname BookIndex@T0130\endcsname{jidian@极点（pole）}
\expandafter\let\csname BookOriginal@trPole\endcsname\trPole
\renewcommand{\trPole}{\BookTerm{T0130}{\csname BookOriginal@trPole\endcsname}}
\expandafter\def\csname BookIndex@T0131\endcsname{xiaoshiji@消失集（vanishing set）}
\expandafter\let\csname BookOriginal@trVanishingSet\endcsname\trVanishingSet
\renewcommand{\trVanishingSet}{\BookTerm{T0131}{\csname BookOriginal@trVanishingSet\endcsname}}
\expandafter\def\csname BookIndex@T0132\endcsname{lingdianji@零点集（zero set）}
\expandafter\let\csname BookOriginal@trZeroSet\endcsname\trZeroSet
\renewcommand{\trZeroSet}{\BookTerm{T0132}{\csname BookOriginal@trZeroSet\endcsname}}
\expandafter\def\csname BookIndex@T0133\endcsname{gen@根（radical）}
\expandafter\let\csname BookOriginal@trRadical\endcsname\trRadical
\renewcommand{\trRadical}{\BookTerm{T0133}{\csname BookOriginal@trRadical\endcsname}}
\expandafter\def\csname BookIndex@T0134\endcsname{genlixiang@根理想（radical ideal）}
\expandafter\let\csname BookOriginal@trRadicalIdeal\endcsname\trRadicalIdeal
\renewcommand{\trRadicalIdeal}{\BookTerm{T0134}{\csname BookOriginal@trRadicalIdeal\endcsname}}
\expandafter\def\csname BookIndex@T0135\endcsname{bibao@闭包（closure）}
\expandafter\let\csname BookOriginal@trClosure\endcsname\trClosure
\renewcommand{\trClosure}{\BookTerm{T0135}{\csname BookOriginal@trClosure\endcsname}}
\expandafter\def\csname BookIndex@T0136\endcsname{zhalisijibibao@扎里斯基闭包（Zariski closure）}
\expandafter\let\csname BookOriginal@trZariskiClosure\endcsname\trZariskiClosure
\renewcommand{\trZariskiClosure}{\BookTerm{T0136}{\csname BookOriginal@trZariskiClosure\endcsname}}
\expandafter\def\csname BookIndex@T0137\endcsname{cu@粗（coarse）}
\expandafter\let\csname BookOriginal@trCoarse\endcsname\trCoarse
\renewcommand{\trCoarse}{\BookTerm{T0137}{\csname BookOriginal@trCoarse\endcsname}}
\expandafter\def\csname BookIndex@T0138\endcsname{yuyouxiantapu@余有限拓扑（cofinite topology）}
\expandafter\let\csname BookOriginal@trCofiniteTopology\endcsname\trCofiniteTopology
\renewcommand{\trCofiniteTopology}{\BookTerm{T0138}{\csname BookOriginal@trCofiniteTopology\endcsname}}
\expandafter\def\csname BookIndex@T0139\endcsname{zikongjiantapu@子空间拓扑（subspace topology）}
\expandafter\let\csname BookOriginal@trSubspaceTopology\endcsname\trSubspaceTopology
\renewcommand{\trSubspaceTopology}{\BookTerm{T0139}{\csname BookOriginal@trSubspaceTopology\endcsname}}
\expandafter\def\csname BookIndex@T0140\endcsname{choumi@稠密（dense）}
\expandafter\let\csname BookOriginal@trDense\endcsname\trDense
\renewcommand{\trDense}{\BookTerm{T0140}{\csname BookOriginal@trDense\endcsname}}
\expandafter\def\csname BookIndex@T0141\endcsname{xing@型（form）}
\expandafter\let\csname BookOriginal@trForm\endcsname\trForm
\renewcommand{\trForm}{\BookTerm{T0141}{\csname BookOriginal@trForm\endcsname}}
\expandafter\def\csname BookIndex@T0142\endcsname{buji@补集（complement）}
\expandafter\let\csname BookOriginal@trComplement\endcsname\trComplement
\renewcommand{\trComplement}{\BookTerm{T0142}{\csname BookOriginal@trComplement\endcsname}}
\expandafter\def\csname BookIndex@T0143\endcsname{jixiaosulixiang@极小素理想（minimal prime ideal）}
\expandafter\let\csname BookOriginal@trMinimalPrimeIdeal\endcsname\trMinimalPrimeIdeal
\renewcommand{\trMinimalPrimeIdeal}{\BookTerm{T0143}{\csname BookOriginal@trMinimalPrimeIdeal\endcsname}}
\expandafter\def\csname BookIndex@T0144\endcsname{jixiaosu@极小素（minimal prime）}
\expandafter\let\csname BookOriginal@trMinimalPrime\endcsname\trMinimalPrime
\renewcommand{\trMinimalPrime}{\BookTerm{T0144}{\csname BookOriginal@trMinimalPrime\endcsname}}
\expandafter\def\csname BookIndex@T0145\endcsname{bukeyuefenliang@不可约分量（irreducible component）}
\expandafter\let\csname BookOriginal@trIrreducibleComponent\endcsname\trIrreducibleComponent
\renewcommand{\trIrreducibleComponent}{\BookTerm{T0145}{\csname BookOriginal@trIrreducibleComponent\endcsname}}
\expandafter\def\csname BookIndex@T0146\endcsname{tehua@特化（specialization）}
\expandafter\let\csname BookOriginal@trSpecialization\endcsname\trSpecialization
\renewcommand{\trSpecialization}{\BookTerm{T0146}{\csname BookOriginal@trSpecialization\endcsname}}
\expandafter\def\csname BookIndex@T0147\endcsname{fanhua@泛化（generization）}
\expandafter\let\csname BookOriginal@trGenerization\endcsname\trGenerization
\renewcommand{\trGenerization}{\BookTerm{T0147}{\csname BookOriginal@trGenerization\endcsname}}
\expandafter\def\csname BookIndex@T0148\endcsname{fandian@泛点（general point）}
\expandafter\let\csname BookOriginal@trGeneralPoint\endcsname\trGeneralPoint
\renewcommand{\trGeneralPoint}{\BookTerm{T0148}{\csname BookOriginal@trGeneralPoint\endcsname}}
\expandafter\def\csname BookIndex@T0149\endcsname{liantongfenliang@连通分量（connected component）}
\expandafter\let\csname BookOriginal@trConnectedComponent\endcsname\trConnectedComponent
\renewcommand{\trConnectedComponent}{\BookTerm{T0149}{\csname BookOriginal@trConnectedComponent\endcsname}}
\expandafter\def\csname BookIndex@T0150\endcsname{jiangliantiaojian@降链条件（descending chain condition）}
\expandafter\let\csname BookOriginal@trDescendingChainCondition\endcsname\trDescendingChainCondition
\renewcommand{\trDescendingChainCondition}{\BookTerm{T0150}{\csname BookOriginal@trDescendingChainCondition\endcsname}}
\expandafter\def\csname BookIndex@T0151\endcsname{shengliantiaojian@升链条件（ascending chain condition）}
\expandafter\let\csname BookOriginal@trAscendingChainCondition\endcsname\trAscendingChainCondition
\renewcommand{\trAscendingChainCondition}{\BookTerm{T0151}{\csname BookOriginal@trAscendingChainCondition\endcsname}}
\expandafter\def\csname BookIndex@T0152\endcsname{midengyuan@幂等元（idempotent）}
\expandafter\let\csname BookOriginal@trIdempotent\endcsname\trIdempotent
\renewcommand{\trIdempotent}{\BookTerm{T0152}{\csname BookOriginal@trIdempotent\endcsname}}
\expandafter\def\csname BookIndex@T0153\endcsname{zhui@锥（cone）}
\expandafter\let\csname BookOriginal@trCone\endcsname\trCone
\renewcommand{\trCone}{\BookTerm{T0153}{\csname BookOriginal@trCone\endcsname}}
\expandafter\def\csname BookIndex@T0154\endcsname{chaoluzi@超滤子（ultrafilter）}
\expandafter\let\csname BookOriginal@trUltrafilter\endcsname\trUltrafilter
\renewcommand{\trUltrafilter}{\BookTerm{T0154}{\csname BookOriginal@trUltrafilter\endcsname}}
\expandafter\def\csname BookIndex@T0155\endcsname{nijinxing@拟紧性（quasicompactness）}
\expandafter\let\csname BookOriginal@trQuasicompactness\endcsname\trQuasicompactness
\renewcommand{\trQuasicompactness}{\BookTerm{T0155}{\csname BookOriginal@trQuasicompactness\endcsname}}
\expandafter\def\csname BookIndex@T0156\endcsname{teyijidi@特异基底（distinguished base）}
\expandafter\let\csname BookOriginal@trDistinguishedBase\endcsname\trDistinguishedBase
\renewcommand{\trDistinguishedBase}{\BookTerm{T0156}{\csname BookOriginal@trDistinguishedBase\endcsname}}
\expandafter\def\csname BookIndex@T0157\endcsname{qianrudian@嵌入点（embedded point）}
\expandafter\let\csname BookOriginal@trEmbeddedPoint\endcsname\trEmbeddedPoint
\renewcommand{\trEmbeddedPoint}{\BookTerm{T0157}{\csname BookOriginal@trEmbeddedPoint\endcsname}}
\expandafter\def\csname BookIndex@T0158\endcsname{danweifenjie@单位分解（partition of unity）}
\expandafter\let\csname BookOriginal@trPartitionOfUnity\endcsname\trPartitionOfUnity
\renewcommand{\trPartitionOfUnity}{\BookTerm{T0158}{\csname BookOriginal@trPartitionOfUnity\endcsname}}
\expandafter\def\csname BookIndex@T0159\endcsname{fangshetongxinyinli@仿射通信引理（Affine Communication Lemma）}
\expandafter\let\csname BookOriginal@trAffineCommunicationLemma\endcsname\trAffineCommunicationLemma
\renewcommand{\trAffineCommunicationLemma}{\BookTerm{T0159}{\csname BookOriginal@trAffineCommunicationLemma\endcsname}}
\expandafter\def\csname BookIndex@T0160\endcsname{erciqumian@二次曲面（quadric surface）}
\expandafter\let\csname BookOriginal@trQuadricSurface\endcsname\trQuadricSurface
\renewcommand{\trQuadricSurface}{\BookTerm{T0160}{\csname BookOriginal@trQuadricSurface\endcsname}}
\expandafter\def\csname BookIndex@T0161\endcsname{rongmao@绒毛（fuzz）}
\expandafter\let\csname BookOriginal@trFuzz\endcsname\trFuzz
\renewcommand{\trFuzz}{\BookTerm{T0161}{\csname BookOriginal@trFuzz\endcsname}}
\expandafter\def\csname BookIndex@T0162\endcsname{xiangbandian@相伴点（associated point）}
\expandafter\let\csname BookOriginal@trAssociatedPoint\endcsname\trAssociatedPoint
\renewcommand{\trAssociatedPoint}{\BookTerm{T0162}{\csname BookOriginal@trAssociatedPoint\endcsname}}
\expandafter\def\csname BookIndex@T0163\endcsname{bizigaixing@闭子概形（closed subscheme）}
\expandafter\let\csname BookOriginal@trClosedSubscheme\endcsname\trClosedSubscheme
\renewcommand{\trClosedSubscheme}{\BookTerm{T0163}{\csname BookOriginal@trClosedSubscheme\endcsname}}
\expandafter\def\csname BookIndex@T0164\endcsname{zhongshu@重数（multiplicity）}
\expandafter\let\csname BookOriginal@trMultiplicity\endcsname\trMultiplicity
\renewcommand{\trMultiplicity}{\BookTerm{T0164}{\csname BookOriginal@trMultiplicity\endcsname}}
\expandafter\def\csname BookIndex@T0165\endcsname{kaizigaixing@开子概形（open subscheme）}
\expandafter\let\csname BookOriginal@trOpenSubscheme\endcsname\trOpenSubscheme
\renewcommand{\trOpenSubscheme}{\BookTerm{T0165}{\csname BookOriginal@trOpenSubscheme\endcsname}}
\expandafter\def\csname BookIndex@T0166\endcsname{zhi@秩（rank）}
\expandafter\let\csname BookOriginal@trRank\endcsname\trRank
\renewcommand{\trRank}{\BookTerm{T0166}{\csname BookOriginal@trRank\endcsname}}
\expandafter\def\csname BookIndex@T0167\endcsname{tezheng@特征（characteristic）}
\expandafter\let\csname BookOriginal@trCharacteristic\endcsname\trCharacteristic
\renewcommand{\trCharacteristic}{\BookTerm{T0167}{\csname BookOriginal@trCharacteristic\endcsname}}
\expandafter\def\csname BookIndex@T0168\endcsname{shizhi@实值（real-valued）}
\expandafter\let\csname BookOriginal@trRealValued\endcsname\trRealValued
\renewcommand{\trRealValued}{\BookTerm{T0168}{\csname BookOriginal@trRealValued\endcsname}}
\expandafter\def\csname BookIndex@T0169\endcsname{zhongguoshengyudingli@中国剩余定理（Chinese Remainder Theorem）}
\expandafter\let\csname BookOriginal@trChineseRemainderTheorem\endcsname\trChineseRemainderTheorem
\renewcommand{\trChineseRemainderTheorem}{\BookTerm{T0169}{\csname BookOriginal@trChineseRemainderTheorem\endcsname}}
\expandafter\def\csname BookIndex@T0170\endcsname{zhulixiang@主理想（principal ideal）}
\expandafter\let\csname BookOriginal@trPrincipalIdeal\endcsname\trPrincipalIdeal
\renewcommand{\trPrincipalIdeal}{\BookTerm{T0170}{\csname BookOriginal@trPrincipalIdeal\endcsname}}
\expandafter\def\csname BookIndex@T0171\endcsname{shangtapu@商拓扑（quotient topology）}
\expandafter\let\csname BookOriginal@trQuotientTopology\endcsname\trQuotientTopology
\renewcommand{\trQuotientTopology}{\BookTerm{T0171}{\csname BookOriginal@trQuotientTopology\endcsname}}
\expandafter\def\csname BookIndex@T0172\endcsname{chaopingmian@超平面（hyperplane）}
\expandafter\let\csname BookOriginal@trHyperplane\endcsname\trHyperplane
\renewcommand{\trHyperplane}{\BookTerm{T0172}{\csname BookOriginal@trHyperplane\endcsname}}
\expandafter\def\csname BookIndex@T0173\endcsname{zhiwen@直纹（ruling）}
\expandafter\let\csname BookOriginal@trRuling\endcsname\trRuling
\renewcommand{\trRuling}{\BookTerm{T0173}{\csname BookOriginal@trRuling\endcsname}}
\expandafter\def\csname BookIndex@T0174\endcsname{sheyingzhixian@射影直线（projective line）}
\expandafter\let\csname BookOriginal@trProjectiveLine\endcsname\trProjectiveLine
\renewcommand{\trProjectiveLine}{\BookTerm{T0174}{\csname BookOriginal@trProjectiveLine\endcsname}}
\expandafter\def\csname BookIndex@T0175\endcsname{fencilixiang@分次理想（graded ideal）}
\expandafter\let\csname BookOriginal@trGradedIdeal\endcsname\trGradedIdeal
\renewcommand{\trGradedIdeal}{\BookTerm{T0175}{\csname BookOriginal@trGradedIdeal\endcsname}}
\expandafter\def\csname BookIndex@T0176\endcsname{fenci@分次（grading）}
\expandafter\let\csname BookOriginal@trGrading\endcsname\trGrading
\renewcommand{\trGrading}{\BookTerm{T0176}{\csname BookOriginal@trGrading\endcsname}}
\expandafter\def\csname BookIndex@T0177\endcsname{jihuan@基环（base ring）}
\expandafter\let\csname BookOriginal@trBaseRing\endcsname\trBaseRing
\renewcommand{\trBaseRing}{\BookTerm{T0177}{\csname BookOriginal@trBaseRing\endcsname}}
\expandafter\def\csname BookIndex@T0178\endcsname{wuguanlixiang@无关理想（irrelevant ideal）}
\expandafter\let\csname BookOriginal@trIrrelevantIdeal\endcsname\trIrrelevantIdeal
\renewcommand{\trIrrelevantIdeal}{\BookTerm{T0178}{\csname BookOriginal@trIrrelevantIdeal\endcsname}}
\expandafter\def\csname BookIndex@T0179\endcsname{duichengdaishu@对称代数（symmetric algebra）}
\expandafter\let\csname BookOriginal@trSymmetricAlgebra\endcsname\trSymmetricAlgebra
\renewcommand{\trSymmetricAlgebra}{\BookTerm{T0179}{\csname BookOriginal@trSymmetricAlgebra\endcsname}}
\expandafter\def\csname BookIndex@T0180\endcsname{nisheying@拟射影（quasiprojective）}
\expandafter\let\csname BookOriginal@trQuasiprojective\endcsname\trQuasiprojective
\renewcommand{\trQuasiprojective}{\BookTerm{T0180}{\csname BookOriginal@trQuasiprojective\endcsname}}
\expandafter\def\csname BookIndex@T0181\endcsname{xiaoshigaixing@消失概形（vanishing scheme）}
\expandafter\let\csname BookOriginal@trVanishingScheme\endcsname\trVanishingScheme
\renewcommand{\trVanishingScheme}{\BookTerm{T0181}{\csname BookOriginal@trVanishingScheme\endcsname}}
\expandafter\def\csname BookIndex@T0182\endcsname{sheyinghua@射影化（projectivization）}
\expandafter\let\csname BookOriginal@trProjectivization\endcsname\trProjectivization
\renewcommand{\trProjectivization}{\BookTerm{T0182}{\csname BookOriginal@trProjectivization\endcsname}}
\expandafter\def\csname BookIndex@T0183\endcsname{wuqiongyuanxian@无穷远线（line at infinity）}
\expandafter\let\csname BookOriginal@trLineAtInfinity\endcsname\trLineAtInfinity
\renewcommand{\trLineAtInfinity}{\BookTerm{T0183}{\csname BookOriginal@trLineAtInfinity\endcsname}}
\expandafter\def\csname BookIndex@T0184\endcsname{wuqiongyuanpingmian@无穷远平面（plane at infinity）}
\expandafter\let\csname BookOriginal@trPlaneAtInfinity\endcsname\trPlaneAtInfinity
\renewcommand{\trPlaneAtInfinity}{\BookTerm{T0184}{\csname BookOriginal@trPlaneAtInfinity\endcsname}}
\expandafter\def\csname BookIndex@T0185\endcsname{erciquxian@二次曲线（conic）}
\expandafter\let\csname BookOriginal@trConic\endcsname\trConic
\renewcommand{\trConic}{\BookTerm{T0185}{\csname BookOriginal@trConic\endcsname}}
\expandafter\def\csname BookIndex@T0186\endcsname{guyouxing@固有性（properness）}
\expandafter\let\csname BookOriginal@trProperness\endcsname\trProperness
\renewcommand{\trProperness}{\BookTerm{T0186}{\csname BookOriginal@trProperness\endcsname}}
\expandafter\def\csname BookIndex@T0187\endcsname{jiyuexing@既约性（reducedness）}
\expandafter\let\csname BookOriginal@trReducedness\endcsname\trReducedness
\renewcommand{\trReducedness}{\BookTerm{T0187}{\csname BookOriginal@trReducedness\endcsname}}
\expandafter\def\csname BookIndex@T0188\endcsname{zhengxing@整性（integrality）}
\expandafter\let\csname BookOriginal@trIntegrality\endcsname\trIntegrality
\renewcommand{\trIntegrality}{\BookTerm{T0188}{\csname BookOriginal@trIntegrality\endcsname}}
\expandafter\def\csname BookIndex@T0189\endcsname{fangshejubu@仿射局部（affine-local）}
\expandafter\let\csname BookOriginal@trAffineLocal\endcsname\trAffineLocal
\renewcommand{\trAffineLocal}{\BookTerm{T0189}{\csname BookOriginal@trAffineLocal\endcsname}}
\expandafter\def\csname BookIndex@T0190\endcsname{zigaixing@子概形（subscheme）}
\expandafter\let\csname BookOriginal@trSubscheme\endcsname\trSubscheme
\renewcommand{\trSubscheme}{\BookTerm{T0190}{\csname BookOriginal@trSubscheme\endcsname}}
\expandafter\def\csname BookIndex@T0191\endcsname{zhengguixing@正规性（normality）}
\expandafter\let\csname BookOriginal@trNormality\endcsname\trNormality
\renewcommand{\trNormality}{\BookTerm{T0191}{\csname BookOriginal@trNormality\endcsname}}
\expandafter\def\csname BookIndex@T0192\endcsname{cishu@次数（degree）}
\expandafter\let\csname BookOriginal@trDegree\endcsname\trDegree
\renewcommand{\trDegree}{\BookTerm{T0192}{\csname BookOriginal@trDegree\endcsname}}
\expandafter\def\csname BookIndex@T0193\endcsname{nisheyingxing@拟射影性（quasiprojectivity）}
\expandafter\let\csname BookOriginal@trQuasiprojectivity\endcsname\trQuasiprojectivity
\renewcommand{\trQuasiprojectivity}{\BookTerm{T0193}{\csname BookOriginal@trQuasiprojectivity\endcsname}}
\expandafter\def\csname BookIndex@T0194\endcsname{zhengbi@整闭（integrally closed）}
\expandafter\let\csname BookOriginal@trIntegrallyClosed\endcsname\trIntegrallyClosed
\renewcommand{\trIntegrallyClosed}{\BookTerm{T0194}{\csname BookOriginal@trIntegrallyClosed\endcsname}}
\expandafter\def\csname BookIndex@T0195\endcsname{zhengbixing@整闭性（integral closedness）}
\expandafter\let\csname BookOriginal@trIntegralClosedness\endcsname\trIntegralClosedness
\renewcommand{\trIntegralClosedness}{\BookTerm{T0195}{\csname BookOriginal@trIntegralClosedness\endcsname}}
\expandafter\def\csname BookIndex@T0196\endcsname{zhengbibao@整闭包（integral closure）}
\expandafter\let\csname BookOriginal@trIntegralClosure\endcsname\trIntegralClosure
\renewcommand{\trIntegralClosure}{\BookTerm{T0196}{\csname BookOriginal@trIntegralClosure\endcsname}}
\expandafter\def\csname BookIndex@T0197\endcsname{weiyifenjie@唯一分解（factorial）}
\expandafter\let\csname BookOriginal@trFactorial\endcsname\trFactorial
\renewcommand{\trFactorial}{\BookTerm{T0197}{\csname BookOriginal@trFactorial\endcsname}}
\expandafter\def\csname BookIndex@T0198\endcsname{weiyifenjiexing@唯一分解性（factoriality）}
\expandafter\let\csname BookOriginal@trFactoriality\endcsname\trFactoriality
\renewcommand{\trFactoriality}{\BookTerm{T0198}{\csname BookOriginal@trFactoriality\endcsname}}
\expandafter\def\csname BookIndex@T0199\endcsname{fenmulixiang@分母理想（ideal of denominators）}
\expandafter\let\csname BookOriginal@trIdealOfDenominators\endcsname\trIdealOfDenominators
\renewcommand{\trIdealOfDenominators}{\BookTerm{T0199}{\csname BookOriginal@trIdealOfDenominators\endcsname}}
\expandafter\def\csname BookIndex@T0200\endcsname{ercixing@二次型（quadratic form）}
\expandafter\let\csname BookOriginal@trQuadraticForm\endcsname\trQuadraticForm
\renewcommand{\trQuadraticForm}{\BookTerm{T0200}{\csname BookOriginal@trQuadraticForm\endcsname}}
\expandafter\def\csname BookIndex@T0201\endcsname{tuoyuanquxian@椭圆曲线（elliptic curve）}
\expandafter\let\csname BookOriginal@trEllipticCurve\endcsname\trEllipticCurve
\renewcommand{\trEllipticCurve}{\BookTerm{T0201}{\csname BookOriginal@trEllipticCurve\endcsname}}
\expandafter\def\csname BookIndex@T0202\endcsname{leiqun@类群（class group）}
\expandafter\let\csname BookOriginal@trClassGroup\endcsname\trClassGroup
\renewcommand{\trClassGroup}{\BookTerm{T0202}{\csname BookOriginal@trClassGroup\endcsname}}
\expandafter\def\csname BookIndex@T0203\endcsname{wupingfangyinzi@无平方因子（square-free）}
\expandafter\let\csname BookOriginal@trSquareFree\endcsname\trSquareFree
\renewcommand{\trSquareFree}{\BookTerm{T0203}{\csname BookOriginal@trSquareFree\endcsname}}
\expandafter\def\csname BookIndex@T0204\endcsname{kefen@可分（separable）}
\expandafter\let\csname BookOriginal@trSeparable\endcsname\trSeparable
\renewcommand{\trSeparable}{\BookTerm{T0204}{\csname BookOriginal@trSeparable\endcsname}}
\expandafter\def\csname BookIndex@T0205\endcsname{gaodu@高度（height）}
\expandafter\let\csname BookOriginal@trHeight\endcsname\trHeight
\renewcommand{\trHeight}{\BookTerm{T0205}{\csname BookOriginal@trHeight\endcsname}}
\expandafter\def\csname BookIndex@T0206\endcsname{manzhi@满秩（full rank）}
\expandafter\let\csname BookOriginal@trFullRank\endcsname\trFullRank
\renewcommand{\trFullRank}{\BookTerm{T0206}{\csname BookOriginal@trFullRank\endcsname}}
\expandafter\def\csname BookIndex@T0207\endcsname{jidazhi@极大秩（maximal rank）}
\expandafter\let\csname BookOriginal@trMaximalRank\endcsname\trMaximalRank
\renewcommand{\trMaximalRank}{\BookTerm{T0207}{\csname BookOriginal@trMaximalRank\endcsname}}
\expandafter\def\csname BookIndex@T0208\endcsname{zhengzexing@正则性（regularity）}
\expandafter\let\csname BookOriginal@trRegularity\endcsname\trRegularity
\renewcommand{\trRegularity}{\BookTerm{T0208}{\csname BookOriginal@trRegularity\endcsname}}
\expandafter\def\csname BookIndex@T0209\endcsname{duoxiangshihuan@多项式环（polynomial ring）}
\expandafter\let\csname BookOriginal@trPolynomialRing\endcsname\trPolynomialRing
\renewcommand{\trPolynomialRing}{\BookTerm{T0209}{\csname BookOriginal@trPolynomialRing\endcsname}}
\expandafter\def\csname BookIndex@T0210\endcsname{duijiaohua@对角化（diagonalize）}
\expandafter\let\csname BookOriginal@trDiagonalize\endcsname\trDiagonalize
\renewcommand{\trDiagonalize}{\BookTerm{T0210}{\csname BookOriginal@trDiagonalize\endcsname}}
\expandafter\def\csname BookIndex@T0211\endcsname{zhongshipingtan@忠实平坦（faithfully flat）}
\expandafter\let\csname BookOriginal@trFaithfullyFlat\endcsname\trFaithfullyFlat
\renewcommand{\trFaithfullyFlat}{\BookTerm{T0211}{\csname BookOriginal@trFaithfullyFlat\endcsname}}
\expandafter\def\csname BookIndex@T0212\endcsname{xiajiang@下降（descent）}
\expandafter\let\csname BookOriginal@trDescent\endcsname\trDescent
\renewcommand{\trDescent}{\BookTerm{T0212}{\csname BookOriginal@trDescent\endcsname}}
\expandafter\def\csname BookIndex@T0213\endcsname{nao@挠（torsion）}
\expandafter\let\csname BookOriginal@trTorsion\endcsname\trTorsion
\renewcommand{\trTorsion}{\BookTerm{T0213}{\csname BookOriginal@trTorsion\endcsname}}
\expandafter\def\csname BookIndex@T0214\endcsname{naozimo@挠子模（torsion submodule）}
\expandafter\let\csname BookOriginal@trTorsionSubmodule\endcsname\trTorsionSubmodule
\renewcommand{\trTorsionSubmodule}{\BookTerm{T0214}{\csname BookOriginal@trTorsionSubmodule\endcsname}}
\expandafter\def\csname BookIndex@T0215\endcsname{teyitapu@特异拓扑（distinguished topology）}
\expandafter\let\csname BookOriginal@trDistinguishedTopology\endcsname\trDistinguishedTopology
\renewcommand{\trDistinguishedTopology}{\BookTerm{T0215}{\csname BookOriginal@trDistinguishedTopology\endcsname}}
\expandafter\def\csname BookIndex@T0216\endcsname{ziyouceng@自由层（free sheaf）}
\expandafter\let\csname BookOriginal@trFreeSheaf\endcsname\trFreeSheaf
\renewcommand{\trFreeSheaf}{\BookTerm{T0216}{\csname BookOriginal@trFreeSheaf\endcsname}}
\expandafter\def\csname BookIndex@T0217\endcsname{fujiexijihe@复解析几何（complex analytic geometry）}
\expandafter\let\csname BookOriginal@trComplexAnalyticGeometry\endcsname\trComplexAnalyticGeometry
\renewcommand{\trComplexAnalyticGeometry}{\BookTerm{T0217}{\csname BookOriginal@trComplexAnalyticGeometry\endcsname}}
\expandafter\def\csname BookIndex@T0218\endcsname{teyifangshejidi@特异仿射基底（distinguished affine base）}
\expandafter\let\csname BookOriginal@trDistinguishedAffineBase\endcsname\trDistinguishedAffineBase
\renewcommand{\trDistinguishedAffineBase}{\BookTerm{T0218}{\csname BookOriginal@trDistinguishedAffineBase\endcsname}}
\expandafter\def\csname BookIndex@T0219\endcsname{pingzhan@平展（etale）}
\expandafter\let\csname BookOriginal@trEtale\endcsname\trEtale
\renewcommand{\trEtale}{\BookTerm{T0219}{\csname BookOriginal@trEtale\endcsname}}
\expandafter\def\csname BookIndex@T0220\endcsname{jing@景（site）}
\expandafter\let\csname BookOriginal@trSite\endcsname\trSite
\renewcommand{\trSite}{\BookTerm{T0220}{\csname BookOriginal@trSite\endcsname}}
\expandafter\def\csname BookIndex@T0221\endcsname{tapusi@拓扑斯（topos）}
\expandafter\let\csname BookOriginal@trTopos\endcsname\trTopos
\renewcommand{\trTopos}{\BookTerm{T0221}{\csname BookOriginal@trTopos\endcsname}}
\expandafter\def\csname BookIndex@T0222\endcsname{feiajimide@非阿基米德（non-Archimedean）}
\expandafter\let\csname BookOriginal@trNonArchimedean\endcsname\trNonArchimedean
\renewcommand{\trNonArchimedean}{\BookTerm{T0222}{\csname BookOriginal@trNonArchimedean\endcsname}}
\expandafter\def\csname BookIndex@T0223\endcsname{quanbuliantong@全不连通（totally disconnected）}
\expandafter\let\csname BookOriginal@trTotallyDisconnected\endcsname\trTotallyDisconnected
\renewcommand{\trTotallyDisconnected}{\BookTerm{T0223}{\csname BookOriginal@trTotallyDisconnected\endcsname}}
\expandafter\def\csname BookIndex@T0224\endcsname{jiyu@基域（base field）}
\expandafter\let\csname BookOriginal@trBaseField\endcsname\trBaseField
\renewcommand{\trBaseField}{\BookTerm{T0224}{\csname BookOriginal@trBaseField\endcsname}}
\expandafter\def\csname BookIndex@T0225\endcsname{xunhuanmo@循环模（cyclic module）}
\expandafter\let\csname BookOriginal@trCyclicModule\endcsname\trCyclicModule
\renewcommand{\trCyclicModule}{\BookTerm{T0225}{\csname BookOriginal@trCyclicModule\endcsname}}
\expandafter\def\csname BookIndex@T0226\endcsname{hechenglie@合成列（composition series）}
\expandafter\let\csname BookOriginal@trCompositionSeries\endcsname\trCompositionSeries
\renewcommand{\trCompositionSeries}{\BookTerm{T0226}{\csname BookOriginal@trCompositionSeries\endcsname}}
\expandafter\def\csname BookIndex@T0227\endcsname{dan@单（simple）}
\expandafter\let\csname BookOriginal@trSimple\endcsname\trSimple
\renewcommand{\trSimple}{\BookTerm{T0227}{\csname BookOriginal@trSimple\endcsname}}
\expandafter\def\csname BookIndex@T0228\endcsname{zhangdu@长度（length）}
\expandafter\let\csname BookOriginal@trLength\endcsname\trLength
\renewcommand{\trLength}{\BookTerm{T0228}{\csname BookOriginal@trLength\endcsname}}
\expandafter\def\csname BookIndex@T0229\endcsname{youxianzhangdu@有限长度（finite length）}
\expandafter\let\csname BookOriginal@trFiniteLength\endcsname\trFiniteLength
\renewcommand{\trFiniteLength}{\BookTerm{T0229}{\csname BookOriginal@trFiniteLength\endcsname}}
\expandafter\def\csname BookIndex@T0230\endcsname{luguo@滤过（filtration /\allowbreak{} filtered）}
\expandafter\let\csname BookOriginal@trFiltration\endcsname\trFiltration
\renewcommand{\trFiltration}{\BookTerm{T0230}{\csname BookOriginal@trFiltration\endcsname}}
\expandafter\def\csname BookIndex@T0231\endcsname{zhengguiziqun@正规子群（normal subgroup）}
\expandafter\let\csname BookOriginal@trNormalSubgroup\endcsname\trNormalSubgroup
\renewcommand{\trNormalSubgroup}{\BookTerm{T0231}{\csname BookOriginal@trNormalSubgroup\endcsname}}
\expandafter\def\csname BookIndex@T0232\endcsname{diertonggoudingli@第二同构定理（second isomorphism theorem）}
\expandafter\let\csname BookOriginal@trSecondIsomorphismTheorem\endcsname\trSecondIsomorphismTheorem
\renewcommand{\trSecondIsomorphismTheorem}{\BookTerm{T0232}{\csname BookOriginal@trSecondIsomorphismTheorem\endcsname}}
\expandafter\def\csname BookIndex@T0233\endcsname{dijiangguina@递降归纳（descending induction）}
\expandafter\let\csname BookOriginal@trDescendingInduction\endcsname\trDescendingInduction
\renewcommand{\trDescendingInduction}{\BookTerm{T0233}{\csname BookOriginal@trDescendingInduction\endcsname}}
\expandafter\def\csname BookIndex@T0234\endcsname{tonggoulei@同构类（isomorphism class）}
\expandafter\let\csname BookOriginal@trIsomorphismClass\endcsname\trIsomorphismClass
\renewcommand{\trIsomorphismClass}{\BookTerm{T0234}{\csname BookOriginal@trIsomorphismClass\endcsname}}
\expandafter\def\csname BookIndex@T0235\endcsname{jiaxi@加细（refine）}
\expandafter\let\csname BookOriginal@trRefine\endcsname\trRefine
\renewcommand{\trRefine}{\BookTerm{T0235}{\csname BookOriginal@trRefine\endcsname}}
\expandafter\def\csname BookIndex@T0236\endcsname{linghua@零化（annihilate）}
\expandafter\let\csname BookOriginal@trAnnihilate\endcsname\trAnnihilate
\renewcommand{\trAnnihilate}{\BookTerm{T0236}{\csname BookOriginal@trAnnihilate\endcsname}}
\expandafter\def\csname BookIndex@T0237\endcsname{fuzhihuan@赋值环（valuation ring）}
\expandafter\let\csname BookOriginal@trValuationRing\endcsname\trValuationRing
\renewcommand{\trValuationRing}{\BookTerm{T0237}{\csname BookOriginal@trValuationRing\endcsname}}
\expandafter\def\csname BookIndex@T0238\endcsname{shouxianmijishu@受限幂级数（restricted power series）}
\expandafter\let\csname BookOriginal@trRestrictedPowerSeries\endcsname\trRestrictedPowerSeries
\renewcommand{\trRestrictedPowerSeries}{\BookTerm{T0238}{\csname BookOriginal@trRestrictedPowerSeries\endcsname}}
\expandafter\def\csname BookIndex@T0239\endcsname{daishubibao@代数闭包（algebraic closure）}
\expandafter\let\csname BookOriginal@trAlgebraicClosure\endcsname\trAlgebraicClosure
\renewcommand{\trAlgebraicClosure}{\BookTerm{T0239}{\csname BookOriginal@trAlgebraicClosure\endcsname}}
\expandafter\def\csname BookIndex@T0240\endcsname{xiangbansulixiang@相伴素理想（associated prime ideal）}
\expandafter\let\csname BookOriginal@trAssociatedPrimeIdeal\endcsname\trAssociatedPrimeIdeal
\renewcommand{\trAssociatedPrimeIdeal}{\BookTerm{T0240}{\csname BookOriginal@trAssociatedPrimeIdeal\endcsname}}
\expandafter\def\csname BookIndex@T0241\endcsname{qianrusu@嵌入素（embedded prime）}
\expandafter\let\csname BookOriginal@trEmbeddedPrime\endcsname\trEmbeddedPrime
\renewcommand{\trEmbeddedPrime}{\BookTerm{T0241}{\csname BookOriginal@trEmbeddedPrime\endcsname}}
\expandafter\def\csname BookIndex@T0242\endcsname{linghualixiang@零化理想（annihilator ideal）}
\expandafter\let\csname BookOriginal@trAnnihilatorIdeal\endcsname\trAnnihilatorIdeal
\renewcommand{\trAnnihilatorIdeal}{\BookTerm{T0242}{\csname BookOriginal@trAnnihilatorIdeal\endcsname}}
\expandafter\def\csname BookIndex@T0243\endcsname{zhunsu@准素（primary）}
\expandafter\let\csname BookOriginal@trPrimary\endcsname\trPrimary
\renewcommand{\trPrimary}{\BookTerm{T0243}{\csname BookOriginal@trPrimary\endcsname}}
\expandafter\def\csname BookIndex@T0244\endcsname{zhunsufenjie@准素分解（primary decomposition）}
\expandafter\let\csname BookOriginal@trPrimaryDecomposition\endcsname\trPrimaryDecomposition
\renewcommand{\trPrimaryDecomposition}{\BookTerm{T0244}{\csname BookOriginal@trPrimaryDecomposition\endcsname}}
\expandafter\def\csname BookIndex@T0245\endcsname{quanfenshihuan@全分式环（total fraction ring）}
\expandafter\let\csname BookOriginal@trTotalFractionRing\endcsname\trTotalFractionRing
\renewcommand{\trTotalFractionRing}{\BookTerm{T0245}{\csname BookOriginal@trTotalFractionRing\endcsname}}
\expandafter\def\csname BookIndex@T0246\endcsname{quanshanghuan@全商环（total quotient ring）}
\expandafter\let\csname BookOriginal@trTotalQuotientRing\endcsname\trTotalQuotientRing
\renewcommand{\trTotalQuotientRing}{\BookTerm{T0246}{\csname BookOriginal@trTotalQuotientRing\endcsname}}
\expandafter\def\csname BookIndex@T0247\endcsname{dingyiyu@定义域（domain of definition）}
\expandafter\let\csname BookOriginal@trDomainOfDefinition\endcsname\trDomainOfDefinition
\renewcommand{\trDomainOfDefinition}{\BookTerm{T0247}{\csname BookOriginal@trDomainOfDefinition\endcsname}}
\expandafter\def\csname BookIndex@T0248\endcsname{budingxingguiji@不定性轨迹（indeterminacy locus）}
\expandafter\let\csname BookOriginal@trIndeterminacyLocus\endcsname\trIndeterminacyLocus
\renewcommand{\trIndeterminacyLocus}{\BookTerm{T0248}{\csname BookOriginal@trIndeterminacyLocus\endcsname}}
\expandafter\def\csname BookIndex@T0249\endcsname{gaixinglunchoumi@概形论稠密（scheme-theoretically dense）}
\expandafter\let\csname BookOriginal@trSchemeTheoreticallyDense\endcsname\trSchemeTheoreticallyDense
\renewcommand{\trSchemeTheoreticallyDense}{\BookTerm{T0249}{\csname BookOriginal@trSchemeTheoreticallyDense\endcsname}}
\expandafter\def\csname BookIndex@T0250\endcsname{wanquanxiangjiao@完全相交（complete intersection）}
\expandafter\let\csname BookOriginal@trCompleteIntersection\endcsname\trCompleteIntersection
\renewcommand{\trCompleteIntersection}{\BookTerm{T0250}{\csname BookOriginal@trCompleteIntersection\endcsname}}
\expandafter\def\csname BookIndex@T0251\endcsname{zhengzexulie@正则序列（regular sequence）}
\expandafter\let\csname BookOriginal@trRegularSequence\endcsname\trRegularSequence
\renewcommand{\trRegularSequence}{\BookTerm{T0251}{\csname BookOriginal@trRegularSequence\endcsname}}
\expandafter\def\csname BookIndex@T0252\endcsname{jubufuhuankongjian@局部赋环空间（locally ringed space）}
\expandafter\let\csname BookOriginal@trLocallyRingedSpace\endcsname\trLocallyRingedSpace
\renewcommand{\trLocallyRingedSpace}{\BookTerm{T0252}{\csname BookOriginal@trLocallyRingedSpace\endcsname}}
\expandafter\def\csname BookIndex@T0253\endcsname{jiegoutaishe@结构态射（structure morphism）}
\expandafter\let\csname BookOriginal@trStructureMorphism\endcsname\trStructureMorphism
\renewcommand{\trStructureMorphism}{\BookTerm{T0253}{\csname BookOriginal@trStructureMorphism\endcsname}}
\expandafter\def\csname BookIndex@T0254\endcsname{zhidian@值点（valued point）}
\expandafter\let\csname BookOriginal@trValuedPoint\endcsname\trValuedPoint
\renewcommand{\trValuedPoint}{\BookTerm{T0254}{\csname BookOriginal@trValuedPoint\endcsname}}
\expandafter\def\csname BookIndex@T0255\endcsname{jiexihua@解析化（analytification）}
\expandafter\let\csname BookOriginal@trAnalytification\endcsname\trAnalytification
\renewcommand{\trAnalytification}{\BookTerm{T0255}{\csname BookOriginal@trAnalytification\endcsname}}
\expandafter\def\csname BookIndex@T0256\endcsname{yucu@预簇（prevariety）}
\expandafter\let\csname BookOriginal@trPrevariety\endcsname\trPrevariety
\renewcommand{\trPrevariety}{\BookTerm{T0256}{\csname BookOriginal@trPrevariety\endcsname}}
\expandafter\def\csname BookIndex@T0257\endcsname{youlidian@有理点（rational point）}
\expandafter\let\csname BookOriginal@trRationalPoint\endcsname\trRationalPoint
\renewcommand{\trRationalPoint}{\BookTerm{T0257}{\csname BookOriginal@trRationalPoint\endcsname}}
\expandafter\def\csname BookIndex@T0258\endcsname{pingfancong@平凡丛（trivial bundle）}
\expandafter\let\csname BookOriginal@trTrivialBundle\endcsname\trTrivialBundle
\renewcommand{\trTrivialBundle}{\BookTerm{T0258}{\csname BookOriginal@trTrivialBundle\endcsname}}
\expandafter\def\csname BookIndex@T0259\endcsname{qiexiangliang@切向量（tangent vector）}
\expandafter\let\csname BookOriginal@trTangentVector\endcsname\trTangentVector
\renewcommand{\trTangentVector}{\BookTerm{T0259}{\csname BookOriginal@trTangentVector\endcsname}}
\expandafter\def\csname BookIndex@T0260\endcsname{fencihuanyingshe@分次环映射（map of graded rings）}
\expandafter\let\csname BookOriginal@trGradedRingMap\endcsname\trGradedRingMap
\renewcommand{\trGradedRingMap}{\BookTerm{T0260}{\csname BookOriginal@trGradedRingMap\endcsname}}
\expandafter\def\csname BookIndex@T0261\endcsname{weiluoneisaizihuan@韦罗内塞子环（veronese subring）}
\expandafter\let\csname BookOriginal@trVeroneseSubring\endcsname\trVeroneseSubring
\renewcommand{\trVeroneseSubring}{\BookTerm{T0261}{\csname BookOriginal@trVeroneseSubring\endcsname}}
\expandafter\def\csname BookIndex@T0262\endcsname{weiluoneisaiqianru@韦罗内塞嵌入（veronese embedding）}
\expandafter\let\csname BookOriginal@trVeroneseEmbedding\endcsname\trVeroneseEmbedding
\renewcommand{\trVeroneseEmbedding}{\BookTerm{T0262}{\csname BookOriginal@trVeroneseEmbedding\endcsname}}
\expandafter\def\csname BookIndex@T0263\endcsname{danxiangshi@单项式（monomial）}
\expandafter\let\csname BookOriginal@trMonomial\endcsname\trMonomial
\renewcommand{\trMonomial}{\BookTerm{T0263}{\csname BookOriginal@trMonomial\endcsname}}
\expandafter\def\csname BookIndex@T0264\endcsname{zhongxinfenci@重新分次（regrading）}
\expandafter\let\csname BookOriginal@trRegrading\endcsname\trRegrading
\renewcommand{\trRegrading}{\BookTerm{T0264}{\csname BookOriginal@trRegrading\endcsname}}
\expandafter\def\csname BookIndex@T0265\endcsname{youliyingshe@有理映射（rational map）}
\expandafter\let\csname BookOriginal@trRationalMap\endcsname\trRationalMap
\renewcommand{\trRationalMap}{\BookTerm{T0265}{\csname BookOriginal@trRationalMap\endcsname}}
\expandafter\def\csname BookIndex@T0266\endcsname{zhipei@支配（dominant）}
\expandafter\let\csname BookOriginal@trDominant\endcsname\trDominant
\renewcommand{\trDominant}{\BookTerm{T0266}{\csname BookOriginal@trDominant\endcsname}}
\expandafter\def\csname BookIndex@T0267\endcsname{shuangyouli@双有理（birational）}
\expandafter\let\csname BookOriginal@trBirational\endcsname\trBirational
\renewcommand{\trBirational}{\BookTerm{T0267}{\csname BookOriginal@trBirational\endcsname}}
\expandafter\def\csname BookIndex@T0268\endcsname{youli@有理（rational）}
\expandafter\let\csname BookOriginal@trRational\endcsname\trRational
\renewcommand{\trRational}{\BookTerm{T0268}{\csname BookOriginal@trRational\endcsname}}
\expandafter\def\csname BookIndex@T0269\endcsname{qunduixiang@群对象（group object）}
\expandafter\let\csname BookOriginal@trGroupObject\endcsname\trGroupObject
\renewcommand{\trGroupObject}{\BookTerm{T0269}{\csname BookOriginal@trGroupObject\endcsname}}
\expandafter\def\csname BookIndex@T0270\endcsname{chengfaqun@乘法群（multiplicative group）}
\expandafter\let\csname BookOriginal@trMultiplicativeGroup\endcsname\trMultiplicativeGroup
\renewcommand{\trMultiplicativeGroup}{\BookTerm{T0270}{\csname BookOriginal@trMultiplicativeGroup\endcsname}}
\expandafter\def\csname BookIndex@T0271\endcsname{huangaixing@环概形（ring scheme）}
\expandafter\let\csname BookOriginal@trRingScheme\endcsname\trRingScheme
\renewcommand{\trRingScheme}{\BookTerm{T0271}{\csname BookOriginal@trRingScheme\endcsname}}
\expandafter\def\csname BookIndex@T0272\endcsname{yibanxianxingqun@一般线性群（general linear group）}
\expandafter\let\csname BookOriginal@trGeneralLinearGroup\endcsname\trGeneralLinearGroup
\renewcommand{\trGeneralLinearGroup}{\BookTerm{T0272}{\csname BookOriginal@trGeneralLinearGroup\endcsname}}
\expandafter\def\csname BookIndex@T0273\endcsname{xinglieshi@行列式（determinant）}
\expandafter\let\csname BookOriginal@trDeterminant\endcsname\trDeterminant
\renewcommand{\trDeterminant}{\BookTerm{T0273}{\csname BookOriginal@trDeterminant\endcsname}}
\expandafter\def\csname BookIndex@T0274\endcsname{huopufudaishu@霍普夫代数（hopf algebra）}
\expandafter\let\csname BookOriginal@trHopfAlgebra\endcsname\trHopfAlgebra
\renewcommand{\trHopfAlgebra}{\BookTerm{T0274}{\csname BookOriginal@trHopfAlgebra\endcsname}}
\expandafter\def\csname BookIndex@T0275\endcsname{yuchengfa@余乘法（comultiplication）}
\expandafter\let\csname BookOriginal@trComultiplication\endcsname\trComultiplication
\renewcommand{\trComultiplication}{\BookTerm{T0275}{\csname BookOriginal@trComultiplication\endcsname}}
\expandafter\def\csname BookIndex@T0276\endcsname{qicu@旗簇（flag variety）}
\expandafter\let\csname BookOriginal@trFlagVariety\endcsname\trFlagVariety
\renewcommand{\trFlagVariety}{\BookTerm{T0276}{\csname BookOriginal@trFlagVariety\endcsname}}
\expandafter\def\csname BookIndex@T0277\endcsname{bufenqicu@部分旗簇（partial flag variety）}
\expandafter\let\csname BookOriginal@trPartialFlagVariety\endcsname\trPartialFlagVariety
\renewcommand{\trPartialFlagVariety}{\BookTerm{T0277}{\csname BookOriginal@trPartialFlagVariety\endcsname}}
\expandafter\def\csname BookIndex@T0278\endcsname{zikongjian@子空间（subspace）}
\expandafter\let\csname BookOriginal@trSubspace\endcsname\trSubspace
\renewcommand{\trSubspace}{\BookTerm{T0278}{\csname BookOriginal@trSubspace\endcsname}}
\expandafter\def\csname BookIndex@T0279\endcsname{fenceng@分层（stratification）}
\expandafter\let\csname BookOriginal@trStratification\endcsname\trStratification
\renewcommand{\trStratification}{\BookTerm{T0279}{\csname BookOriginal@trStratification\endcsname}}
\expandafter\def\csname BookIndex@T0280\endcsname{fanshi@范式（normal form）}
\expandafter\let\csname BookOriginal@trNormalForm\endcsname\trNormalForm
\renewcommand{\trNormalForm}{\BookTerm{T0280}{\csname BookOriginal@trNormalForm\endcsname}}
\expandafter\def\csname BookIndex@T0281\endcsname{xingyuejianjietixing@行约简阶梯形（row-reduced echelon form）}
\expandafter\let\csname BookOriginal@trRowReducedEchelonForm\endcsname\trRowReducedEchelonForm
\renewcommand{\trRowReducedEchelonForm}{\BookTerm{T0281}{\csname BookOriginal@trRowReducedEchelonForm\endcsname}}
\expandafter\def\csname BookIndex@T0282\endcsname{shubotebaoqiang@舒伯特胞腔（Schubert cell）}
\expandafter\let\csname BookOriginal@trSchubertCell\endcsname\trSchubertCell
\renewcommand{\trSchubertCell}{\BookTerm{T0282}{\csname BookOriginal@trSchubertCell\endcsname}}
\expandafter\def\csname BookIndex@T0283\endcsname{pulukeqianru@普吕克嵌入（plucker embedding）}
\expandafter\let\csname BookOriginal@trPluckerEmbedding\endcsname\trPluckerEmbedding
\renewcommand{\trPluckerEmbedding}{\BookTerm{T0283}{\csname BookOriginal@trPluckerEmbedding\endcsname}}
\expandafter\def\csname BookIndex@T0284\endcsname{zhengkuozhang@整扩张（integral extension）}
\expandafter\let\csname BookOriginal@trIntegralExtension\endcsname\trIntegralExtension
\renewcommand{\trIntegralExtension}{\BookTerm{T0284}{\csname BookOriginal@trIntegralExtension\endcsname}}
\expandafter\def\csname BookIndex@T0285\endcsname{xiaoqudingli@消去定理（Cancellation Theorem）}
\expandafter\let\csname BookOriginal@trCancellationTheorem\endcsname\trCancellationTheorem
\renewcommand{\trCancellationTheorem}{\BookTerm{T0285}{\csname BookOriginal@trCancellationTheorem\endcsname}}
\expandafter\def\csname BookIndex@T0286\endcsname{yagebusengen@雅各布森根（jacobson radical）}
\expandafter\let\csname BookOriginal@trJacobsonRadical\endcsname\trJacobsonRadical
\renewcommand{\trJacobsonRadical}{\BookTerm{T0286}{\csname BookOriginal@trJacobsonRadical\endcsname}}
\expandafter\def\csname BookIndex@T0287\endcsname{shouyiduoxiangshi@首一多项式（monic polynomial）}
\expandafter\let\csname BookOriginal@trMonicPolynomial\endcsname\trMonicPolynomial
\renewcommand{\trMonicPolynomial}{\BookTerm{T0287}{\csname BookOriginal@trMonicPolynomial\endcsname}}
\expandafter\def\csname BookIndex@T0288\endcsname{bansuijuzhen@伴随矩阵（adjugate matrix）}
\expandafter\let\csname BookOriginal@trAdjugateMatrix\endcsname\trAdjugateMatrix
\renewcommand{\trAdjugateMatrix}{\BookTerm{T0288}{\csname BookOriginal@trAdjugateMatrix\endcsname}}
\expandafter\def\csname BookIndex@T0289\endcsname{zidaishu@子代数（subalgebra）}
\expandafter\let\csname BookOriginal@trSubalgebra\endcsname\trSubalgebra
\renewcommand{\trSubalgebra}{\BookTerm{T0289}{\csname BookOriginal@trSubalgebra\endcsname}}
\expandafter\def\csname BookIndex@T0290\endcsname{bi@闭（closed）}
\expandafter\let\csname BookOriginal@trClosed\endcsname\trClosed
\renewcommand{\trClosed}{\BookTerm{T0290}{\csname BookOriginal@trClosed\endcsname}}
\expandafter\def\csname BookIndex@T0291\endcsname{kegouzao@可构造（constructible）}
\expandafter\let\csname BookOriginal@trConstructible\endcsname\trConstructible
\renewcommand{\trConstructible}{\BookTerm{T0291}{\csname BookOriginal@trConstructible\endcsname}}
\expandafter\def\csname BookIndex@T0292\endcsname{xiaoyuanlilunjibendingli@消元理论基本定理（Fundamental Theorem of Elimination Theory）}
\expandafter\let\csname BookOriginal@trFundamentalTheoremOfEliminationTheory\endcsname\trFundamentalTheoremOfEliminationTheory
\renewcommand{\trFundamentalTheoremOfEliminationTheory}{\BookTerm{T0292}{\csname BookOriginal@trFundamentalTheoremOfEliminationTheory\endcsname}}
\expandafter\def\csname BookIndex@T0293\endcsname{jieshi@结式（resultant）}
\expandafter\let\csname BookOriginal@trResultant\endcsname\trResultant
\renewcommand{\trResultant}{\BookTerm{T0293}{\csname BookOriginal@trResultant\endcsname}}
\expandafter\def\csname BookIndex@T0294\endcsname{liangcixiaoqu@量词消去（elimination of quantifiers）}
\expandafter\let\csname BookOriginal@trEliminationOfQuantifiers\endcsname\trEliminationOfQuantifiers
\renewcommand{\trEliminationOfQuantifiers}{\BookTerm{T0294}{\csname BookOriginal@trEliminationOfQuantifiers\endcsname}}
\expandafter\def\csname BookIndex@T0295\endcsname{niyouxian@拟有限（quasifinite）}
\expandafter\let\csname BookOriginal@trQuasifinite\endcsname\trQuasifinite
\renewcommand{\trQuasifinite}{\BookTerm{T0295}{\csname BookOriginal@trQuasifinite\endcsname}}
\expandafter\def\csname BookIndex@T0296\endcsname{zhencishu@真次数（true degree）}
\expandafter\let\csname BookOriginal@trTrueDegree\endcsname\trTrueDegree
\renewcommand{\trTrueDegree}{\BookTerm{T0296}{\csname BookOriginal@trTrueDegree\endcsname}}
\expandafter\def\csname BookIndex@T0297\endcsname{xianweiguiji@纤维轨迹（fibral locus）}
\expandafter\let\csname BookOriginal@trFibralLocus\endcsname\trFibralLocus
\renewcommand{\trFibralLocus}{\BookTerm{T0297}{\csname BookOriginal@trFibralLocus\endcsname}}
\expandafter\def\csname BookIndex@T0298\endcsname{shibiyu@实闭域（real closed field）}
\expandafter\let\csname BookOriginal@trRealClosedField\endcsname\trRealClosedField
\renewcommand{\trRealClosedField}{\BookTerm{T0298}{\csname BookOriginal@trRealClosedField\endcsname}}
\expandafter\def\csname BookIndex@T0299\endcsname{youxian@有限（finite）}
\expandafter\let\csname BookOriginal@trFinite\endcsname\trFinite
\renewcommand{\trFinite}{\BookTerm{T0299}{\csname BookOriginal@trFinite\endcsname}}
\expandafter\def\csname BookIndex@T0300\endcsname{zitongtai@自同态（endomorphism）}
\expandafter\let\csname BookOriginal@trEndomorphism\endcsname\trEndomorphism
\renewcommand{\trEndomorphism}{\BookTerm{T0300}{\csname BookOriginal@trEndomorphism\endcsname}}
\expandafter\def\csname BookIndex@T0301\endcsname{zhalisijizhudingli@扎里斯基主定理（Zariski main theorem）}
\expandafter\let\csname BookOriginal@trZariskiMainTheorem\endcsname\trZariskiMainTheorem
\renewcommand{\trZariskiMainTheorem}{\BookTerm{T0301}{\csname BookOriginal@trZariskiMainTheorem\endcsname}}
\expandafter\def\csname BookIndex@T0302\endcsname{fenqifugai@分歧覆盖（branched cover）}
\expandafter\let\csname BookOriginal@trBranchedCover\endcsname\trBranchedCover
\renewcommand{\trBranchedCover}{\BookTerm{T0302}{\csname BookOriginal@trBranchedCover\endcsname}}
\expandafter\def\csname BookIndex@T0303\endcsname{jiedian@结点（node）}
\expandafter\let\csname BookOriginal@trNode\endcsname\trNode
\renewcommand{\trNode}{\BookTerm{T0303}{\csname BookOriginal@trNode\endcsname}}
\expandafter\def\csname BookIndex@T0304\endcsname{zhongyanyingshe@重言映射（tautological map）}
\expandafter\let\csname BookOriginal@trTautologicalMap\endcsname\trTautologicalMap
\renewcommand{\trTautologicalMap}{\BookTerm{T0304}{\csname BookOriginal@trTautologicalMap\endcsname}}
\expandafter\def\csname BookIndex@T0305\endcsname{youxianduiyi@有限对一（finite-to-one）}
\expandafter\let\csname BookOriginal@trFiniteToOne\endcsname\trFiniteToOne
\renewcommand{\trFiniteToOne}{\BookTerm{T0305}{\csname BookOriginal@trFiniteToOne\endcsname}}
\expandafter\def\csname BookIndex@T0306\endcsname{changshejiashe@常设假设（standing hypothesis）}
\expandafter\let\csname BookOriginal@trStandingHypothesis\endcsname\trStandingHypothesis
\renewcommand{\trStandingHypothesis}{\BookTerm{T0306}{\csname BookOriginal@trStandingHypothesis\endcsname}}
\expandafter\def\csname BookIndex@T0307\endcsname{toushexitong@投射系统（projective system）}
\expandafter\let\csname BookOriginal@trProjectiveSystem\endcsname\trProjectiveSystem
\renewcommand{\trProjectiveSystem}{\BookTerm{T0307}{\csname BookOriginal@trProjectiveSystem\endcsname}}
\expandafter\def\csname BookIndex@T0308\endcsname{kai@开（open）}
\expandafter\let\csname BookOriginal@trOpen\endcsname\trOpen
\renewcommand{\trOpen}{\BookTerm{T0308}{\csname BookOriginal@trOpen\endcsname}}
\expandafter\def\csname BookIndex@T0309\endcsname{lixiangceng@理想层（ideal sheaf /\allowbreak{} sheaf of ideals）}
\expandafter\let\csname BookOriginal@trIdealSheaf\endcsname\trIdealSheaf
\renewcommand{\trIdealSheaf}{\BookTerm{T0309}{\csname BookOriginal@trIdealSheaf\endcsname}}
\expandafter\def\csname BookIndex@T0310\endcsname{gaixinglun@概形论（scheme-theoretic）}
\expandafter\let\csname BookOriginal@trSchemeTheoretic\endcsname\trSchemeTheoretic
\renewcommand{\trSchemeTheoretic}{\BookTerm{T0310}{\csname BookOriginal@trSchemeTheoretic\endcsname}}
\expandafter\def\csname BookIndex@T0311\endcsname{chongshi@充实（enrichment）}
\expandafter\let\csname BookOriginal@trEnrichment\endcsname\trEnrichment
\renewcommand{\trEnrichment}{\BookTerm{T0311}{\csname BookOriginal@trEnrichment\endcsname}}
\expandafter\def\csname BookIndex@T0312\endcsname{niuqusanciquxian@扭曲三次曲线（twisted cubic）}
\expandafter\let\csname BookOriginal@trTwistedCubic\endcsname\trTwistedCubic
\renewcommand{\trTwistedCubic}{\BookTerm{T0312}{\csname BookOriginal@trTwistedCubic\endcsname}}
\expandafter\def\csname BookIndex@T0313\endcsname{youlizhengguiquxian@有理正规曲线（rational normal curve）}
\expandafter\let\csname BookOriginal@trRationalNormalCurve\endcsname\trRationalNormalCurve
\renewcommand{\trRationalNormalCurve}{\BookTerm{T0313}{\csname BookOriginal@trRationalNormalCurve\endcsname}}
\expandafter\def\csname BookIndex@T0314\endcsname{xianxingkongjian@线性空间（linear space）}
\expandafter\let\csname BookOriginal@trLinearSpace\endcsname\trLinearSpace
\renewcommand{\trLinearSpace}{\BookTerm{T0314}{\csname BookOriginal@trLinearSpace\endcsname}}
\expandafter\def\csname BookIndex@T0315\endcsname{jiaquansheyingkongjian@加权射影空间（weighted projective space）}
\expandafter\let\csname BookOriginal@trWeightedProjectiveSpace\endcsname\trWeightedProjectiveSpace
\renewcommand{\trWeightedProjectiveSpace}{\BookTerm{T0315}{\csname BookOriginal@trWeightedProjectiveSpace\endcsname}}
\expandafter\def\csname BookIndex@T0316\endcsname{zishi@子式（minor）}
\expandafter\let\csname BookOriginal@trMinor\endcsname\trMinor
\renewcommand{\trMinor}{\BookTerm{T0316}{\csname BookOriginal@trMinor\endcsname}}
\expandafter\def\csname BookIndex@T0317\endcsname{sheyingwanbeihua@射影完备化（projective completion）}
\expandafter\let\csname BookOriginal@trProjectiveCompletion\endcsname\trProjectiveCompletion
\renewcommand{\trProjectiveCompletion}{\BookTerm{T0317}{\csname BookOriginal@trProjectiveCompletion\endcsname}}
\expandafter\def\csname BookIndex@T0318\endcsname{youxiaokadiyechuzi@有效卡蒂耶除子（effective cartier divisor）}
\expandafter\let\csname BookOriginal@trEffectiveCartierDivisor\endcsname\trEffectiveCartierDivisor
\renewcommand{\trEffectiveCartierDivisor}{\BookTerm{T0318}{\csname BookOriginal@trEffectiveCartierDivisor\endcsname}}
\expandafter\def\csname BookIndex@T0319\endcsname{zhengzeqianru@正则嵌入（regular embedding）}
\expandafter\let\csname BookOriginal@trRegularEmbedding\endcsname\trRegularEmbedding
\renewcommand{\trRegularEmbedding}{\BookTerm{T0319}{\csname BookOriginal@trRegularEmbedding\endcsname}}
\expandafter\def\csname BookIndex@T0320\endcsname{jubuzhu@局部主（locally principal）}
\expandafter\let\csname BookOriginal@trLocallyPrincipal\endcsname\trLocallyPrincipal
\renewcommand{\trLocallyPrincipal}{\BookTerm{T0320}{\csname BookOriginal@trLocallyPrincipal\endcsname}}
\expandafter\def\csname BookIndex@T0321\endcsname{qiepian@切片（slicing）}
\expandafter\let\csname BookOriginal@trSlicing\endcsname\trSlicing
\renewcommand{\trSlicing}{\BookTerm{T0321}{\csname BookOriginal@trSlicing\endcsname}}
\expandafter\def\csname BookIndex@T0322\endcsname{jiyuehua@既约化（reduction）}
\expandafter\let\csname BookOriginal@trReduction\endcsname\trReduction
\renewcommand{\trReduction}{\BookTerm{T0322}{\csname BookOriginal@trReduction\endcsname}}
\expandafter\def\csname BookIndex@T0323\endcsname{gaixinglunxiang@概形论像（scheme-theoretic image）}
\expandafter\let\csname BookOriginal@trSchemeTheoreticImage\endcsname\trSchemeTheoreticImage
\renewcommand{\trSchemeTheoreticImage}{\BookTerm{T0323}{\csname BookOriginal@trSchemeTheoreticImage\endcsname}}
\expandafter\def\csname BookIndex@T0324\endcsname{zhalisijiceng@扎里斯基层（Zariski sheaf）}
\expandafter\let\csname BookOriginal@trZariskiSheaf\endcsname\trZariskiSheaf
\renewcommand{\trZariskiSheaf}{\BookTerm{T0324}{\csname BookOriginal@trZariskiSheaf\endcsname}}
\expandafter\def\csname BookIndex@T0325\endcsname{zihanzi@子函子（subfunctor）}
\expandafter\let\csname BookOriginal@trSubfunctor\endcsname\trSubfunctor
\renewcommand{\trSubfunctor}{\BookTerm{T0325}{\csname BookOriginal@trSubfunctor\endcsname}}
\expandafter\def\csname BookIndex@T0326\endcsname{zhuqixingkongjian@主齐性空间（torsor）}
\expandafter\let\csname BookOriginal@trTorsor\endcsname\trTorsor
\renewcommand{\trTorsor}{\BookTerm{T0326}{\csname BookOriginal@trTorsor\endcsname}}
\expandafter\def\csname BookIndex@T0327\endcsname{tousheyouxiantapu@投射有限拓扑（profinite topology）}
\expandafter\let\csname BookOriginal@trProfiniteTopology\endcsname\trProfiniteTopology
\renewcommand{\trProfiniteTopology}{\BookTerm{T0327}{\csname BookOriginal@trProfiniteTopology\endcsname}}
\expandafter\def\csname BookIndex@T0328\endcsname{kuozhang@扩张（extension）}
\expandafter\let\csname BookOriginal@trExtension\endcsname\trExtension
\renewcommand{\trExtension}{\BookTerm{T0328}{\csname BookOriginal@trExtension\endcsname}}
\expandafter\def\csname BookIndex@T0329\endcsname{jihedian@几何点（geometric point）}
\expandafter\let\csname BookOriginal@trGeometricPoint\endcsname\trGeometricPoint
\renewcommand{\trGeometricPoint}{\BookTerm{T0329}{\csname BookOriginal@trGeometricPoint\endcsname}}
\expandafter\def\csname BookIndex@T0330\endcsname{jihexianwei@几何纤维（geometric fiber）}
\expandafter\let\csname BookOriginal@trGeometricFiber\endcsname\trGeometricFiber
\renewcommand{\trGeometricFiber}{\BookTerm{T0330}{\csname BookOriginal@trGeometricFiber\endcsname}}
\expandafter\def\csname BookIndex@T0331\endcsname{jiheliantong@几何连通（geometrically connected）}
\expandafter\let\csname BookOriginal@trGeometricallyConnected\endcsname\trGeometricallyConnected
\renewcommand{\trGeometricallyConnected}{\BookTerm{T0331}{\csname BookOriginal@trGeometricallyConnected\endcsname}}
\expandafter\def\csname BookIndex@T0332\endcsname{jihebukeyue@几何不可约（geometrically irreducible）}
\expandafter\let\csname BookOriginal@trGeometricallyIrreducible\endcsname\trGeometricallyIrreducible
\renewcommand{\trGeometricallyIrreducible}{\BookTerm{T0332}{\csname BookOriginal@trGeometricallyIrreducible\endcsname}}
\expandafter\def\csname BookIndex@T0333\endcsname{jihezheng@几何整（geometrically integral）}
\expandafter\let\csname BookOriginal@trGeometricallyIntegral\endcsname\trGeometricallyIntegral
\renewcommand{\trGeometricallyIntegral}{\BookTerm{T0333}{\csname BookOriginal@trGeometricallyIntegral\endcsname}}
\expandafter\def\csname BookIndex@T0334\endcsname{jihejiyue@几何既约（geometrically reduced）}
\expandafter\let\csname BookOriginal@trGeometricallyReduced\endcsname\trGeometricallyReduced
\renewcommand{\trGeometricallyReduced}{\BookTerm{T0334}{\csname BookOriginal@trGeometricallyReduced\endcsname}}
\expandafter\def\csname BookIndex@T0335\endcsname{fandanshe@泛单射（universally injective）}
\expandafter\let\csname BookOriginal@trUniversallyInjective\endcsname\trUniversallyInjective
\renewcommand{\trUniversallyInjective}{\BookTerm{T0335}{\csname BookOriginal@trUniversallyInjective\endcsname}}
\expandafter\def\csname BookIndex@T0336\endcsname{fanbi@泛闭（universally closed）}
\expandafter\let\csname BookOriginal@trUniversallyClosed\endcsname\trUniversallyClosed
\renewcommand{\trUniversallyClosed}{\BookTerm{T0336}{\csname BookOriginal@trUniversallyClosed\endcsname}}
\expandafter\def\csname BookIndex@T0337\endcsname{fankai@泛开（universally open）}
\expandafter\let\csname BookOriginal@trUniversallyOpen\endcsname\trUniversallyOpen
\renewcommand{\trUniversallyOpen}{\BookTerm{T0337}{\csname BookOriginal@trUniversallyOpen\endcsname}}
\expandafter\def\csname BookIndex@T0338\endcsname{chunbukefen@纯不可分（purely inseparable）}
\expandafter\let\csname BookOriginal@trPurelyInseparable\endcsname\trPurelyInseparable
\renewcommand{\trPurelyInseparable}{\BookTerm{T0338}{\csname BookOriginal@trPurelyInseparable\endcsname}}
\expandafter\def\csname BookIndex@T0339\endcsname{kefenbi@可分闭（separably closed）}
\expandafter\let\csname BookOriginal@trSeparablyClosed\endcsname\trSeparablyClosed
\renewcommand{\trSeparablyClosed}{\BookTerm{T0339}{\csname BookOriginal@trSeparablyClosed\endcsname}}
\expandafter\def\csname BookIndex@T0340\endcsname{wanquanbibao@完全闭包（perfect closure）}
\expandafter\let\csname BookOriginal@trPerfectClosure\endcsname\trPerfectClosure
\renewcommand{\trPerfectClosure}{\BookTerm{T0340}{\csname BookOriginal@trPerfectClosure\endcsname}}
\expandafter\def\csname BookIndex@T0341\endcsname{wanquanyu@完全域（perfect field）}
\expandafter\let\csname BookOriginal@trPerfectField\endcsname\trPerfectField
\renewcommand{\trPerfectField}{\BookTerm{T0341}{\csname BookOriginal@trPerfectField\endcsname}}
\expandafter\def\csname BookIndex@T0342\endcsname{kefenshengcheng@可分生成（separably generated）}
\expandafter\let\csname BookOriginal@trSeparablyGenerated\endcsname\trSeparablyGenerated
\renewcommand{\trSeparablyGenerated}{\BookTerm{T0342}{\csname BookOriginal@trSeparablyGenerated\endcsname}}
\expandafter\def\csname BookIndex@T0343\endcsname{chaoyueji@超越基（transcendence base）}
\expandafter\let\csname BookOriginal@trTranscendenceBase\endcsname\trTranscendenceBase
\renewcommand{\trTranscendenceBase}{\BookTerm{T0343}{\csname BookOriginal@trTranscendenceBase\endcsname}}
\expandafter\def\csname BookIndex@T0344\endcsname{fenlichaoyueji@分离超越基（separating transcendence base）}
\expandafter\let\csname BookOriginal@trSeparatingTranscendenceBase\endcsname\trSeparatingTranscendenceBase
\renewcommand{\trSeparatingTranscendenceBase}{\BookTerm{T0344}{\csname BookOriginal@trSeparatingTranscendenceBase\endcsname}}
\expandafter\def\csname BookIndex@T0345\endcsname{bukefencishu@不可分次数（inseparable degree）}
\expandafter\let\csname BookOriginal@trInseparableDegree\endcsname\trInseparableDegree
\renewcommand{\trInseparableDegree}{\BookTerm{T0345}{\csname BookOriginal@trInseparableDegree\endcsname}}
\expandafter\def\csname BookIndex@T0346\endcsname{jixiaoduoxiangshi@极小多项式（minimal polynomial）}
\expandafter\let\csname BookOriginal@trMinimalPolynomial\endcsname\trMinimalPolynomial
\renewcommand{\trMinimalPolynomial}{\BookTerm{T0346}{\csname BookOriginal@trMinimalPolynomial\endcsname}}
\expandafter\def\csname BookIndex@T0347\endcsname{jububiqianru@局部闭嵌入（locally closed embedding）}
\expandafter\let\csname BookOriginal@trLocallyClosedEmbedding\endcsname\trLocallyClosedEmbedding
\renewcommand{\trLocallyClosedEmbedding}{\BookTerm{T0347}{\csname BookOriginal@trLocallyClosedEmbedding\endcsname}}
\expandafter\def\csname BookIndex@T0348\endcsname{heli@合理（reasonable）}
\expandafter\let\csname BookOriginal@trReasonable\endcsname\trReasonable
\renewcommand{\trReasonable}{\BookTerm{T0348}{\csname BookOriginal@trReasonable\endcsname}}
\expandafter\def\csname BookIndex@T0349\endcsname{zu@族（family）}
\expandafter\let\csname BookOriginal@trFamily\endcsname\trFamily
\renewcommand{\trFamily}{\BookTerm{T0349}{\csname BookOriginal@trFamily\endcsname}}
\expandafter\def\csname BookIndex@T0350\endcsname{jigaixing@基概形（base scheme）}
\expandafter\let\csname BookOriginal@trBaseScheme\endcsname\trBaseScheme
\renewcommand{\trBaseScheme}{\BookTerm{T0350}{\csname BookOriginal@trBaseScheme\endcsname}}
\expandafter\def\csname BookIndex@T0351\endcsname{gaixinglunyuanxiang@概形论原像（scheme-theoretic preimage）}
\expandafter\let\csname BookOriginal@trSchemeTheoreticPreimage\endcsname\trSchemeTheoreticPreimage
\renewcommand{\trSchemeTheoreticPreimage}{\BookTerm{T0351}{\csname BookOriginal@trSchemeTheoreticPreimage\endcsname}}
\expandafter\def\csname BookIndex@T0352\endcsname{gaixinglunxianwei@概形论纤维（scheme-theoretic fiber）}
\expandafter\let\csname BookOriginal@trSchemeTheoreticFiber\endcsname\trSchemeTheoreticFiber
\renewcommand{\trSchemeTheoreticFiber}{\BookTerm{T0352}{\csname BookOriginal@trSchemeTheoreticFiber\endcsname}}
\expandafter\def\csname BookIndex@T0353\endcsname{yibanxianwei@一般纤维（generic fiber）}
\expandafter\let\csname BookOriginal@trGenericFiber\endcsname\trGenericFiber
\renewcommand{\trGenericFiber}{\BookTerm{T0353}{\csname BookOriginal@trGenericFiber\endcsname}}
\expandafter\def\csname BookIndex@T0354\endcsname{yibanyouxian@一般有限（generically finite）}
\expandafter\let\csname BookOriginal@trGenericallyFinite\endcsname\trGenericallyFinite
\renewcommand{\trGenericallyFinite}{\BookTerm{T0354}{\csname BookOriginal@trGenericallyFinite\endcsname}}
\expandafter\def\csname BookIndex@T0355\endcsname{liwaichuzi@例外除子（exceptional divisor）}
\expandafter\let\csname BookOriginal@trExceptionalDivisor\endcsname\trExceptionalDivisor
\renewcommand{\trExceptionalDivisor}{\BookTerm{T0355}{\csname BookOriginal@trExceptionalDivisor\endcsname}}
\expandafter\def\csname BookIndex@T0356\endcsname{jibianhuantu@基变换图（base change diagram）}
\expandafter\let\csname BookOriginal@trBaseChangeDiagram\endcsname\trBaseChangeDiagram
\renewcommand{\trBaseChangeDiagram}{\BookTerm{T0356}{\csname BookOriginal@trBaseChangeDiagram\endcsname}}
\expandafter\def\csname BookIndex@T0357\endcsname{huanji@换基（change of base）}
\expandafter\let\csname BookOriginal@trChangeOfBase\endcsname\trChangeOfBase
\renewcommand{\trChangeOfBase}{\BookTerm{T0357}{\csname BookOriginal@trChangeOfBase\endcsname}}
\expandafter\def\csname BookIndex@T0358\endcsname{paowuxian@抛物线（parabola）}
\expandafter\let\csname BookOriginal@trParabola\endcsname\trParabola
\renewcommand{\trParabola}{\BookTerm{T0358}{\csname BookOriginal@trParabola\endcsname}}
\expandafter\def\csname BookIndex@T0359\endcsname{biaozhunkaiji@标准开集（standard open set）}
\expandafter\let\csname BookOriginal@trStandardOpenSet\endcsname\trStandardOpenSet
\renewcommand{\trStandardOpenSet}{\BookTerm{T0359}{\csname BookOriginal@trStandardOpenSet\endcsname}}
\expandafter\def\csname BookIndex@T0360\endcsname{shuyu@数域（number field）}
\expandafter\let\csname BookOriginal@trNumberField\endcsname\trNumberField
\renewcommand{\trNumberField}{\BookTerm{T0360}{\csname BookOriginal@trNumberField\endcsname}}
\expandafter\def\csname BookIndex@T0361\endcsname{zhengshuhuan@整数环（ring of integers）}
\expandafter\let\csname BookOriginal@trRingOfIntegers\endcsname\trRingOfIntegers
\renewcommand{\trRingOfIntegers}{\BookTerm{T0361}{\csname BookOriginal@trRingOfIntegers\endcsname}}
\expandafter\def\csname BookIndex@T0362\endcsname{jiandian@尖点（cusp）}
\expandafter\let\csname BookOriginal@trCusp\endcsname\trCusp
\renewcommand{\trCusp}{\BookTerm{T0362}{\csname BookOriginal@trCusp\endcsname}}
\expandafter\def\csname BookIndex@T0363\endcsname{shuangqiedian@双切点（tacnode）}
\expandafter\let\csname BookOriginal@trTacnode\endcsname\trTacnode
\renewcommand{\trTacnode}{\BookTerm{T0363}{\csname BookOriginal@trTacnode\endcsname}}
\expandafter\def\csname BookIndex@T0364\endcsname{wuzuobiao@无坐标（coordinate-free）}
\expandafter\let\csname BookOriginal@trCoordinateFree\endcsname\trCoordinateFree
\renewcommand{\trCoordinateFree}{\BookTerm{T0364}{\csname BookOriginal@trCoordinateFree\endcsname}}
\expandafter\def\csname BookIndex@T0365\endcsname{niehepingmian@捏合平面（pinched plane）}
\expandafter\let\csname BookOriginal@trPinchedPlane\endcsname\trPinchedPlane
\renewcommand{\trPinchedPlane}{\BookTerm{T0365}{\csname BookOriginal@trPinchedPlane\endcsname}}
\expandafter\def\csname BookIndex@T0366\endcsname{tuxiangtaishe@图像态射（graph morphism）}
\expandafter\let\csname BookOriginal@trGraphMorphism\endcsname\trGraphMorphism
\renewcommand{\trGraphMorphism}{\BookTerm{T0366}{\csname BookOriginal@trGraphMorphism\endcsname}}
\expandafter\def\csname BookIndex@T0367\endcsname{zicu@子簇（subvariety）}
\expandafter\let\csname BookOriginal@trSubvariety\endcsname\trSubvariety
\renewcommand{\trSubvariety}{\BookTerm{T0367}{\csname BookOriginal@trSubvariety\endcsname}}
\expandafter\def\csname BookIndex@T0368\endcsname{xiaoquxingzhi@消去性质（cancellation property）}
\expandafter\let\csname BookOriginal@trCancellationProperty\endcsname\trCancellationProperty
\renewcommand{\trCancellationProperty}{\BookTerm{T0368}{\csname BookOriginal@trCancellationProperty\endcsname}}
\expandafter\def\csname BookIndex@T0369\endcsname{fangsheduijiao@仿射对角（affine-diagonal）}
\expandafter\let\csname BookOriginal@trAffineDiagonal\endcsname\trAffineDiagonal
\renewcommand{\trAffineDiagonal}{\BookTerm{T0369}{\csname BookOriginal@trAffineDiagonal\endcsname}}
\expandafter\def\csname BookIndex@T0370\endcsname{bukeyuexing@不可约性（irreducibility）}
\expandafter\let\csname BookOriginal@trIrreducibility\endcsname\trIrreducibility
\renewcommand{\trIrreducibility}{\BookTerm{T0370}{\csname BookOriginal@trIrreducibility\endcsname}}
\expandafter\def\csname BookIndex@T0371\endcsname{jibendian@基本点（fundamental points）}
\expandafter\let\csname BookOriginal@trFundamentalPoint\endcsname\trFundamentalPoint
\renewcommand{\trFundamentalPoint}{\BookTerm{T0371}{\csname BookOriginal@trFundamentalPoint\endcsname}}
\expandafter\def\csname BookIndex@T0372\endcsname{quncu@群簇（group variety）}
\expandafter\let\csname BookOriginal@trGroupVariety\endcsname\trGroupVariety
\renewcommand{\trGroupVariety}{\BookTerm{T0372}{\csname BookOriginal@trGroupVariety\endcsname}}
\expandafter\def\csname BookIndex@T0373\endcsname{daishuqun@代数群（algebraic group）}
\expandafter\let\csname BookOriginal@trAlgebraicGroup\endcsname\trAlgebraicGroup
\renewcommand{\trAlgebraicGroup}{\BookTerm{T0373}{\csname BookOriginal@trAlgebraicGroup\endcsname}}
\expandafter\def\csname BookIndex@T0374\endcsname{abeiercu@阿贝尔簇（abelian variety）}
\expandafter\let\csname BookOriginal@trAbelianVariety\endcsname\trAbelianVariety
\renewcommand{\trAbelianVariety}{\BookTerm{T0374}{\csname BookOriginal@trAbelianVariety\endcsname}}
\expandafter\def\csname BookIndex@T0375\endcsname{fuhuanmian@复环面（complex torus）}
\expandafter\let\csname BookOriginal@trComplexTorus\endcsname\trComplexTorus
\renewcommand{\trComplexTorus}{\BookTerm{T0375}{\csname BookOriginal@trComplexTorus\endcsname}}
\expandafter\def\csname BookIndex@T0376\endcsname{ge@格（lattice）}
\expandafter\let\csname BookOriginal@trLattice\endcsname\trLattice
\renewcommand{\trLattice}{\BookTerm{T0376}{\csname BookOriginal@trLattice\endcsname}}
\expandafter\def\csname BookIndex@T0377\endcsname{danliantong@单连通（simply connected）}
\expandafter\let\csname BookOriginal@trSimplyConnected\endcsname\trSimplyConnected
\renewcommand{\trSimplyConnected}{\BookTerm{T0377}{\csname BookOriginal@trSimplyConnected\endcsname}}
\expandafter\def\csname BookIndex@T0378\endcsname{fanfugai@泛覆盖（universal cover）}
\expandafter\let\csname BookOriginal@trUniversalCover\endcsname\trUniversalCover
\renewcommand{\trUniversalCover}{\BookTerm{T0378}{\csname BookOriginal@trUniversalCover\endcsname}}
\expandafter\def\csname BookIndex@T0379\endcsname{pingyi@平移（translation）}
\expandafter\let\csname BookOriginal@trTranslation\endcsname\trTranslation
\renewcommand{\trTranslation}{\BookTerm{T0379}{\csname BookOriginal@trTranslation\endcsname}}
\expandafter\def\csname BookIndex@T0380\endcsname{danweidian@单位点（identity point）}
\expandafter\let\csname BookOriginal@trIdentityPoint\endcsname\trIdentityPoint
\renewcommand{\trIdentityPoint}{\BookTerm{T0380}{\csname BookOriginal@trIdentityPoint\endcsname}}
\expandafter\def\csname BookIndex@T0381\endcsname{wanbei@完备（complete）}
\expandafter\let\csname BookOriginal@trComplete\endcsname\trComplete
\renewcommand{\trComplete}{\BookTerm{T0381}{\csname BookOriginal@trComplete\endcsname}}
\expandafter\def\csname BookIndex@T0382\endcsname{biyingshe@闭映射（closed map）}
\expandafter\let\csname BookOriginal@trClosedMap\endcsname\trClosedMap
\renewcommand{\trClosedMap}{\BookTerm{T0382}{\csname BookOriginal@trClosedMap\endcsname}}
\expandafter\def\csname BookIndex@T0383\endcsname{ningju@凝聚（coherent）}
\expandafter\let\csname BookOriginal@trCoh\endcsname\trCoh
\renewcommand{\trCoh}{\BookTerm{T0383}{\csname BookOriginal@trCoh\endcsname}}
\expandafter\def\csname BookIndex@T0384\endcsname{zhalisijitapu@扎里斯基拓扑（Zariski topology）}
\expandafter\let\csname BookOriginal@trZariskiTopology\endcsname\trZariskiTopology
\renewcommand{\trZariskiTopology}{\BookTerm{T0384}{\csname BookOriginal@trZariskiTopology\endcsname}}
\expandafter\def\csname BookIndex@T0385\endcsname{yibandian@一般点（generic point）}
\expandafter\let\csname BookOriginal@trGenericPoint\endcsname\trGenericPoint
\renewcommand{\trGenericPoint}{\BookTerm{T0385}{\csname BookOriginal@trGenericPoint\endcsname}}
\expandafter\def\csname BookIndex@T0386\endcsname{bidian@闭点（closed point）}
\expandafter\let\csname BookOriginal@trClosedPoint\endcsname\trClosedPoint
\renewcommand{\trClosedPoint}{\BookTerm{T0386}{\csname BookOriginal@trClosedPoint\endcsname}}
\expandafter\def\csname BookIndex@T0387\endcsname{miling@幂零（nilpotent）}
\expandafter\let\csname BookOriginal@trNilpotent\endcsname\trNilpotent
\renewcommand{\trNilpotent}{\BookTerm{T0387}{\csname BookOriginal@trNilpotent\endcsname}}
\expandafter\def\csname BookIndex@T0388\endcsname{jiyue@既约（reduced）}
\expandafter\let\csname BookOriginal@trReduced\endcsname\trReduced
\renewcommand{\trReduced}{\BookTerm{T0388}{\csname BookOriginal@trReduced\endcsname}}
\expandafter\def\csname BookIndex@T0389\endcsname{zheng@整（integral）}
\expandafter\let\csname BookOriginal@trIntegral\endcsname\trIntegral
\renewcommand{\trIntegral}{\BookTerm{T0389}{\csname BookOriginal@trIntegral\endcsname}}
\expandafter\def\csname BookIndex@T0390\endcsname{zhenggui@正规（normal）}
\expandafter\let\csname BookOriginal@trNormal\endcsname\trNormal
\renewcommand{\trNormal}{\BookTerm{T0390}{\csname BookOriginal@trNormal\endcsname}}
\expandafter\def\csname BookIndex@T0391\endcsname{zhengze@正则（regular）}
\expandafter\let\csname BookOriginal@trRegular\endcsname\trRegular
\renewcommand{\trRegular}{\BookTerm{T0391}{\csname BookOriginal@trRegular\endcsname}}
\expandafter\def\csname BookIndex@T0392\endcsname{guanghua@光滑（smooth）}
\expandafter\let\csname BookOriginal@trSmooth\endcsname\trSmooth
\renewcommand{\trSmooth}{\BookTerm{T0392}{\csname BookOriginal@trSmooth\endcsname}}
\expandafter\def\csname BookIndex@T0393\endcsname{weishu@维数（dimension）}
\expandafter\let\csname BookOriginal@trDimension\endcsname\trDimension
\renewcommand{\trDimension}{\BookTerm{T0393}{\csname BookOriginal@trDimension\endcsname}}
\expandafter\def\csname BookIndex@T0394\endcsname{yuweishu@余维数（codimension）}
\expandafter\let\csname BookOriginal@trCodimension\endcsname\trCodimension
\renewcommand{\trCodimension}{\BookTerm{T0394}{\csname BookOriginal@trCodimension\endcsname}}
\expandafter\def\csname BookIndex@T0395\endcsname{chunwei@纯维（pure dimensional）}
\expandafter\let\csname BookOriginal@trPureDimensional\endcsname\trPureDimensional
\renewcommand{\trPureDimensional}{\BookTerm{T0395}{\csname BookOriginal@trPureDimensional\endcsname}}
\expandafter\def\csname BookIndex@T0396\endcsname{chunweishu@纯维数（pure dimension）}
\expandafter\let\csname BookOriginal@trPureDimension\endcsname\trPureDimension
\renewcommand{\trPureDimension}{\BookTerm{T0396}{\csname BookOriginal@trPureDimension\endcsname}}
\expandafter\def\csname BookIndex@T0397\endcsname{dengwei@等维（equidimensional）}
\expandafter\let\csname BookOriginal@trEquidimensional\endcsname\trEquidimensional
\renewcommand{\trEquidimensional}{\BookTerm{T0397}{\csname BookOriginal@trEquidimensional\endcsname}}
\expandafter\def\csname BookIndex@T0398\endcsname{nweiti@$n$ 维体（n-fold）}
\expandafter\let\csname BookOriginal@trNFold\endcsname\trNFold
\renewcommand{\trNFold}{\BookTerm{T0398}{\csname BookOriginal@trNFold\endcsname}}
\expandafter\def\csname BookIndex@T0399\endcsname{nuotezhengguihua@诺特正规化（Noether normalization）}
\expandafter\let\csname BookOriginal@trNoetherNormalization\endcsname\trNoetherNormalization
\renewcommand{\trNoetherNormalization}{\BookTerm{T0399}{\csname BookOriginal@trNoetherNormalization\endcsname}}
\expandafter\def\csname BookIndex@T0400\endcsname{daishuwuguan@代数无关（algebraically independent）}
\expandafter\let\csname BookOriginal@trAlgebraicallyIndependent\endcsname\trAlgebraicallyIndependent
\renewcommand{\trAlgebraicallyIndependent}{\BookTerm{T0400}{\csname BookOriginal@trAlgebraicallyIndependent\endcsname}}
\expandafter\def\csname BookIndex@T0401\endcsname{daishuxiangguan@代数相关（algebraically dependent）}
\expandafter\let\csname BookOriginal@trAlgebraicallyDependent\endcsname\trAlgebraicallyDependent
\renewcommand{\trAlgebraicallyDependent}{\BookTerm{T0401}{\csname BookOriginal@trAlgebraicallyDependent\endcsname}}
\expandafter\def\csname BookIndex@T0402\endcsname{hechengyu@合成域（compositum）}
\expandafter\let\csname BookOriginal@trCompositum\endcsname\trCompositum
\renewcommand{\trCompositum}{\BookTerm{T0402}{\csname BookOriginal@trCompositum\endcsname}}
\expandafter\def\csname BookIndex@T0403\endcsname{chunchaoyue@纯超越（purely transcendental）}
\expandafter\let\csname BookOriginal@trPurelyTranscendental\endcsname\trPurelyTranscendental
\renewcommand{\trPurelyTranscendental}{\BookTerm{T0403}{\csname BookOriginal@trPurelyTranscendental\endcsname}}
\expandafter\def\csname BookIndex@T0404\endcsname{lianshi@链式（catenary）}
\expandafter\let\csname BookOriginal@trCatenary\endcsname\trCatenary
\renewcommand{\trCatenary}{\BookTerm{T0404}{\csname BookOriginal@trCatenary\endcsname}}
\expandafter\def\csname BookIndex@T0405\endcsname{banlianxuxing@半连续性（semicontinuity）}
\expandafter\let\csname BookOriginal@trSemicontinuity\endcsname\trSemicontinuity
\renewcommand{\trSemicontinuity}{\BookTerm{T0405}{\csname BookOriginal@trSemicontinuity\endcsname}}
\expandafter\def\csname BookIndex@T0406\endcsname{shangbanlianxu@上半连续（upper semicontinuous）}
\expandafter\let\csname BookOriginal@trUpperSemicontinuous\endcsname\trUpperSemicontinuous
\renewcommand{\trUpperSemicontinuous}{\BookTerm{T0406}{\csname BookOriginal@trUpperSemicontinuous\endcsname}}
\expandafter\def\csname BookIndex@T0407\endcsname{xiabanlianxu@下半连续（lower semicontinuous）}
\expandafter\let\csname BookOriginal@trLowerSemicontinuous\endcsname\trLowerSemicontinuous
\renewcommand{\trLowerSemicontinuous}{\BookTerm{T0407}{\csname BookOriginal@trLowerSemicontinuous\endcsname}}
\expandafter\def\csname BookIndex@T0408\endcsname{guanlianduiying@关联对应（incidence correspondence）}
\expandafter\let\csname BookOriginal@trIncidenceCorrespondence\endcsname\trIncidenceCorrespondence
\renewcommand{\trIncidenceCorrespondence}{\BookTerm{T0408}{\csname BookOriginal@trIncidenceCorrespondence\endcsname}}
\expandafter\def\csname BookIndex@T0409\endcsname{guanliancu@关联簇（incidence variety）}
\expandafter\let\csname BookOriginal@trIncidenceVariety\endcsname\trIncidenceVariety
\renewcommand{\trIncidenceVariety}{\BookTerm{T0409}{\csname BookOriginal@trIncidenceVariety\endcsname}}
\expandafter\def\csname BookIndex@T0410\endcsname{zhengguibibao@正规闭包（normal closure）}
\expandafter\let\csname BookOriginal@trNormalClosure\endcsname\trNormalClosure
\renewcommand{\trNormalClosure}{\BookTerm{T0410}{\csname BookOriginal@trNormalClosure\endcsname}}
\expandafter\def\csname BookIndex@T0411\endcsname{yakebijuzhen@雅可比矩阵（Jacobian matrix）}
\expandafter\let\csname BookOriginal@trJacobianMatrix\endcsname\trJacobianMatrix
\renewcommand{\trJacobianMatrix}{\BookTerm{T0411}{\csname BookOriginal@trJacobianMatrix\endcsname}}
\expandafter\def\csname BookIndex@T0412\endcsname{feichangfan@非常泛（very general）}
\expandafter\let\csname BookOriginal@trVeryGeneral\endcsname\trVeryGeneral
\renewcommand{\trVeryGeneral}{\BookTerm{T0412}{\csname BookOriginal@trVeryGeneral\endcsname}}
\expandafter\def\csname BookIndex@T0413\endcsname{zhulixiangdingli@主理想定理（principal ideal theorem）}
\expandafter\let\csname BookOriginal@trPrincipalIdealTheorem\endcsname\trPrincipalIdealTheorem
\renewcommand{\trPrincipalIdealTheorem}{\BookTerm{T0413}{\csname BookOriginal@trPrincipalIdealTheorem\endcsname}}
\expandafter\def\csname BookIndex@T0414\endcsname{canshuxi@参数系（system of parameters）}
\expandafter\let\csname BookOriginal@trSystemOfParameters\endcsname\trSystemOfParameters
\renewcommand{\trSystemOfParameters}{\BookTerm{T0414}{\csname BookOriginal@trSystemOfParameters\endcsname}}
\expandafter\def\csname BookIndex@T0415\endcsname{fuhaomi@符号幂（symbolic power）}
\expandafter\let\csname BookOriginal@trSymbolicPower\endcsname\trSymbolicPower
\renewcommand{\trSymbolicPower}{\BookTerm{T0415}{\csname BookOriginal@trSymbolicPower\endcsname}}
\expandafter\def\csname BookIndex@T0416\endcsname{danyouli@单有理（uniruled）}
\expandafter\let\csname BookOriginal@trUniruled\endcsname\trUniruled
\renewcommand{\trUniruled}{\BookTerm{T0416}{\csname BookOriginal@trUniruled\endcsname}}
\expandafter\def\csname BookIndex@T0417\endcsname{suyuan@素元（prime element）}
\expandafter\let\csname BookOriginal@trPrimeElement\endcsname\trPrimeElement
\renewcommand{\trPrimeElement}{\BookTerm{T0417}{\csname BookOriginal@trPrimeElement\endcsname}}
\expandafter\def\csname BookIndex@T0418\endcsname{yuqiexiangliang@余切向量（cotangent vector）}
\expandafter\let\csname BookOriginal@trCotangentVector\endcsname\trCotangentVector
\renewcommand{\trCotangentVector}{\BookTerm{T0418}{\csname BookOriginal@trCotangentVector\endcsname}}
\expandafter\def\csname BookIndex@T0419\endcsname{daozi@导子（derivation）}
\expandafter\let\csname BookOriginal@trDerivation\endcsname\trDerivation
\renewcommand{\trDerivation}{\BookTerm{T0419}{\csname BookOriginal@trDerivation\endcsname}}
\expandafter\def\csname BookIndex@T0420\endcsname{qiexian@切线（tangent line）}
\expandafter\let\csname BookOriginal@trTangentLine\endcsname\trTangentLine
\renewcommand{\trTangentLine}{\BookTerm{T0420}{\csname BookOriginal@trTangentLine\endcsname}}
\expandafter\def\csname BookIndex@T0421\endcsname{yuzhi@余秩（corank）}
\expandafter\let\csname BookOriginal@trCorank\endcsname\trCorank
\renewcommand{\trCorank}{\BookTerm{T0421}{\csname BookOriginal@trCorank\endcsname}}
\expandafter\def\csname BookIndex@T0422\endcsname{xiangbanfenci@相伴分次（associated graded）}
\expandafter\let\csname BookOriginal@trAssociatedGraded\endcsname\trAssociatedGraded
\renewcommand{\trAssociatedGraded}{\BookTerm{T0422}{\csname BookOriginal@trAssociatedGraded\endcsname}}
\expandafter\def\csname BookIndex@T0423\endcsname{yakebipanbiefa@雅可比判别法（Jacobian criterion）}
\expandafter\let\csname BookOriginal@trJacobianCriterion\endcsname\trJacobianCriterion
\renewcommand{\trJacobianCriterion}{\BookTerm{T0423}{\csname BookOriginal@trJacobianCriterion\endcsname}}
\expandafter\def\csname BookIndex@T0424\endcsname{feiqiyi@非奇异（nonsingular）}
\expandafter\let\csname BookOriginal@trNonsingular\endcsname\trNonsingular
\renewcommand{\trNonsingular}{\BookTerm{T0424}{\csname BookOriginal@trNonsingular\endcsname}}
\expandafter\def\csname BookIndex@T0425\endcsname{qiyi@奇异（singular）}
\expandafter\let\csname BookOriginal@trSingular\endcsname\trSingular
\renewcommand{\trSingular}{\BookTerm{T0425}{\csname BookOriginal@trSingular\endcsname}}
\expandafter\def\csname BookIndex@T0426\endcsname{qidian@奇点（singular point）}
\expandafter\let\csname BookOriginal@trSingularPoint\endcsname\trSingularPoint
\renewcommand{\trSingularPoint}{\BookTerm{T0426}{\csname BookOriginal@trSingularPoint\endcsname}}
\expandafter\def\csname BookIndex@T0427\endcsname{hunhetezheng@混合特征（mixed characteristic）}
\expandafter\let\csname BookOriginal@trMixedCharacteristic\endcsname\trMixedCharacteristic
\renewcommand{\trMixedCharacteristic}{\BookTerm{T0427}{\csname BookOriginal@trMixedCharacteristic\endcsname}}
\expandafter\def\csname BookIndex@T0428\endcsname{qiepingmian@切平面（tangent plane）}
\expandafter\let\csname BookOriginal@trTangentPlane\endcsname\trTangentPlane
\renewcommand{\trTangentPlane}{\BookTerm{T0428}{\csname BookOriginal@trTangentPlane\endcsname}}
\expandafter\def\csname BookIndex@T0429\endcsname{toushe@投射（projection /\allowbreak{} projective）}
\expandafter\let\csname BookOriginal@trProjection\endcsname\trProjection
\renewcommand{\trProjection}{\BookTerm{T0429}{\csname BookOriginal@trProjection\endcsname}}
\expandafter\def\csname BookIndex@T0430\endcsname{duioucu@对偶簇（dual variety）}
\expandafter\let\csname BookOriginal@trDualVariety\endcsname\trDualVariety
\renewcommand{\trDualVariety}{\BookTerm{T0430}{\csname BookOriginal@trDualVariety\endcsname}}
\expandafter\def\csname BookIndex@T0431\endcsname{duiousheyingkongjian@对偶射影空间（dual projective space）}
\expandafter\let\csname BookOriginal@trDualProjectiveSpace\endcsname\trDualProjectiveSpace
\renewcommand{\trDualProjectiveSpace}{\BookTerm{T0431}{\csname BookOriginal@trDualProjectiveSpace\endcsname}}
\expandafter\def\csname BookIndex@T0432\endcsname{wanyouchaopingmian@万有超平面（universal hyperplane）}
\expandafter\let\csname BookOriginal@trUniversalHyperplane\endcsname\trUniversalHyperplane
\renewcommand{\trUniversalHyperplane}{\BookTerm{T0432}{\csname BookOriginal@trUniversalHyperplane\endcsname}}
\expandafter\def\csname BookIndex@T0433\endcsname{danzhihuayuan@单值化元（uniformizer）}
\expandafter\let\csname BookOriginal@trUniformizer\endcsname\trUniformizer
\renewcommand{\trUniformizer}{\BookTerm{T0433}{\csname BookOriginal@trUniformizer\endcsname}}
\expandafter\def\csname BookIndex@T0434\endcsname{fuzhi@赋值（valuation）}
\expandafter\let\csname BookOriginal@trValuation\endcsname\trValuation
\renewcommand{\trValuation}{\BookTerm{T0434}{\csname BookOriginal@trValuation\endcsname}}
\expandafter\def\csname BookIndex@T0435\endcsname{lingdian@零点（zeroes）}
\expandafter\let\csname BookOriginal@trZeros\endcsname\trZeros
\renewcommand{\trZeros}{\BookTerm{T0435}{\csname BookOriginal@trZeros\endcsname}}
\expandafter\def\csname BookIndex@T0436\endcsname{daidejinzhenghuan@戴德金整环（dedekind domain）}
\expandafter\let\csname BookOriginal@trDedekindDomain\endcsname\trDedekindDomain
\renewcommand{\trDedekindDomain}{\BookTerm{T0436}{\csname BookOriginal@trDedekindDomain\endcsname}}
\expandafter\def\csname BookIndex@T0437\endcsname{zhuanyihanshu@转移函数（transition function）}
\expandafter\let\csname BookOriginal@trTransitionFunction\endcsname\trTransitionFunction
\renewcommand{\trTransitionFunction}{\BookTerm{T0437}{\csname BookOriginal@trTransitionFunction\endcsname}}
\expandafter\def\csname BookIndex@T0438\endcsname{pingfanhua@平凡化（trivialization）}
\expandafter\let\csname BookOriginal@trTrivialization\endcsname\trTrivialization
\renewcommand{\trTrivialization}{\BookTerm{T0438}{\csname BookOriginal@trTrivialization\endcsname}}
\expandafter\def\csname BookIndex@T0439\endcsname{pingfan@平凡（trivial）}
\expandafter\let\csname BookOriginal@trTrivial\endcsname\trTrivial
\renewcommand{\trTrivial}{\BookTerm{T0439}{\csname BookOriginal@trTrivial\endcsname}}
\expandafter\def\csname BookIndex@T0440\endcsname{jubuziyou@局部自由（locally free）}
\expandafter\let\csname BookOriginal@trLocallyFree\endcsname\trLocallyFree
\renewcommand{\trLocallyFree}{\BookTerm{T0440}{\csname BookOriginal@trLocallyFree\endcsname}}
\expandafter\def\csname BookIndex@T0441\endcsname{fenshilixiang@分式理想（fractional ideal）}
\expandafter\let\csname BookOriginal@trFractionalIdeal\endcsname\trFractionalIdeal
\renewcommand{\trFractionalIdeal}{\BookTerm{T0441}{\csname BookOriginal@trFractionalIdeal\endcsname}}
\expandafter\def\csname BookIndex@T0442\endcsname{waidaishu@外代数（exterior algebra）}
\expandafter\let\csname BookOriginal@trExteriorAlgebra\endcsname\trExteriorAlgebra
\renewcommand{\trExteriorAlgebra}{\BookTerm{T0442}{\csname BookOriginal@trExteriorAlgebra\endcsname}}
\expandafter\def\csname BookIndex@T0443\endcsname{wanmeipeidui@完美配对（perfect pairing）}
\expandafter\let\csname BookOriginal@trPerfectPairing\endcsname\trPerfectPairing
\renewcommand{\trPerfectPairing}{\BookTerm{T0443}{\csname BookOriginal@trPerfectPairing\endcsname}}
\expandafter\def\csname BookIndex@T0444\endcsname{zifanceng@自反层（reflexive sheaf）}
\expandafter\let\csname BookOriginal@trReflexiveSheaf\endcsname\trReflexiveSheaf
\renewcommand{\trReflexiveSheaf}{\BookTerm{T0444}{\csname BookOriginal@trReflexiveSheaf\endcsname}}
\expandafter\def\csname BookIndex@T0445\endcsname{jiyingshe@迹映射（trace map）}
\expandafter\let\csname BookOriginal@trTraceMap\endcsname\trTraceMap
\renewcommand{\trTraceMap}{\BookTerm{T0445}{\csname BookOriginal@trTraceMap\endcsname}}
\expandafter\def\csname BookIndex@T0446\endcsname{baohehua@饱和化（saturation）}
\expandafter\let\csname BookOriginal@trSaturation\endcsname\trSaturation
\renewcommand{\trSaturation}{\BookTerm{T0446}{\csname BookOriginal@trSaturation\endcsname}}
\expandafter\def\csname BookIndex@T0447\endcsname{jidian@基点（base point）}
\expandafter\let\csname BookOriginal@trBasePoint\endcsname\trBasePoint
\renewcommand{\trBasePoint}{\BookTerm{T0447}{\csname BookOriginal@trBasePoint\endcsname}}
\expandafter\def\csname BookIndex@T0448\endcsname{jidianguiji@基点轨迹（base locus）}
\expandafter\let\csname BookOriginal@trBaseLocus\endcsname\trBaseLocus
\renewcommand{\trBaseLocus}{\BookTerm{T0448}{\csname BookOriginal@trBaseLocus\endcsname}}
\expandafter\def\csname BookIndex@T0449\endcsname{wujidian@无基点（base-point-free）}
\expandafter\let\csname BookOriginal@trBasePointFree\endcsname\trBasePointFree
\renewcommand{\trBasePointFree}{\BookTerm{T0449}{\csname BookOriginal@trBasePointFree\endcsname}}
\expandafter\def\csname BookIndex@T0450\endcsname{xianxingxi@线性系（linear system/\allowbreak{}series）}
\expandafter\let\csname BookOriginal@trLinearSystem\endcsname\trLinearSystem
\renewcommand{\trLinearSystem}{\BookTerm{T0450}{\csname BookOriginal@trLinearSystem\endcsname}}
\expandafter\def\csname BookIndex@T0451\endcsname{tuihua@退化（degenerate）}
\expandafter\let\csname BookOriginal@trDegenerate\endcsname\trDegenerate
\renewcommand{\trDegenerate}{\BookTerm{T0451}{\csname BookOriginal@trDegenerate\endcsname}}
\expandafter\def\csname BookIndex@T0452\endcsname{heji@盒积（box product）}
\expandafter\let\csname BookOriginal@trBoxProduct\endcsname\trBoxProduct
\renewcommand{\trBoxProduct}{\BookTerm{T0452}{\csname BookOriginal@trBoxProduct\endcsname}}
\expandafter\def\csname BookIndex@T0453\endcsname{gexian@割线（secant）}
\expandafter\let\csname BookOriginal@trSecant\endcsname\trSecant
\renewcommand{\trSecant}{\BookTerm{T0453}{\csname BookOriginal@trSecant\endcsname}}
\expandafter\def\csname BookIndex@T0454\endcsname{zhu@主（principal）}
\expandafter\let\csname BookOriginal@trPrincipal\endcsname\trPrincipal
\renewcommand{\trPrincipal}{\BookTerm{T0454}{\csname BookOriginal@trPrincipal\endcsname}}
\expandafter\def\csname BookIndex@T0455\endcsname{qiechu@切除（excision）}
\expandafter\let\csname BookOriginal@trExcision\endcsname\trExcision
\renewcommand{\trExcision}{\BookTerm{T0455}{\csname BookOriginal@trExcision\endcsname}}
\expandafter\def\csname BookIndex@T0456\endcsname{sheyingbianhuan@射影变换（projective transformation）}
\expandafter\let\csname BookOriginal@trProjectiveTransformation\endcsname\trProjectiveTransformation
\renewcommand{\trProjectiveTransformation}{\BookTerm{T0456}{\csname BookOriginal@trProjectiveTransformation\endcsname}}
\expandafter\def\csname BookIndex@T0457\endcsname{fenshixianxingbianhuan@分式线性变换（fractional linear transformation）}
\expandafter\let\csname BookOriginal@trFractionalLinearTransformation\endcsname\trFractionalLinearTransformation
\renewcommand{\trFractionalLinearTransformation}{\BookTerm{T0457}{\csname BookOriginal@trFractionalLinearTransformation\endcsname}}
\expandafter\def\csname BookIndex@T0458\endcsname{wendingziqun@稳定子群（stabilizer）}
\expandafter\let\csname BookOriginal@trStabilizer\endcsname\trStabilizer
\renewcommand{\trStabilizer}{\BookTerm{T0458}{\csname BookOriginal@trStabilizer\endcsname}}
\expandafter\def\csname BookIndex@T0459\endcsname{zhengguicong@正规丛（normal bundle）}
\expandafter\let\csname BookOriginal@trNormalBundle\endcsname\trNormalBundle
\renewcommand{\trNormalBundle}{\BookTerm{T0459}{\csname BookOriginal@trNormalBundle\endcsname}}
\expandafter\def\csname BookIndex@T0460\endcsname{baohe@饱和（saturated）}
\expandafter\let\csname BookOriginal@trSaturated\endcsname\trSaturated
\renewcommand{\trSaturated}{\BookTerm{T0460}{\csname BookOriginal@trSaturated\endcsname}}
\expandafter\def\csname BookIndex@T0461\endcsname{zhengtishengchengde@整体生成的（globally generated）}
\expandafter\let\csname BookOriginal@trGloballyGenerated\endcsname\trGloballyGenerated
\renewcommand{\trGloballyGenerated}{\BookTerm{T0461}{\csname BookOriginal@trGloballyGenerated\endcsname}}
\expandafter\def\csname BookIndex@T0462\endcsname{fengfuxing@丰富性（ampleness）}
\expandafter\let\csname BookOriginal@trAmpleness\endcsname\trAmpleness
\renewcommand{\trAmpleness}{\BookTerm{T0462}{\csname BookOriginal@trAmpleness\endcsname}}
\expandafter\def\csname BookIndex@T0463\endcsname{hechong@合冲（syzygy）}
\expandafter\let\csname BookOriginal@trSyzygy\endcsname\trSyzygy
\renewcommand{\trSyzygy}{\BookTerm{T0463}{\csname BookOriginal@trSyzygy\endcsname}}
\expandafter\def\csname BookIndex@T0464\endcsname{zhengtishengcheng@整体生成（global generation）}
\expandafter\let\csname BookOriginal@trGlobalGeneration\endcsname\trGlobalGeneration
\renewcommand{\trGlobalGeneration}{\BookTerm{T0464}{\csname BookOriginal@trGlobalGeneration\endcsname}}
\expandafter\def\csname BookIndex@T0465\endcsname{banfengfu@半丰富（semiample）}
\expandafter\let\csname BookOriginal@trSemiample\endcsname\trSemiample
\renewcommand{\trSemiample}{\BookTerm{T0465}{\csname BookOriginal@trSemiample\endcsname}}
\expandafter\def\csname BookIndex@T0466\endcsname{jifengfuxing@极丰富性（very ampleness）}
\expandafter\let\csname BookOriginal@trVeryAmpleness\endcsname\trVeryAmpleness
\renewcommand{\trVeryAmpleness}{\BookTerm{T0466}{\csname BookOriginal@trVeryAmpleness\endcsname}}
\expandafter\def\csname BookIndex@T0467\endcsname{emiteduliang@厄米特度量（Hermitian metric）}
\expandafter\let\csname BookOriginal@trHermitianMetric\endcsname\trHermitianMetric
\renewcommand{\trHermitianMetric}{\BookTerm{T0467}{\csname BookOriginal@trHermitianMetric\endcsname}}
\expandafter\def\csname BookIndex@T0468\endcsname{zilixiangceng@子理想层（subideal sheaf）}
\expandafter\let\csname BookOriginal@trSubidealSheaf\endcsname\trSubidealSheaf
\renewcommand{\trSubidealSheaf}{\BookTerm{T0468}{\csname BookOriginal@trSubidealSheaf\endcsname}}
\expandafter\def\csname BookIndex@T0469\endcsname{zhumo@主模（principal module）}
\expandafter\let\csname BookOriginal@trPrincipalModule\endcsname\trPrincipalModule
\renewcommand{\trPrincipalModule}{\BookTerm{T0469}{\csname BookOriginal@trPrincipalModule\endcsname}}
\expandafter\def\csname BookIndex@T0470\endcsname{zhongyanshangcong@重言商丛（tautological quotient bundle）}
\expandafter\let\csname BookOriginal@trTautologicalQuotientBundle\endcsname\trTautologicalQuotientBundle
\renewcommand{\trTautologicalQuotientBundle}{\BookTerm{T0470}{\csname BookOriginal@trTautologicalQuotientBundle\endcsname}}
\expandafter\def\csname BookIndex@T0471\endcsname{zhongyanzicong@重言子丛（tautological subbundle）}
\expandafter\let\csname BookOriginal@trTautologicalSubbundle\endcsname\trTautologicalSubbundle
\renewcommand{\trTautologicalSubbundle}{\BookTerm{T0471}{\csname BookOriginal@trTautologicalSubbundle\endcsname}}
\expandafter\def\csname BookIndex@T0472\endcsname{fanzhenghelie@泛正合列（universal exact sequence）}
\expandafter\let\csname BookOriginal@trUniversalExactSequence\endcsname\trUniversalExactSequence
\renewcommand{\trUniversalExactSequence}{\BookTerm{T0472}{\csname BookOriginal@trUniversalExactSequence\endcsname}}
\expandafter\def\csname BookIndex@T0473\endcsname{gelasimancong@格拉斯曼丛（Grassmannian bundle）}
\expandafter\let\csname BookOriginal@trGrassmannianBundle\endcsname\trGrassmannianBundle
\renewcommand{\trGrassmannianBundle}{\BookTerm{T0473}{\csname BookOriginal@trGrassmannianBundle\endcsname}}
\expandafter\def\csname BookIndex@T0474\endcsname{ercichaoqumian@二次超曲面（quadric hypersurface）}
\expandafter\let\csname BookOriginal@trQuadricHypersurface\endcsname\trQuadricHypersurface
\renewcommand{\trQuadricHypersurface}{\BookTerm{T0474}{\csname BookOriginal@trQuadricHypersurface\endcsname}}
\expandafter\def\csname BookIndex@T0475\endcsname{chunzhangliang@纯张量（pure tensor）}
\expandafter\let\csname BookOriginal@trPureTensor\endcsname\trPureTensor
\renewcommand{\trPureTensor}{\BookTerm{T0475}{\csname BookOriginal@trPureTensor\endcsname}}
\expandafter\def\csname BookIndex@T0476\endcsname{daishuceng@代数层（sheaf of algebras）}
\expandafter\let\csname BookOriginal@trSheafOfAlgebras\endcsname\trSheafOfAlgebras
\renewcommand{\trSheafOfAlgebras}{\BookTerm{T0476}{\csname BookOriginal@trSheafOfAlgebras\endcsname}}
\expandafter\def\csname BookIndex@T0477\endcsname{fencidaishuceng@分次代数层（sheaf of graded algebras）}
\expandafter\let\csname BookOriginal@trSheafOfGradedAlgebras\endcsname\trSheafOfGradedAlgebras
\renewcommand{\trSheafOfGradedAlgebras}{\BookTerm{T0477}{\csname BookOriginal@trSheafOfGradedAlgebras\endcsname}}
\expandafter\def\csname BookIndex@T0478\endcsname{quankongjian@全空间（total space）}
\expandafter\let\csname BookOriginal@trTotalSpace\endcsname\trTotalSpace
\renewcommand{\trTotalSpace}{\BookTerm{T0478}{\csname BookOriginal@trTotalSpace\endcsname}}
\expandafter\def\csname BookIndex@T0479\endcsname{zhongyancong@重言丛（tautological bundle）}
\expandafter\let\csname BookOriginal@trTautologicalBundle\endcsname\trTautologicalBundle
\renewcommand{\trTautologicalBundle}{\BookTerm{T0479}{\csname BookOriginal@trTautologicalBundle\endcsname}}
\expandafter\def\csname BookIndex@T0480\endcsname{xicebuluhequmian@希策布鲁赫曲面（Hirzebruch surface）}
\expandafter\let\csname BookOriginal@trHirzebruchSurface\endcsname\trHirzebruchSurface
\renewcommand{\trHirzebruchSurface}{\BookTerm{T0480}{\csname BookOriginal@trHirzebruchSurface\endcsname}}
\expandafter\def\csname BookIndex@T0481\endcsname{nifangshe@拟仿射（quasiaffine）}
\expandafter\let\csname BookOriginal@trQuasiaffine\endcsname\trQuasiaffine
\renewcommand{\trQuasiaffine}{\BookTerm{T0481}{\csname BookOriginal@trQuasiaffine\endcsname}}
\expandafter\def\csname BookIndex@T0482\endcsname{nitonggou@拟同构（quasiisomorphism）}
\expandafter\let\csname BookOriginal@trQuasiisomorphism\endcsname\trQuasiisomorphism
\renewcommand{\trQuasiisomorphism}{\BookTerm{T0482}{\csname BookOriginal@trQuasiisomorphism\endcsname}}
\expandafter\def\csname BookIndex@T0483\endcsname{daochufanchou@导出范畴（derived category）}
\expandafter\let\csname BookOriginal@trDerivedCategory\endcsname\trDerivedCategory
\renewcommand{\trDerivedCategory}{\BookTerm{T0483}{\csname BookOriginal@trDerivedCategory\endcsname}}
\expandafter\def\csname BookIndex@T0484\endcsname{duiouhuaceng@对偶化层（dualizing sheaf）}
\expandafter\let\csname BookOriginal@trDualizingSheaf\endcsname\trDualizingSheaf
\renewcommand{\trDualizingSheaf}{\BookTerm{T0484}{\csname BookOriginal@trDualizingSheaf\endcsname}}
\expandafter\def\csname BookIndex@T0485\endcsname{laurentduoxiangshi@Laurent 多项式（Laurent polynomial）}
\expandafter\let\csname BookOriginal@trLaurentPolynomial\endcsname\trLaurentPolynomial
\renewcommand{\trLaurentPolynomial}{\BookTerm{T0485}{\csname BookOriginal@trLaurentPolynomial\endcsname}}
\expandafter\def\csname BookIndex@T0486\endcsname{oulashixingshu@欧拉示性数（Euler characteristic）}
\expandafter\let\csname BookOriginal@trEulerCharacteristic\endcsname\trEulerCharacteristic
\renewcommand{\trEulerCharacteristic}{\BookTerm{T0486}{\csname BookOriginal@trEulerCharacteristic\endcsname}}
\expandafter\def\csname BookIndex@T0487\endcsname{suanshukuige@算术亏格（arithmetic genus）}
\expandafter\let\csname BookOriginal@trArithmeticGenus\endcsname\trArithmeticGenus
\renewcommand{\trArithmeticGenus}{\BookTerm{T0487}{\csname BookOriginal@trArithmeticGenus\endcsname}}
\expandafter\def\csname BookIndex@T0488\endcsname{shuzhiyouxiao@数值有效（nef）}
\expandafter\let\csname BookOriginal@trNef\endcsname\trNef
\renewcommand{\trNef}{\BookTerm{T0488}{\csname BookOriginal@trNef\endcsname}}
\expandafter\def\csname BookIndex@T0489\endcsname{shuzhipingfan@数值平凡（numerically trivial）}
\expandafter\let\csname BookOriginal@trNumericallyTrivial\endcsname\trNumericallyTrivial
\renewcommand{\trNumericallyTrivial}{\BookTerm{T0489}{\csname BookOriginal@trNumericallyTrivial\endcsname}}
\expandafter\def\csname BookIndex@T0490\endcsname{shuzhidengjia@数值等价（numerically equivalent）}
\expandafter\let\csname BookOriginal@trNumericallyEquivalent\endcsname\trNumericallyEquivalent
\renewcommand{\trNumericallyEquivalent}{\BookTerm{T0490}{\csname BookOriginal@trNumericallyEquivalent\endcsname}}
\expandafter\def\csname BookIndex@T0491\endcsname{fangshexing@仿射性（affineness）}
\expandafter\let\csname BookOriginal@trAffineness\endcsname\trAffineness
\renewcommand{\trAffineness}{\BookTerm{T0491}{\csname BookOriginal@trAffineness\endcsname}}
\expandafter\def\csname BookIndex@T0492\endcsname{yuanwang@愿望（Desideratum）}
\expandafter\let\csname BookOriginal@trDesideratum\endcsname\trDesideratum
\renewcommand{\trDesideratum}{\BookTerm{T0492}{\csname BookOriginal@trDesideratum\endcsname}}
\expandafter\def\csname BookIndex@T0493\endcsname{daishutijixingshi@代数体积形式（algebraic volume form）}
\expandafter\let\csname BookOriginal@trAlgebraicVolumeForm\endcsname\trAlgebraicVolumeForm
\renewcommand{\trAlgebraicVolumeForm}{\BookTerm{T0493}{\csname BookOriginal@trAlgebraicVolumeForm\endcsname}}
\expandafter\def\csname BookIndex@T0494\endcsname{lingceng@零层（zero-sheaf）}
\expandafter\let\csname BookOriginal@trZeroSheaf\endcsname\trZeroSheaf
\renewcommand{\trZeroSheaf}{\BookTerm{T0494}{\csname BookOriginal@trZeroSheaf\endcsname}}
\expandafter\def\csname BookIndex@T0495\endcsname{chaopingmianlei@超平面类（hyperplane class）}
\expandafter\let\csname BookOriginal@trHyperplaneClass\endcsname\trHyperplaneClass
\renewcommand{\trHyperplaneClass}{\BookTerm{T0495}{\csname BookOriginal@trHyperplaneClass\endcsname}}
\expandafter\def\csname BookIndex@T0496\endcsname{tapukuige@拓扑亏格（topological genus）}
\expandafter\let\csname BookOriginal@trTopologicalGenus\endcsname\trTopologicalGenus
\renewcommand{\trTopologicalGenus}{\BookTerm{T0496}{\csname BookOriginal@trTopologicalGenus\endcsname}}
\expandafter\def\csname BookIndex@T0497\endcsname{gaojieqiantui@高阶前推（higher pushforward）}
\expandafter\let\csname BookOriginal@trHigherPushforward\endcsname\trHigherPushforward
\renewcommand{\trHigherPushforward}{\BookTerm{T0497}{\csname BookOriginal@trHigherPushforward\endcsname}}
\expandafter\def\csname BookIndex@T0498\endcsname{gaojiezhixiang@高阶直像（higher direct image）}
\expandafter\let\csname BookOriginal@trHigherDirectImage\endcsname\trHigherDirectImage
\renewcommand{\trHigherDirectImage}{\BookTerm{T0498}{\csname BookOriginal@trHigherDirectImage\endcsname}}
\expandafter\def\csname BookIndex@T0499\endcsname{chaotuoyuan@超椭圆（hyperelliptic）}
\expandafter\let\csname BookOriginal@trHyperelliptic\endcsname\trHyperelliptic
\renewcommand{\trHyperelliptic}{\BookTerm{T0499}{\csname BookOriginal@trHyperelliptic\endcsname}}
\expandafter\def\csname BookIndex@T0500\endcsname{fenzhidian@分支点（branch point）}
\expandafter\let\csname BookOriginal@trBranchPoint\endcsname\trBranchPoint
\renewcommand{\trBranchPoint}{\BookTerm{T0500}{\csname BookOriginal@trBranchPoint\endcsname}}
\expandafter\def\csname BookIndex@T0501\endcsname{fenqi@分歧（ramification）}
\expandafter\let\csname BookOriginal@trRamification\endcsname\trRamification
\renewcommand{\trRamification}{\BookTerm{T0501}{\csname BookOriginal@trRamification\endcsname}}
\expandafter\def\csname BookIndex@T0502\endcsname{erzhongfugai@二重覆盖（double cover）}
\expandafter\let\csname BookOriginal@trDoubleCover\endcsname\trDoubleCover
\renewcommand{\trDoubleCover}{\BookTerm{T0502}{\csname BookOriginal@trDoubleCover\endcsname}}
\expandafter\def\csname BookIndex@T0503\endcsname{duihe@对合（involution）}
\expandafter\let\csname BookOriginal@trInvolution\endcsname\trInvolution
\renewcommand{\trInvolution}{\BookTerm{T0503}{\csname BookOriginal@trInvolution\endcsname}}
\expandafter\def\csname BookIndex@T0504\endcsname{pingmiansiciquxian@平面四次曲线（plane quartic）}
\expandafter\let\csname BookOriginal@trPlaneQuartic\endcsname\trPlaneQuartic
\renewcommand{\trPlaneQuartic}{\BookTerm{T0504}{\csname BookOriginal@trPlaneQuartic\endcsname}}
\expandafter\def\csname BookIndex@T0505\endcsname{sanjiao@三角（trigonal）}
\expandafter\let\csname BookOriginal@trTrigonal\endcsname\trTrigonal
\renewcommand{\trTrigonal}{\BookTerm{T0505}{\csname BookOriginal@trTrigonal\endcsname}}
\expandafter\def\csname BookIndex@T0506\endcsname{jiaobi@交比（cross-ratio）}
\expandafter\let\csname BookOriginal@trCrossRatio\endcsname\trCrossRatio
\renewcommand{\trCrossRatio}{\BookTerm{T0506}{\csname BookOriginal@trCrossRatio\endcsname}}
\expandafter\def\csname BookIndex@T0507\endcsname{guai@拐（flex）}
\expandafter\let\csname BookOriginal@trFlex\endcsname\trFlex
\renewcommand{\trFlex}{\BookTerm{T0507}{\csname BookOriginal@trFlex\endcsname}}
\expandafter\def\csname BookIndex@T0508\endcsname{weiersitelasifanshi@魏尔斯特拉斯范式（Weierstrass normal form）}
\expandafter\let\csname BookOriginal@trWeierstrassNormalForm\endcsname\trWeierstrassNormalForm
\renewcommand{\trWeierstrassNormalForm}{\BookTerm{T0508}{\csname BookOriginal@trWeierstrassNormalForm\endcsname}}
\expandafter\def\csname BookIndex@T0509\endcsname{moxingshi@模形式（modular form）}
\expandafter\let\csname BookOriginal@trModularForm\endcsname\trModularForm
\renewcommand{\trModularForm}{\BookTerm{T0509}{\csname BookOriginal@trModularForm\endcsname}}
\expandafter\def\csname BookIndex@T0510\endcsname{moquxian@模曲线（modular curve）}
\expandafter\let\csname BookOriginal@trModularCurve\endcsname\trModularCurve
\renewcommand{\trModularCurve}{\BookTerm{T0510}{\csname BookOriginal@trModularCurve\endcsname}}
\expandafter\def\csname BookIndex@T0511\endcsname{erzhongpingmian@二重平面（double plane）}
\expandafter\let\csname BookOriginal@trDoublePlane\endcsname\trDoublePlane
\renewcommand{\trDoublePlane}{\BookTerm{T0511}{\csname BookOriginal@trDoublePlane\endcsname}}
\expandafter\def\csname BookIndex@T0512\endcsname{qunlu@群律（group law）}
\expandafter\let\csname BookOriginal@trGroupLaw\endcsname\trGroupLaw
\renewcommand{\trGroupLaw}{\BookTerm{T0512}{\csname BookOriginal@trGroupLaw\endcsname}}
\expandafter\def\csname BookIndex@T0513\endcsname{xiangjiaoji@相交积（intersection product）}
\expandafter\let\csname BookOriginal@trIntersectionProduct\endcsname\trIntersectionProduct
\renewcommand{\trIntersectionProduct}{\BookTerm{T0513}{\csname BookOriginal@trIntersectionProduct\endcsname}}
\expandafter\def\csname BookIndex@T0514\endcsname{xiangjiaoshu@相交数（intersection number）}
\expandafter\let\csname BookOriginal@trIntersectionNumber\endcsname\trIntersectionNumber
\renewcommand{\trIntersectionNumber}{\BookTerm{T0514}{\csname BookOriginal@trIntersectionNumber\endcsname}}
\expandafter\def\csname BookIndex@T0515\endcsname{xiangjiao@相交（intersection）}
\expandafter\let\csname BookOriginal@trIntersection\endcsname\trIntersection
\renewcommand{\trIntersection}{\BookTerm{T0515}{\csname BookOriginal@trIntersection\endcsname}}
\expandafter\def\csname BookIndex@T0516\endcsname{xianweihua@纤维化（fibration）}
\expandafter\let\csname BookOriginal@trFibration\endcsname\trFibration
\renewcommand{\trFibration}{\BookTerm{T0516}{\csname BookOriginal@trFibration\endcsname}}
\expandafter\def\csname BookIndex@T0517\endcsname{zhenbianhuan@真变换（proper transform）}
\expandafter\let\csname BookOriginal@trProperTransform\endcsname\trProperTransform
\renewcommand{\trProperTransform}{\BookTerm{T0517}{\csname BookOriginal@trProperTransform\endcsname}}
\expandafter\def\csname BookIndex@T0518\endcsname{zixiangjiao@自相交（self-intersection）}
\expandafter\let\csname BookOriginal@trSelfIntersection\endcsname\trSelfIntersection
\renewcommand{\trSelfIntersection}{\BookTerm{T0518}{\csname BookOriginal@trSelfIntersection\endcsname}}
\expandafter\def\csname BookIndex@T0519\endcsname{xiangjiaoxing@相交型（intersection form）}
\expandafter\let\csname BookOriginal@trIntersectionForm\endcsname\trIntersectionForm
\renewcommand{\trIntersectionForm}{\BookTerm{T0519}{\csname BookOriginal@trIntersectionForm\endcsname}}
\expandafter\def\csname BookIndex@T0520\endcsname{bansuigongshi@伴随公式（adjunction formula）}
\expandafter\let\csname BookOriginal@trAdjunctionFormula\endcsname\trAdjunctionFormula
\renewcommand{\trAdjunctionFormula}{\BookTerm{T0520}{\csname BookOriginal@trAdjunctionFormula\endcsname}}
\expandafter\def\csname BookIndex@T0521\endcsname{fuhaocha@符号差（signature）}
\expandafter\let\csname BookOriginal@trSignature\endcsname\trSignature
\renewcommand{\trSignature}{\BookTerm{T0521}{\csname BookOriginal@trSignature\endcsname}}
\expandafter\def\csname BookIndex@T0522\endcsname{yidongyinli@移动引理（moving lemma）}
\expandafter\let\csname BookOriginal@trMovingLemma\endcsname\trMovingLemma
\renewcommand{\trMovingLemma}{\BookTerm{T0522}{\csname BookOriginal@trMovingLemma\endcsname}}
\expandafter\def\csname BookIndex@T0523\endcsname{yushuipingluguo@余水平滤过（coniveau filtration）}
\expandafter\let\csname BookOriginal@trConiveauFiltration\endcsname\trConiveauFiltration
\renewcommand{\trConiveauFiltration}{\BookTerm{T0523}{\csname BookOriginal@trConiveauFiltration\endcsname}}
\expandafter\def\csname BookIndex@T0524\endcsname{youxiao@有效（effective）}
\expandafter\let\csname BookOriginal@trEffective\endcsname\trEffective
\renewcommand{\trEffective}{\BookTerm{T0524}{\csname BookOriginal@trEffective\endcsname}}
\expandafter\def\csname BookIndex@T0525\endcsname{yuqieceng@余切层（cotangent sheaf）}
\expandafter\let\csname BookOriginal@trCotangentSheaf\endcsname\trCotangentSheaf
\renewcommand{\trCotangentSheaf}{\BookTerm{T0525}{\csname BookOriginal@trCotangentSheaf\endcsname}}
\expandafter\def\csname BookIndex@T0526\endcsname{qieceng@切层（tangent sheaf）}
\expandafter\let\csname BookOriginal@trTangentSheaf\endcsname\trTangentSheaf
\renewcommand{\trTangentSheaf}{\BookTerm{T0526}{\csname BookOriginal@trTangentSheaf\endcsname}}
\expandafter\def\csname BookIndex@T0527\endcsname{yuzhengguimo@余正规模（conormal module）}
\expandafter\let\csname BookOriginal@trConormalSheaf\endcsname\trConormalSheaf
\renewcommand{\trConormalSheaf}{\BookTerm{T0527}{\csname BookOriginal@trConormalSheaf\endcsname}}
\expandafter\def\csname BookIndex@T0528\endcsname{yuzhenggui@余正规（conormal）}
\expandafter\let\csname BookOriginal@trConormal\endcsname\trConormal
\renewcommand{\trConormal}{\BookTerm{T0528}{\csname BookOriginal@trConormal\endcsname}}
\expandafter\def\csname BookIndex@T0529\endcsname{jihekuige@几何亏格（geometric genus）}
\expandafter\let\csname BookOriginal@trGeometricGenus\endcsname\trGeometricGenus
\renewcommand{\trGeometricGenus}{\BookTerm{T0529}{\csname BookOriginal@trGeometricGenus\endcsname}}
\expandafter\def\csname BookIndex@T0530\endcsname{qixingkongjian@齐性空间（homogeneous space）}
\expandafter\let\csname BookOriginal@trHomogeneousSpace\endcsname\trHomogeneousSpace
\renewcommand{\trHomogeneousSpace}{\BookTerm{T0530}{\csname BookOriginal@trHomogeneousSpace\endcsname}}
\expandafter\def\csname BookIndex@T0531\endcsname{fenqiguiji@分歧轨迹（ramification locus）}
\expandafter\let\csname BookOriginal@trRamificationLocus\endcsname\trRamificationLocus
\renewcommand{\trRamificationLocus}{\BookTerm{T0531}{\csname BookOriginal@trRamificationLocus\endcsname}}
\expandafter\def\csname BookIndex@T0532\endcsname{fenzhiguiji@分支轨迹（branch locus）}
\expandafter\let\csname BookOriginal@trBranchLocus\endcsname\trBranchLocus
\renewcommand{\trBranchLocus}{\BookTerm{T0532}{\csname BookOriginal@trBranchLocus\endcsname}}
\expandafter\def\csname BookIndex@T0533\endcsname{fenqichuzi@分歧除子（ramification divisor）}
\expandafter\let\csname BookOriginal@trRamificationDivisor\endcsname\trRamificationDivisor
\renewcommand{\trRamificationDivisor}{\BookTerm{T0533}{\csname BookOriginal@trRamificationDivisor\endcsname}}
\expandafter\def\csname BookIndex@T0534\endcsname{fenzhichuzi@分支除子（branch divisor）}
\expandafter\let\csname BookOriginal@trBranchDivisor\endcsname\trBranchDivisor
\renewcommand{\trBranchDivisor}{\BookTerm{T0534}{\csname BookOriginal@trBranchDivisor\endcsname}}
\expandafter\def\csname BookIndex@T0535\endcsname{fenqizhishu@分歧指数（ramification index）}
\expandafter\let\csname BookOriginal@trRamificationIndex\endcsname\trRamificationIndex
\renewcommand{\trRamificationIndex}{\BookTerm{T0535}{\csname BookOriginal@trRamificationIndex\endcsname}}
\expandafter\def\csname BookIndex@T0536\endcsname{xun@驯（tame）}
\expandafter\let\csname BookOriginal@trTame\endcsname\trTame
\renewcommand{\trTame}{\BookTerm{T0536}{\csname BookOriginal@trTame\endcsname}}
\expandafter\def\csname BookIndex@T0537\endcsname{ye@野（wild）}
\expandafter\let\csname BookOriginal@trWild\endcsname\trWild
\renewcommand{\trWild}{\BookTerm{T0537}{\csname BookOriginal@trWild\endcsname}}
\expandafter\def\csname BookIndex@T0538\endcsname{wufenzhi@无分支（unbranched）}
\expandafter\let\csname BookOriginal@trUnbranched\endcsname\trUnbranched
\renewcommand{\trUnbranched}{\BookTerm{T0538}{\csname BookOriginal@trUnbranched\endcsname}}
\expandafter\def\csname BookIndex@T0539\endcsname{yibankefen@一般可分（generically separable）}
\expandafter\let\csname BookOriginal@trGenericallySeparable\endcsname\trGenericallySeparable
\renewcommand{\trGenericallySeparable}{\BookTerm{T0539}{\csname BookOriginal@trGenericallySeparable\endcsname}}
\expandafter\def\csname BookIndex@T0540\endcsname{chajilixiang@差积理想（different ideal）}
\expandafter\let\csname BookOriginal@trDifferentIdeal\endcsname\trDifferentIdeal
\renewcommand{\trDifferentIdeal}{\BookTerm{T0540}{\csname BookOriginal@trDifferentIdeal\endcsname}}
\expandafter\def\csname BookIndex@T0541\endcsname{panbieshilixiang@判别式理想（discriminant ideal）}
\expandafter\let\csname BookOriginal@trDiscriminantIdeal\endcsname\trDiscriminantIdeal
\renewcommand{\trDiscriminantIdeal}{\BookTerm{T0541}{\csname BookOriginal@trDiscriminantIdeal\endcsname}}
\expandafter\def\csname BookIndex@T0542\endcsname{panbieshi@判别式（discriminant）}
\expandafter\let\csname BookOriginal@trDiscriminant\endcsname\trDiscriminant
\renewcommand{\trDiscriminant}{\BookTerm{T0542}{\csname BookOriginal@trDiscriminant\endcsname}}
\expandafter\def\csname BookIndex@T0543\endcsname{xiaojiexiao@小解消（small resolution）}
\expandafter\let\csname BookOriginal@trSmallResolution\endcsname\trSmallResolution
\renewcommand{\trSmallResolution}{\BookTerm{T0543}{\csname BookOriginal@trSmallResolution\endcsname}}
\expandafter\def\csname BookIndex@T0544\endcsname{guanzhuanglinyu@管状邻域（tubular neighborhood）}
\expandafter\let\csname BookOriginal@trTubularNeighborhood\endcsname\trTubularNeighborhood
\renewcommand{\trTubularNeighborhood}{\BookTerm{T0544}{\csname BookOriginal@trTubularNeighborhood\endcsname}}
\expandafter\def\csname BookIndex@T0545\endcsname{xianwei@纤维（fiber）}
\expandafter\let\csname BookOriginal@trFiber\endcsname\trFiber
\renewcommand{\trFiber}{\BookTerm{T0545}{\csname BookOriginal@trFiber\endcsname}}
\expandafter\def\csname BookIndex@T0546\endcsname{xianweiji@纤维积（fiber product）}
\expandafter\let\csname BookOriginal@trFiberProduct\endcsname\trFiberProduct
\renewcommand{\trFiberProduct}{\BookTerm{T0546}{\csname BookOriginal@trFiberProduct\endcsname}}
\expandafter\def\csname BookIndex@T0547\endcsname{jibianhuan@基变换（base change）}
\expandafter\let\csname BookOriginal@trBaseChange\endcsname\trBaseChange
\renewcommand{\trBaseChange}{\BookTerm{T0547}{\csname BookOriginal@trBaseChange\endcsname}}
\expandafter\def\csname BookIndex@T0548\endcsname{guyou@固有（proper）}
\expandafter\let\csname BookOriginal@trProper\endcsname\trProper
\renewcommand{\trProper}{\BookTerm{T0548}{\csname BookOriginal@trProper\endcsname}}
\expandafter\def\csname BookIndex@T0549\endcsname{fenli@分离（separated）}
\expandafter\let\csname BookOriginal@trSeparated\endcsname\trSeparated
\renewcommand{\trSeparated}{\BookTerm{T0549}{\csname BookOriginal@trSeparated\endcsname}}
\expandafter\def\csname BookIndex@T0550\endcsname{youxianxing@有限型（finite type）}
\expandafter\let\csname BookOriginal@trFiniteType\endcsname\trFiniteType
\renewcommand{\trFiniteType}{\BookTerm{T0550}{\csname BookOriginal@trFiniteType\endcsname}}
\expandafter\def\csname BookIndex@T0551\endcsname{jubu@局部（locally）}
\expandafter\let\csname BookOriginal@trLocally\endcsname\trLocally
\renewcommand{\trLocally}{\BookTerm{T0551}{\csname BookOriginal@trLocally\endcsname}}
\expandafter\def\csname BookIndex@T0552\endcsname{linyu@邻域（neighborhood）}
\expandafter\let\csname BookOriginal@trNeighborhood\endcsname\trNeighborhood
\renewcommand{\trNeighborhood}{\BookTerm{T0552}{\csname BookOriginal@trNeighborhood\endcsname}}
\expandafter\def\csname BookIndex@T0553\endcsname{fugai@覆盖（cover）}
\expandafter\let\csname BookOriginal@trCover\endcsname\trCover
\renewcommand{\trCover}{\BookTerm{T0553}{\csname BookOriginal@trCover\endcsname}}
\expandafter\def\csname BookIndex@T0554\endcsname{kaifugai@开覆盖（open cover）}
\expandafter\let\csname BookOriginal@trOpenCover\endcsname\trOpenCover
\renewcommand{\trOpenCover}{\BookTerm{T0554}{\csname BookOriginal@trOpenCover\endcsname}}
\expandafter\def\csname BookIndex@T0555\endcsname{lisantapu@离散拓扑（discrete topology）}
\expandafter\let\csname BookOriginal@trDiscreteTopology\endcsname\trDiscreteTopology
\renewcommand{\trDiscreteTopology}{\BookTerm{T0555}{\csname BookOriginal@trDiscreteTopology\endcsname}}
\expandafter\def\csname BookIndex@T0556\endcsname{jingdiantapu@经典拓扑（classical topology）}
\expandafter\let\csname BookOriginal@trClassicalTopology\endcsname\trClassicalTopology
\renewcommand{\trClassicalTopology}{\BookTerm{T0556}{\csname BookOriginal@trClassicalTopology\endcsname}}
\expandafter\def\csname BookIndex@T0557\endcsname{tapuqun@拓扑群（topological group）}
\expandafter\let\csname BookOriginal@trTopologicalGroup\endcsname\trTopologicalGroup
\renewcommand{\trTopologicalGroup}{\BookTerm{T0557}{\csname BookOriginal@trTopologicalGroup\endcsname}}
\expandafter\def\csname BookIndex@T0558\endcsname{youjie@有界（bounded）}
\expandafter\let\csname BookOriginal@trBounded\endcsname\trBounded
\renewcommand{\trBounded}{\BookTerm{T0558}{\csname BookOriginal@trBounded\endcsname}}
\expandafter\def\csname BookIndex@T0559\endcsname{lianxu@连续（continuous）}
\expandafter\let\csname BookOriginal@trContinuous\endcsname\trContinuous
\renewcommand{\trContinuous}{\BookTerm{T0559}{\csname BookOriginal@trContinuous\endcsname}}
\expandafter\def\csname BookIndex@T0560\endcsname{kaijinru@开浸入（open immersion）}
\expandafter\let\csname BookOriginal@trOpenImmersion\endcsname\trOpenImmersion
\renewcommand{\trOpenImmersion}{\BookTerm{T0560}{\csname BookOriginal@trOpenImmersion\endcsname}}
\expandafter\def\csname BookIndex@T0561\endcsname{biqianru@闭嵌入（closed immersion /\allowbreak{} closed embedding）}
\expandafter\let\csname BookOriginal@trClosedImmersion\endcsname\trClosedImmersion
\renewcommand{\trClosedImmersion}{\BookTerm{T0561}{\csname BookOriginal@trClosedImmersion\endcsname}}
\expandafter\def\csname BookIndex@T0562\endcsname{jinru@浸入（immersion）}
\expandafter\let\csname BookOriginal@trImmersion\endcsname\trImmersion
\renewcommand{\trImmersion}{\BookTerm{T0562}{\csname BookOriginal@trImmersion\endcsname}}
\expandafter\def\csname BookIndex@T0563\endcsname{kaiqianru@开嵌入（open embedding）}
\expandafter\let\csname BookOriginal@trOpenEmbedding\endcsname\trOpenEmbedding
\renewcommand{\trOpenEmbedding}{\BookTerm{T0563}{\csname BookOriginal@trOpenEmbedding\endcsname}}
\expandafter\def\csname BookIndex@T0564\endcsname{biqianru@闭嵌入（closed immersion /\allowbreak{} closed embedding）}
\expandafter\let\csname BookOriginal@trClosedEmbedding\endcsname\trClosedEmbedding
\renewcommand{\trClosedEmbedding}{\BookTerm{T0564}{\csname BookOriginal@trClosedEmbedding\endcsname}}
\expandafter\def\csname BookIndex@T0565\endcsname{fengfu@丰富（ample）}
\expandafter\let\csname BookOriginal@trAmple\endcsname\trAmple
\renewcommand{\trAmple}{\BookTerm{T0565}{\csname BookOriginal@trAmple\endcsname}}
\expandafter\def\csname BookIndex@T0566\endcsname{jifengfu@极丰富（very ample）}
\expandafter\let\csname BookOriginal@trVeryAmple\endcsname\trVeryAmple
\renewcommand{\trVeryAmple}{\BookTerm{T0566}{\csname BookOriginal@trVeryAmple\endcsname}}
\expandafter\def\csname BookIndex@T0567\endcsname{chuzi@除子（divisor）}
\expandafter\let\csname BookOriginal@trDivisor\endcsname\trDivisor
\renewcommand{\trDivisor}{\BookTerm{T0567}{\csname BookOriginal@trDivisor\endcsname}}
\expandafter\def\csname BookIndex@T0568\endcsname{kadiyechuzi@卡蒂耶除子（cartier divisor）}
\expandafter\let\csname BookOriginal@trCartierDivisor\endcsname\trCartierDivisor
\renewcommand{\trCartierDivisor}{\BookTerm{T0568}{\csname BookOriginal@trCartierDivisor\endcsname}}
\expandafter\def\csname BookIndex@T0569\endcsname{weiyichuzi@韦伊除子（weil divisor）}
\expandafter\let\csname BookOriginal@trWeilDivisor\endcsname\trWeilDivisor
\renewcommand{\trWeilDivisor}{\BookTerm{T0569}{\csname BookOriginal@trWeilDivisor\endcsname}}
\expandafter\def\csname BookIndex@T0570\endcsname{youxiaochuzi@有效除子（effective divisor）}
\expandafter\let\csname BookOriginal@trEffectiveDivisor\endcsname\trEffectiveDivisor
\renewcommand{\trEffectiveDivisor}{\BookTerm{T0570}{\csname BookOriginal@trEffectiveDivisor\endcsname}}
\expandafter\def\csname BookIndex@T0571\endcsname{keniceng@可逆层（invertible sheaf）}
\expandafter\let\csname BookOriginal@trInvertibleSheaf\endcsname\trInvertibleSheaf
\renewcommand{\trInvertibleSheaf}{\BookTerm{T0571}{\csname BookOriginal@trInvertibleSheaf\endcsname}}
\expandafter\def\csname BookIndex@T0572\endcsname{xiancong@线丛（line bundle）}
\expandafter\let\csname BookOriginal@trLineBundle\endcsname\trLineBundle
\renewcommand{\trLineBundle}{\BookTerm{T0572}{\csname BookOriginal@trLineBundle\endcsname}}
\expandafter\def\csname BookIndex@T0573\endcsname{xiangliangcong@向量丛（vector bundle）}
\expandafter\let\csname BookOriginal@trVectorBundle\endcsname\trVectorBundle
\renewcommand{\trVectorBundle}{\BookTerm{T0573}{\csname BookOriginal@trVectorBundle\endcsname}}
\expandafter\def\csname BookIndex@T0574\endcsname{qiecong@切丛（tangent bundle）}
\expandafter\let\csname BookOriginal@trTangentBundle\endcsname\trTangentBundle
\renewcommand{\trTangentBundle}{\BookTerm{T0574}{\csname BookOriginal@trTangentBundle\endcsname}}
\expandafter\def\csname BookIndex@T0575\endcsname{jubuziyouceng@局部自由层（locally free sheaf）}
\expandafter\let\csname BookOriginal@trLocallyFreeSheaf\endcsname\trLocallyFreeSheaf
\renewcommand{\trLocallyFreeSheaf}{\BookTerm{T0575}{\csname BookOriginal@trLocallyFreeSheaf\endcsname}}
\expandafter\def\csname BookIndex@T0576\endcsname{zhengtijiemian@整体截面（global sections）}
\expandafter\let\csname BookOriginal@trGlobalSections\endcsname\trGlobalSections
\renewcommand{\trGlobalSections}{\BookTerm{T0576}{\csname BookOriginal@trGlobalSections\endcsname}}
\expandafter\def\csname BookIndex@T0577\endcsname{shangtongdiao@上同调（cohomology）}
\expandafter\let\csname BookOriginal@trCohomology\endcsname\trCohomology
\renewcommand{\trCohomology}{\BookTerm{T0577}{\csname BookOriginal@trCohomology\endcsname}}
\expandafter\def\csname BookIndex@T0578\endcsname{qieheshangtongdiao@切赫上同调（cech cohomology）}
\expandafter\let\csname BookOriginal@trCechCohomology\endcsname\trCechCohomology
\renewcommand{\trCechCohomology}{\BookTerm{T0578}{\csname BookOriginal@trCechCohomology\endcsname}}
\expandafter\def\csname BookIndex@T0579\endcsname{daochuhanzi@导出函子（derived functor）}
\expandafter\let\csname BookOriginal@trDerivedFunctor\endcsname\trDerivedFunctor
\renewcommand{\trDerivedFunctor}{\BookTerm{T0579}{\csname BookOriginal@trDerivedFunctor\endcsname}}
\expandafter\def\csname BookIndex@T0580\endcsname{pingtan@平坦（flat）}
\expandafter\let\csname BookOriginal@trFlat\endcsname\trFlat
\renewcommand{\trFlat}{\BookTerm{T0580}{\csname BookOriginal@trFlat\endcsname}}
\expandafter\def\csname BookIndex@T0581\endcsname{wunao@无挠（torsion-free）}
\expandafter\let\csname BookOriginal@trTorsionFree\endcsname\trTorsionFree
\renewcommand{\trTorsionFree}{\BookTerm{T0581}{\csname BookOriginal@trTorsionFree\endcsname}}
\expandafter\def\csname BookIndex@T0582\endcsname{zhengguihua@正规化（normalization）}
\expandafter\let\csname BookOriginal@trNormalization\endcsname\trNormalization
\renewcommand{\trNormalization}{\BookTerm{T0582}{\csname BookOriginal@trNormalization\endcsname}}
\expandafter\def\csname BookIndex@T0583\endcsname{zhangkai@胀开（blow-up）}
\expandafter\let\csname BookOriginal@trBlowup\endcsname\trBlowup
\renewcommand{\trBlowup}{\BookTerm{T0583}{\csname BookOriginal@trBlowup\endcsname}}
\expandafter\def\csname BookIndex@T0584\endcsname{mokongjian@模空间（moduli space）}
\expandafter\let\csname BookOriginal@trModuliSpace\endcsname\trModuliSpace
\renewcommand{\trModuliSpace}{\BookTerm{T0584}{\csname BookOriginal@trModuliSpace\endcsname}}
\expandafter\def\csname BookIndex@T0585\endcsname{gelasimanliuxing@格拉斯曼流形（Grassmannian）}
\expandafter\let\csname BookOriginal@trGrassmannian\endcsname\trGrassmannian
\renewcommand{\trGrassmannian}{\BookTerm{T0585}{\csname BookOriginal@trGrassmannian\endcsname}}
\expandafter\def\csname BookIndex@T0586\endcsname{daishubi@代数闭（algebraically closed）}
\expandafter\let\csname BookOriginal@trAlgClosed\endcsname\trAlgClosed
\renewcommand{\trAlgClosed}{\BookTerm{T0586}{\csname BookOriginal@trAlgClosed\endcsname}}
\expandafter\def\csname BookIndex@T0587\endcsname{daishubiyu@代数闭域（algebraically closed field）}
\expandafter\let\csname BookOriginal@trAlgClosedField\endcsname\trAlgClosedField
\renewcommand{\trAlgClosedField}{\BookTerm{T0587}{\csname BookOriginal@trAlgClosedField\endcsname}}
\expandafter\def\csname BookIndex@T0588\endcsname{cu@簇（variety）}
\expandafter\let\csname BookOriginal@trVariety\endcsname\trVariety
\renewcommand{\trVariety}{\BookTerm{T0588}{\csname BookOriginal@trVariety\endcsname}}
\expandafter\def\csname BookIndex@T0589\endcsname{quxian@曲线（curve）}
\expandafter\let\csname BookOriginal@trCurve\endcsname\trCurve
\renewcommand{\trCurve}{\BookTerm{T0589}{\csname BookOriginal@trCurve\endcsname}}
\expandafter\def\csname BookIndex@T0590\endcsname{qumian@曲面（surface）}
\expandafter\let\csname BookOriginal@trSurface\endcsname\trSurface
\renewcommand{\trSurface}{\BookTerm{T0590}{\csname BookOriginal@trSurface\endcsname}}
\expandafter\def\csname BookIndex@T0591\endcsname{pikaqun@皮卡群（Picard group）}
\expandafter\let\csname BookOriginal@trPicardGroup\endcsname\trPicardGroup
\renewcommand{\trPicardGroup}{\BookTerm{T0591}{\csname BookOriginal@trPicardGroup\endcsname}}
\expandafter\def\csname BookIndex@T0592\endcsname{jubuhuan@局部环（local ring）}
\expandafter\let\csname BookOriginal@trLocalRing\endcsname\trLocalRing
\renewcommand{\trLocalRing}{\BookTerm{T0592}{\csname BookOriginal@trLocalRing\endcsname}}
\expandafter\def\csname BookIndex@T0593\endcsname{jubuhua@局部化（localization）}
\expandafter\let\csname BookOriginal@trLocalization\endcsname\trLocalization
\renewcommand{\trLocalization}{\BookTerm{T0593}{\csname BookOriginal@trLocalization\endcsname}}
\expandafter\def\csname BookIndex@T0594\endcsname{sulixiang@素理想（prime ideal）}
\expandafter\let\csname BookOriginal@trPrimeIdeal\endcsname\trPrimeIdeal
\renewcommand{\trPrimeIdeal}{\BookTerm{T0594}{\csname BookOriginal@trPrimeIdeal\endcsname}}
\expandafter\def\csname BookIndex@T0595\endcsname{jidalixiang@极大理想（maximal ideal）}
\expandafter\let\csname BookOriginal@trMaximalIdeal\endcsname\trMaximalIdeal
\renewcommand{\trMaximalIdeal}{\BookTerm{T0595}{\csname BookOriginal@trMaximalIdeal\endcsname}}
\expandafter\def\csname BookIndex@T0596\endcsname{lixiang@理想（ideal）}
\expandafter\let\csname BookOriginal@trIdeal\endcsname\trIdeal
\renewcommand{\trIdeal}{\BookTerm{T0596}{\csname BookOriginal@trIdeal\endcsname}}
\expandafter\def\csname BookIndex@T0597\endcsname{yuqiekongjian@余切空间（cotangent space）}
\expandafter\let\csname BookOriginal@trCotangentSpace\endcsname\trCotangentSpace
\renewcommand{\trCotangentSpace}{\BookTerm{T0597}{\csname BookOriginal@trCotangentSpace\endcsname}}
\expandafter\def\csname BookIndex@T0598\endcsname{shengyuyu@剩余域（residue field）}
\expandafter\let\csname BookOriginal@trResidueField\endcsname\trResidueField
\renewcommand{\trResidueField}{\BookTerm{T0598}{\csname BookOriginal@trResidueField\endcsname}}
\expandafter\def\csname BookIndex@T0599\endcsname{hanshuyu@函数域（function field）}
\expandafter\let\csname BookOriginal@trFunctionField\endcsname\trFunctionField
\renewcommand{\trFunctionField}{\BookTerm{T0599}{\csname BookOriginal@trFunctionField\endcsname}}
\expandafter\def\csname BookIndex@T0600\endcsname{lisanfuzhihuan@离散赋值环（discrete valuation ring）}
\expandafter\let\csname BookOriginal@trDVR\endcsname\trDVR
\renewcommand{\trDVR}{\BookTerm{T0600}{\csname BookOriginal@trDVR\endcsname}}
\expandafter\def\csname BookIndex@T0601\endcsname{chaoyuecishu@超越次数（transcendence degree）}
\expandafter\let\csname BookOriginal@trTrDeg\endcsname\trTrDeg
\renewcommand{\trTrDeg}{\BookTerm{T0601}{\csname BookOriginal@trTrDeg\endcsname}}
\expandafter\def\csname BookIndex@T0602\endcsname{keluerweishu@克鲁尔维数（krull dimension）}
\expandafter\let\csname BookOriginal@trKrullDimension\endcsname\trKrullDimension
\renewcommand{\trKrullDimension}{\BookTerm{T0602}{\csname BookOriginal@trKrullDimension\endcsname}}
\expandafter\def\csname BookIndex@T0603\endcsname{shendu@深度（depth）}
\expandafter\let\csname BookOriginal@trDepth\endcsname\trDepth
\renewcommand{\trDepth}{\BookTerm{T0603}{\csname BookOriginal@trDepth\endcsname}}
\expandafter\def\csname BookIndex@T0604\endcsname{keenmaikaoli@科恩--麦考利（cm）}
\expandafter\let\csname BookOriginal@trCM\endcsname\trCM
\renewcommand{\trCM}{\BookTerm{T0604}{\csname BookOriginal@trCM\endcsname}}
\expandafter\def\csname BookIndex@T0605\endcsname{weifen@微分（differential）}
\expandafter\let\csname BookOriginal@trDifferential\endcsname\trDifferential
\renewcommand{\trDifferential}{\BookTerm{T0605}{\csname BookOriginal@trDifferential\endcsname}}
\expandafter\def\csname BookIndex@T0606\endcsname{kaileiweifen@凯勒微分（Kähler differentials）}
\expandafter\let\csname BookOriginal@trKahlerDifferentials\endcsname\trKahlerDifferentials
\renewcommand{\trKahlerDifferentials}{\BookTerm{T0606}{\csname BookOriginal@trKahlerDifferentials\endcsname}}
\expandafter\def\csname BookIndex@T0607\endcsname{yuqiecong@余切丛（cotangent bundle）}
\expandafter\let\csname BookOriginal@trCotangentBundle\endcsname\trCotangentBundle
\renewcommand{\trCotangentBundle}{\BookTerm{T0607}{\csname BookOriginal@trCotangentBundle\endcsname}}
\expandafter\def\csname BookIndex@T0608\endcsname{qiekongjian@切空间（tangent space）}
\expandafter\let\csname BookOriginal@trTangentSpace\endcsname\trTangentSpace
\renewcommand{\trTangentSpace}{\BookTerm{T0608}{\csname BookOriginal@trTangentSpace\endcsname}}
\expandafter\def\csname BookIndex@T0609\endcsname{xiangjiaolilun@相交理论（intersection theory）}
\expandafter\let\csname BookOriginal@trIntersectionTheory\endcsname\trIntersectionTheory
\renewcommand{\trIntersectionTheory}{\BookTerm{T0609}{\csname BookOriginal@trIntersectionTheory\endcsname}}
\expandafter\def\csname BookIndex@T0610\endcsname{fencihuan@分次环（graded ring）}
\expandafter\let\csname BookOriginal@trGradedRing\endcsname\trGradedRing
\renewcommand{\trGradedRing}{\BookTerm{T0610}{\csname BookOriginal@trGradedRing\endcsname}}
\expandafter\def\csname BookIndex@T0611\endcsname{fencimo@分次模（graded module）}
\expandafter\let\csname BookOriginal@trGradedModule\endcsname\trGradedModule
\renewcommand{\trGradedModule}{\BookTerm{T0611}{\csname BookOriginal@trGradedModule\endcsname}}
\expandafter\def\csname BookIndex@T0612\endcsname{qici@齐次（homogeneous）}
\expandafter\let\csname BookOriginal@trHomogeneous\endcsname\trHomogeneous
\renewcommand{\trHomogeneous}{\BookTerm{T0612}{\csname BookOriginal@trHomogeneous\endcsname}}
\expandafter\def\csname BookIndex@T0613\endcsname{qungaixing@群概形（group scheme）}
\expandafter\let\csname BookOriginal@trGroupScheme\endcsname\trGroupScheme
\renewcommand{\trGroupScheme}{\BookTerm{T0613}{\csname BookOriginal@trGroupScheme\endcsname}}
\expandafter\def\csname BookIndex@T0614\endcsname{kuige@亏格（genus）}
\expandafter\let\csname BookOriginal@trGenus\endcsname\trGenus
\renewcommand{\trGenus}{\BookTerm{T0614}{\csname BookOriginal@trGenus\endcsname}}
\expandafter\def\csname BookIndex@T0615\endcsname{limanluohe@黎曼--罗赫（Riemann roch）}
\expandafter\let\csname BookOriginal@trRiemannRoch\endcsname\trRiemannRoch
\renewcommand{\trRiemannRoch}{\BookTerm{T0615}{\csname BookOriginal@trRiemannRoch\endcsname}}
\expandafter\def\csname BookIndex@T0616\endcsname{saierduiou@塞尔对偶（Serre duality）}
\expandafter\let\csname BookOriginal@trSerreDuality\endcsname\trSerreDuality
\renewcommand{\trSerreDuality}{\BookTerm{T0616}{\csname BookOriginal@trSerreDuality\endcsname}}
\expandafter\def\csname BookIndex@T0617\endcsname{duiou@对偶（duality /\allowbreak{} dual）}
\expandafter\let\csname BookOriginal@trDuality\endcsname\trDuality
\renewcommand{\trDuality}{\BookTerm{T0617}{\csname BookOriginal@trDuality\endcsname}}
\expandafter\def\csname BookIndex@T0618\endcsname{jubuxiaofanchou@局部小范畴（locally small category）}
\expandafter\let\csname BookOriginal@trLocallySmallCategory\endcsname\trLocallySmallCategory
\renewcommand{\trLocallySmallCategory}{\BookTerm{T0618}{\csname BookOriginal@trLocallySmallCategory\endcsname}}
\expandafter\def\csname BookIndex@T0619\endcsname{dengzi@等子（equalizer）}
\expandafter\let\csname BookOriginal@trEqualizer\endcsname\trEqualizer
\renewcommand{\trEqualizer}{\BookTerm{T0619}{\csname BookOriginal@trEqualizer\endcsname}}
\expandafter\def\csname BookIndex@T0620\endcsname{lei@类（class）}
\expandafter\let\csname BookOriginal@trClass\endcsname\trClass
\renewcommand{\trClass}{\BookTerm{T0620}{\csname BookOriginal@trClass\endcsname}}
\expandafter\def\csname BookIndex@T0621\endcsname{hanzixing@函子性（functoriality）}
\expandafter\let\csname BookOriginal@trFunctoriality\endcsname\trFunctoriality
\renewcommand{\trFunctoriality}{\BookTerm{T0621}{\csname BookOriginal@trFunctoriality\endcsname}}
\expandafter\def\csname BookIndex@T0622\endcsname{daolu@道路（path）}
\expandafter\let\csname BookOriginal@trPath\endcsname\trPath
\renewcommand{\trPath}{\BookTerm{T0622}{\csname BookOriginal@trPath\endcsname}}
\expandafter\def\csname BookIndex@T0623\endcsname{tonglun@同伦（homotopy）}
\expandafter\let\csname BookOriginal@trHomotopy\endcsname\trHomotopy
\renewcommand{\trHomotopy}{\BookTerm{T0623}{\csname BookOriginal@trHomotopy\endcsname}}
\expandafter\def\csname BookIndex@T0624\endcsname{liantong@连通（connected）}
\expandafter\let\csname BookOriginal@trConnected\endcsname\trConnected
\renewcommand{\trConnected}{\BookTerm{T0624}{\csname BookOriginal@trConnected\endcsname}}
\expandafter\def\csname BookIndex@T0625\endcsname{xianxingbianhuan@线性变换（linear transformation）}
\expandafter\let\csname BookOriginal@trLinearTransformation\endcsname\trLinearTransformation
\renewcommand{\trLinearTransformation}{\BookTerm{T0625}{\csname BookOriginal@trLinearTransformation\endcsname}}
\expandafter\def\csname BookIndex@T0626\endcsname{keni@可逆（invertible）}
\expandafter\let\csname BookOriginal@trInvertible\endcsname\trInvertible
\renewcommand{\trInvertible}{\BookTerm{T0626}{\csname BookOriginal@trInvertible\endcsname}}
\expandafter\def\csname BookIndex@T0627\endcsname{jihelun@集合论（set-theoretic）}
\expandafter\let\csname BookOriginal@trSetTheoretic\endcsname\trSetTheoretic
\renewcommand{\trSetTheoretic}{\BookTerm{T0627}{\csname BookOriginal@trSetTheoretic\endcsname}}
\expandafter\def\csname BookIndex@T0628\endcsname{beilun@悖论（paradox）}
\expandafter\let\csname BookOriginal@trParadox\endcsname\trParadox
\renewcommand{\trParadox}{\BookTerm{T0628}{\csname BookOriginal@trParadox\endcsname}}
\expandafter\def\csname BookIndex@T0629\endcsname{daishujihe@代数几何（algebraic geometry）}
\expandafter\let\csname BookOriginal@trAlgebraicGeometry\endcsname\trAlgebraicGeometry
\renewcommand{\trAlgebraicGeometry}{\BookTerm{T0629}{\csname BookOriginal@trAlgebraicGeometry\endcsname}}
\expandafter\def\csname BookIndex@T0630\endcsname{gaixinglilun@概形理论（scheme theory）}
\expandafter\let\csname BookOriginal@trSchemeTheory\endcsname\trSchemeTheory
\renewcommand{\trSchemeTheory}{\BookTerm{T0630}{\csname BookOriginal@trSchemeTheory\endcsname}}
\expandafter\def\csname BookIndex@T0631\endcsname{zhengdian@正典（canon）}
\expandafter\let\csname BookOriginal@trCanon\endcsname\trCanon
\renewcommand{\trCanon}{\BookTerm{T0631}{\csname BookOriginal@trCanon\endcsname}}
\expandafter\def\csname BookIndex@T0632\endcsname{fangshe@仿射（affine）}
\expandafter\let\csname BookOriginal@trAffine\endcsname\trAffine
\renewcommand{\trAffine}{\BookTerm{T0632}{\csname BookOriginal@trAffine\endcsname}}
\expandafter\def\csname BookIndex@T0633\endcsname{fu@复（complex）}
\expandafter\let\csname BookOriginal@trComplex\endcsname\trComplex
\renewcommand{\trComplex}{\BookTerm{T0633}{\csname BookOriginal@trComplex\endcsname}}
\expandafter\def\csname BookIndex@T0634\endcsname{shi@实（real）}
\expandafter\let\csname BookOriginal@trReal\endcsname\trReal
\renewcommand{\trReal}{\BookTerm{T0634}{\csname BookOriginal@trReal\endcsname}}
\expandafter\def\csname BookIndex@T0635\endcsname{liuxing@流形（manifold）}
\expandafter\let\csname BookOriginal@trManifold\endcsname\trManifold
\renewcommand{\trManifold}{\BookTerm{T0635}{\csname BookOriginal@trManifold\endcsname}}
\expandafter\def\csname BookIndex@T0636\endcsname{guanghuaxing@光滑性（smoothness）}
\expandafter\let\csname BookOriginal@trSmoothness\endcsname\trSmoothness
\renewcommand{\trSmoothness}{\BookTerm{T0636}{\csname BookOriginal@trSmoothness\endcsname}}
\expandafter\def\csname BookIndex@T0637\endcsname{weifenjihe@微分几何（differential geometry）}
\expandafter\let\csname BookOriginal@trDifferentialGeometry\endcsname\trDifferentialGeometry
\renewcommand{\trDifferentialGeometry}{\BookTerm{T0637}{\csname BookOriginal@trDifferentialGeometry\endcsname}}
\expandafter\def\csname BookIndex@T0638\endcsname{fujihe@复几何（complex geometry）}
\expandafter\let\csname BookOriginal@trComplexGeometry\endcsname\trComplexGeometry
\renewcommand{\trComplexGeometry}{\BookTerm{T0638}{\csname BookOriginal@trComplexGeometry\endcsname}}
\expandafter\def\csname BookIndex@T0639\endcsname{fenxi@分析（analysis）}
\expandafter\let\csname BookOriginal@trAnalysis\endcsname\trAnalysis
\renewcommand{\trAnalysis}{\BookTerm{T0639}{\csname BookOriginal@trAnalysis\endcsname}}
\expandafter\def\csname BookIndex@T0640\endcsname{shulun@数论（number theory）}
\expandafter\let\csname BookOriginal@trNumberTheory\endcsname\trNumberTheory
\renewcommand{\trNumberTheory}{\BookTerm{T0640}{\csname BookOriginal@trNumberTheory\endcsname}}
\expandafter\def\csname BookIndex@T0641\endcsname{daishutapu@代数拓扑（algebraic topology）}
\expandafter\let\csname BookOriginal@trAlgebraicTopology\endcsname\trAlgebraicTopology
\renewcommand{\trAlgebraicTopology}{\BookTerm{T0641}{\csname BookOriginal@trAlgebraicTopology\endcsname}}
\expandafter\def\csname BookIndex@T0642\endcsname{tongdiaodaishu@同调代数（homological algebra）}
\expandafter\let\csname BookOriginal@trHomologicalAlgebra\endcsname\trHomologicalAlgebra
\renewcommand{\trHomologicalAlgebra}{\BookTerm{T0642}{\csname BookOriginal@trHomologicalAlgebra\endcsname}}
\expandafter\def\csname BookIndex@T0643\endcsname{zuobiaopian@坐标片（coordinate patch）}
\expandafter\let\csname BookOriginal@trCoordinatePatch\endcsname\trCoordinatePatch
\renewcommand{\trCoordinatePatch}{\BookTerm{T0643}{\csname BookOriginal@trCoordinatePatch\endcsname}}
\expandafter\def\csname BookIndex@T0644\endcsname{fuhuankongjian@赋环空间（ringed space）}
\expandafter\let\csname BookOriginal@trRingedSpace\endcsname\trRingedSpace
\renewcommand{\trRingedSpace}{\BookTerm{T0644}{\csname BookOriginal@trRingedSpace\endcsname}}
\expandafter\def\csname BookIndex@T0645\endcsname{fujiexikongjian@复解析空间（complex analytic space）}
\expandafter\let\csname BookOriginal@trComplexAnalyticSpace\endcsname\trComplexAnalyticSpace
\renewcommand{\trComplexAnalyticSpace}{\BookTerm{T0645}{\csname BookOriginal@trComplexAnalyticSpace\endcsname}}
\expandafter\def\csname BookIndex@T0646\endcsname{daishukongjian@代数空间（algebraic space）}
\expandafter\let\csname BookOriginal@trAlgebraicSpace\endcsname\trAlgebraicSpace
\renewcommand{\trAlgebraicSpace}{\BookTerm{T0646}{\csname BookOriginal@trAlgebraicSpace\endcsname}}
\expandafter\def\csname BookIndex@T0647\endcsname{guixing@轨形（orbifold）}
\expandafter\let\csname BookOriginal@trOrbifold\endcsname\trOrbifold
\renewcommand{\trOrbifold}{\BookTerm{T0647}{\csname BookOriginal@trOrbifold\endcsname}}
\expandafter\def\csname BookIndex@T0648\endcsname{die@叠（stack）}
\expandafter\let\csname BookOriginal@trStack\endcsname\trStack
\renewcommand{\trStack}{\BookTerm{T0648}{\csname BookOriginal@trStack\endcsname}}
\expandafter\def\csname BookIndex@T0649\endcsname{gangxingjiexikongjian@刚性解析空间（rigid analytic space）}
\expandafter\let\csname BookOriginal@trRigidAnalyticSpace\endcsname\trRigidAnalyticSpace
\renewcommand{\trRigidAnalyticSpace}{\BookTerm{T0649}{\csname BookOriginal@trRigidAnalyticSpace\endcsname}}
\expandafter\def\csname BookIndex@T0650\endcsname{xingshigaixing@形式概形（formal scheme）}
\expandafter\let\csname BookOriginal@trFormalScheme\endcsname\trFormalScheme
\renewcommand{\trFormalScheme}{\BookTerm{T0650}{\csname BookOriginal@trFormalScheme\endcsname}}
\expandafter\def\csname BookIndex@T0651\endcsname{mitianyinli@米田引理（Yoneda lemma）}
\expandafter\let\csname BookOriginal@trYonedaLemma\endcsname\trYonedaLemma
\renewcommand{\trYonedaLemma}{\BookTerm{T0651}{\csname BookOriginal@trYonedaLemma\endcsname}}
\expandafter\def\csname BookIndex@T0652\endcsname{ji@积（product）}
\expandafter\let\csname BookOriginal@trProduct\endcsname\trProduct
\renewcommand{\trProduct}{\BookTerm{T0652}{\csname BookOriginal@trProduct\endcsname}}
\expandafter\def\csname BookIndex@T0653\endcsname{jidi@基底（base）}
\expandafter\let\csname BookOriginal@trBase\endcsname\trBase
\renewcommand{\trBase}{\BookTerm{T0653}{\csname BookOriginal@trBase\endcsname}}
\expandafter\def\csname BookIndex@T0654\endcsname{culun@簇论（theory of varieties）}
\expandafter\let\csname BookOriginal@trVarietyTheory\endcsname\trVarietyTheory
\renewcommand{\trVarietyTheory}{\BookTerm{T0654}{\csname BookOriginal@trVarietyTheory\endcsname}}
\expandafter\def\csname BookIndex@T0655\endcsname{limanqumian@黎曼曲面（Riemann surface）}
\expandafter\let\csname BookOriginal@trRiemannSurface\endcsname\trRiemannSurface
\renewcommand{\trRiemannSurface}{\BookTerm{T0655}{\csname BookOriginal@trRiemannSurface\endcsname}}
\expandafter\def\csname BookIndex@T0656\endcsname{zuhedaishujihe@组合代数几何（combinatorial algebraic geometry）}
\expandafter\let\csname BookOriginal@trCombinatorialAlgGeom\endcsname\trCombinatorialAlgGeom
\renewcommand{\trCombinatorialAlgGeom}{\BookTerm{T0656}{\csname BookOriginal@trCombinatorialAlgGeom\endcsname}}
\expandafter\def\csname BookIndex@T0657\endcsname{jisuandaishujihe@计算代数几何（computational algebraic geometry）}
\expandafter\let\csname BookOriginal@trComputationalAlgGeom\endcsname\trComputationalAlgGeom
\renewcommand{\trComputationalAlgGeom}{\BookTerm{T0657}{\csname BookOriginal@trComputationalAlgGeom\endcsname}}
\expandafter\def\csname BookIndex@T0658\endcsname{quanchun@全纯（holomorphic）}
\expandafter\let\csname BookOriginal@trHolomorphic\endcsname\trHolomorphic
\renewcommand{\trHolomorphic}{\BookTerm{T0658}{\csname BookOriginal@trHolomorphic\endcsname}}
\expandafter\def\csname BookIndex@T0659\endcsname{jiexi@解析（analytic）}
\expandafter\let\csname BookOriginal@trAnalytic\endcsname\trAnalytic
\renewcommand{\trAnalytic}{\BookTerm{T0659}{\csname BookOriginal@trAnalytic\endcsname}}
\expandafter\def\csname BookIndex@T0660\endcsname{kewei@可微（differentiable）}
\expandafter\let\csname BookOriginal@trDifferentiable\endcsname\trDifferentiable
\renewcommand{\trDifferentiable}{\BookTerm{T0660}{\csname BookOriginal@trDifferentiable\endcsname}}
\expandafter\def\csname BookIndex@T0661\endcsname{yuzhou@宇宙（universe）}
\expandafter\let\csname BookOriginal@trUniverse\endcsname\trUniverse
\renewcommand{\trUniverse}{\BookTerm{T0661}{\csname BookOriginal@trUniverse\endcsname}}
\expandafter\def\csname BookIndex@T0662\endcsname{nijin@拟紧（quasicompact）}
\expandafter\let\csname BookOriginal@trQuasicompact\endcsname\trQuasicompact
\renewcommand{\trQuasicompact}{\BookTerm{T0662}{\csname BookOriginal@trQuasicompact\endcsname}}
\expandafter\def\csname BookIndex@T0663\endcsname{nifenli@拟分离（quasiseparated）}
\expandafter\let\csname BookOriginal@trQuasiSeparated\endcsname\trQuasiSeparated
\renewcommand{\trQuasiSeparated}{\BookTerm{T0663}{\csname BookOriginal@trQuasiSeparated\endcsname}}
\expandafter\def\csname BookIndex@T0664\endcsname{nifenlixing@拟分离性（quasiseparatedness）}
\expandafter\let\csname BookOriginal@trQuasiseparatedness\endcsname\trQuasiseparatedness
\renewcommand{\trQuasiseparatedness}{\BookTerm{T0664}{\csname BookOriginal@trQuasiseparatedness\endcsname}}
\expandafter\def\csname BookIndex@T0665\endcsname{fuzhipanbiefa@赋值判别法（valuative criteria）}
\expandafter\let\csname BookOriginal@trValuativeCriteria\endcsname\trValuativeCriteria
\renewcommand{\trValuativeCriteria}{\BookTerm{T0665}{\csname BookOriginal@trValuativeCriteria\endcsname}}
\expandafter\def\csname BookIndex@T0666\endcsname{jiaohuan@交换（commutative）}
\expandafter\let\csname BookOriginal@trCommutative\endcsname\trCommutative
\renewcommand{\trCommutative}{\BookTerm{T0666}{\csname BookOriginal@trCommutative\endcsname}}
\expandafter\def\csname BookIndex@T0667\endcsname{danweiyuan@单位元（unit）}
\expandafter\let\csname BookOriginal@trUnit\endcsname\trUnit
\renewcommand{\trUnit}{\BookTerm{T0667}{\csname BookOriginal@trUnit\endcsname}}
\expandafter\def\csname BookIndex@T0668\endcsname{linghuan@零环（zero ring）}
\expandafter\let\csname BookOriginal@trZeroRing\endcsname\trZeroRing
\renewcommand{\trZeroRing}{\BookTerm{T0668}{\csname BookOriginal@trZeroRing\endcsname}}
\expandafter\def\csname BookIndex@T0669\endcsname{zhenghuan@整环（integral domain）}
\expandafter\let\csname BookOriginal@trIntegralDomain\endcsname\trIntegralDomain
\renewcommand{\trIntegralDomain}{\BookTerm{T0669}{\csname BookOriginal@trIntegralDomain\endcsname}}
\expandafter\def\csname BookIndex@T0670\endcsname{zhongduixiang@终对象（final object）}
\expandafter\let\csname BookOriginal@trFinalObject\endcsname\trFinalObject
\renewcommand{\trFinalObject}{\BookTerm{T0670}{\csname BookOriginal@trFinalObject\endcsname}}
\expandafter\def\csname BookIndex@T0671\endcsname{chushiduixiang@初始对象（initial object）}
\expandafter\let\csname BookOriginal@trInitialObject\endcsname\trInitialObject
\renewcommand{\trInitialObject}{\BookTerm{T0671}{\csname BookOriginal@trInitialObject\endcsname}}
\expandafter\def\csname BookIndex@T0672\endcsname{xuanzegongli@选择公理（Axiom of Choice）}
\expandafter\let\csname BookOriginal@trAxiomOfChoice\endcsname\trAxiomOfChoice
\renewcommand{\trAxiomOfChoice}{\BookTerm{T0672}{\csname BookOriginal@trAxiomOfChoice\endcsname}}
\expandafter\def\csname BookIndex@T0673\endcsname{zuoenyinli@佐恩引理（zorns lemma）}
\expandafter\let\csname BookOriginal@trZornsLemma\endcsname\trZornsLemma
\renewcommand{\trZornsLemma}{\BookTerm{T0673}{\csname BookOriginal@trZornsLemma\endcsname}}
\expandafter\def\csname BookIndex@T0674\endcsname{zhenlixiang@真理想（proper ideal）}
\expandafter\let\csname BookOriginal@trProperIdeal\endcsname\trProperIdeal
\renewcommand{\trProperIdeal}{\BookTerm{T0674}{\csname BookOriginal@trProperIdeal\endcsname}}
\expandafter\def\csname BookIndex@T0675\endcsname{neishe@内射（injective）}
\expandafter\let\csname BookOriginal@trInjective\endcsname\trInjective
\renewcommand{\trInjective}{\BookTerm{T0675}{\csname BookOriginal@trInjective\endcsname}}
\expandafter\def\csname BookIndex@T0676\endcsname{toushe@投射（projection /\allowbreak{} projective）}
\expandafter\let\csname BookOriginal@trProjectiveObj\endcsname\trProjectiveObj
\renewcommand{\trProjectiveObj}{\BookTerm{T0676}{\csname BookOriginal@trProjectiveObj\endcsname}}
\expandafter\def\csname BookIndex@T0677\endcsname{sheying@射影（projective）}
\expandafter\let\csname BookOriginal@trProjective\endcsname\trProjective
\renewcommand{\trProjective}{\BookTerm{T0677}{\csname BookOriginal@trProjective\endcsname}}
\expandafter\def\csname BookIndex@T0678\endcsname{sheyingkongjian@射影空间（projective space）}
\expandafter\let\csname BookOriginal@trProjectiveSpace\endcsname\trProjectiveSpace
\renewcommand{\trProjectiveSpace}{\BookTerm{T0678}{\csname BookOriginal@trProjectiveSpace\endcsname}}
\expandafter\def\csname BookIndex@T0679\endcsname{zhulixiangzhenghuan@主理想整环（principal ideal domain）}
\expandafter\let\csname BookOriginal@trPID\endcsname\trPID
\renewcommand{\trPID}{\BookTerm{T0679}{\csname BookOriginal@trPID\endcsname}}
\expandafter\def\csname BookIndex@T0680\endcsname{weiyifenjiezhenghuan@唯一分解整环（unique factorization domain）}
\expandafter\let\csname BookOriginal@trUFD\endcsname\trUFD
\renewcommand{\trUFD}{\BookTerm{T0680}{\csname BookOriginal@trUFD\endcsname}}
\expandafter\def\csname BookIndex@T0681\endcsname{zhangliangji@张量积（tensor product）}
\expandafter\let\csname BookOriginal@trTensorProduct\endcsname\trTensorProduct
\renewcommand{\trTensorProduct}{\BookTerm{T0681}{\csname BookOriginal@trTensorProduct\endcsname}}
\expandafter\def\csname BookIndex@T0682\endcsname{zhenghelie@正合列（exact sequence）}
\expandafter\let\csname BookOriginal@trExactSequence\endcsname\trExactSequence
\renewcommand{\trExactSequence}{\BookTerm{T0682}{\csname BookOriginal@trExactSequence\endcsname}}
\expandafter\def\csname BookIndex@T0683\endcsname{zhihe@直和（direct sum）}
\expandafter\let\csname BookOriginal@trDirectSum\endcsname\trDirectSum
\renewcommand{\trDirectSum}{\BookTerm{T0683}{\csname BookOriginal@trDirectSum\endcsname}}
\expandafter\def\csname BookIndex@T0684\endcsname{shuangshe@双射（bijection）}
\expandafter\let\csname BookOriginal@trBijection\endcsname\trBijection
\renewcommand{\trBijection}{\BookTerm{T0684}{\csname BookOriginal@trBijection\endcsname}}
\expandafter\def\csname BookIndex@T0685\endcsname{danshe@单射（injective）}
\expandafter\let\csname BookOriginal@trInjectiveMap\endcsname\trInjectiveMap
\renewcommand{\trInjectiveMap}{\BookTerm{T0685}{\csname BookOriginal@trInjectiveMap\endcsname}}
\expandafter\def\csname BookIndex@T0686\endcsname{manshe@满射（surjective）}
\expandafter\let\csname BookOriginal@trSurjectiveMap\endcsname\trSurjectiveMap
\renewcommand{\trSurjectiveMap}{\BookTerm{T0686}{\csname BookOriginal@trSurjectiveMap\endcsname}}
\expandafter\def\csname BookIndex@T0687\endcsname{lahui@拉回（pullback）}
\expandafter\let\csname BookOriginal@trPullback\endcsname\trPullback
\renewcommand{\trPullback}{\BookTerm{T0687}{\csname BookOriginal@trPullback\endcsname}}
\expandafter\def\csname BookIndex@T0688\endcsname{manzifanchou@满子范畴（full subcategory）}
\expandafter\let\csname BookOriginal@trFullSubcategory\endcsname\trFullSubcategory
\renewcommand{\trFullSubcategory}{\BookTerm{T0688}{\csname BookOriginal@trFullSubcategory\endcsname}}
\expandafter\def\csname BookIndex@T0689\endcsname{fanfanchou@反范畴（opposite category）}
\expandafter\let\csname BookOriginal@trOppositeCategory\endcsname\trOppositeCategory
\renewcommand{\trOppositeCategory}{\BookTerm{T0689}{\csname BookOriginal@trOppositeCategory\endcsname}}
\expandafter\def\csname BookIndex@T0690\endcsname{ziranbianhuan@自然变换（natural transformation）}
\expandafter\let\csname BookOriginal@trNaturalTransformation\endcsname\trNaturalTransformation
\renewcommand{\trNaturalTransformation}{\BookTerm{T0690}{\csname BookOriginal@trNaturalTransformation\endcsname}}
\expandafter\def\csname BookIndex@T0691\endcsname{zirantonggou@自然同构（natural isomorphism）}
\expandafter\let\csname BookOriginal@trNaturalIsomorphism\endcsname\trNaturalIsomorphism
\renewcommand{\trNaturalIsomorphism}{\BookTerm{T0691}{\csname BookOriginal@trNaturalIsomorphism\endcsname}}
\expandafter\def\csname BookIndex@T0692\endcsname{fanchoudedengjia@范畴的等价（equivalence of categories）}
\expandafter\let\csname BookOriginal@trEquivalenceOfCategories\endcsname\trEquivalenceOfCategories
\renewcommand{\trEquivalenceOfCategories}{\BookTerm{T0692}{\csname BookOriginal@trEquivalenceOfCategories\endcsname}}
\expandafter\def\csname BookIndex@T0693\endcsname{dianhanzi@点函子（functor of points）}
\expandafter\let\csname BookOriginal@trFunctorOfPoints\endcsname\trFunctorOfPoints
\renewcommand{\trFunctorOfPoints}{\BookTerm{T0693}{\csname BookOriginal@trFunctorOfPoints\endcsname}}
\expandafter\def\csname BookIndex@T0694\endcsname{dijihe@底集合（underlying set）}
\expandafter\let\csname BookOriginal@trUnderlyingSet\endcsname\trUnderlyingSet
\renewcommand{\trUnderlyingSet}{\BookTerm{T0694}{\csname BookOriginal@trUnderlyingSet\endcsname}}
\expandafter\def\csname BookIndex@T0695\endcsname{lianxuyingshe@连续映射（continuous map）}
\expandafter\let\csname BookOriginal@trContinuousMap\endcsname\trContinuousMap
\renewcommand{\trContinuousMap}{\BookTerm{T0695}{\csname BookOriginal@trContinuousMap\endcsname}}
\expandafter\def\csname BookIndex@T0696\endcsname{youxianwei@有限维（finite-dimensional）}
\expandafter\let\csname BookOriginal@trFiniteDimensional\endcsname\trFiniteDimensional
\renewcommand{\trFiniteDimensional}{\BookTerm{T0696}{\csname BookOriginal@trFiniteDimensional\endcsname}}
\expandafter\def\csname BookIndex@T0697\endcsname{duiou@对偶（duality /\allowbreak{} dual）}
\expandafter\let\csname BookOriginal@trDual\endcsname\trDual
\renewcommand{\trDual}{\BookTerm{T0697}{\csname BookOriginal@trDual\endcsname}}
\expandafter\def\csname BookIndex@T0698\endcsname{xiaofanchou@小范畴（small category）}
\expandafter\let\csname BookOriginal@trSmallCategory\endcsname\trSmallCategory
\renewcommand{\trSmallCategory}{\BookTerm{T0698}{\csname BookOriginal@trSmallCategory\endcsname}}
\expandafter\def\csname BookIndex@T0699\endcsname{jiaohuantu@交换图（commutative diagram）}
\expandafter\let\csname BookOriginal@trCommutativeDiagram\endcsname\trCommutativeDiagram
\renewcommand{\trCommutativeDiagram}{\BookTerm{T0699}{\csname BookOriginal@trCommutativeDiagram\endcsname}}
\expandafter\def\csname BookIndex@T0700\endcsname{lingduixiang@零对象（zero object）}
\expandafter\let\csname BookOriginal@trZeroObject\endcsname\trZeroObject
\renewcommand{\trZeroObject}{\BookTerm{T0700}{\csname BookOriginal@trZeroObject\endcsname}}
\expandafter\def\csname BookIndex@T0701\endcsname{lingyinzi@零因子（zerodivisor）}
\expandafter\let\csname BookOriginal@trZerodivisor\endcsname\trZerodivisor
\renewcommand{\trZerodivisor}{\BookTerm{T0701}{\csname BookOriginal@trZerodivisor\endcsname}}
\expandafter\def\csname BookIndex@T0702\endcsname{feilingyinzi@非零因子（non-zerodivisor）}
\expandafter\let\csname BookOriginal@trNonZerodivisor\endcsname\trNonZerodivisor
\renewcommand{\trNonZerodivisor}{\BookTerm{T0702}{\csname BookOriginal@trNonZerodivisor\endcsname}}
\expandafter\def\csname BookIndex@T0703\endcsname{youxianshengcheng@有限生成（finitely generated）}
\expandafter\let\csname BookOriginal@trFinitelyGenerated\endcsname\trFinitelyGenerated
\renewcommand{\trFinitelyGenerated}{\BookTerm{T0703}{\csname BookOriginal@trFinitelyGenerated\endcsname}}
\expandafter\def\csname BookIndex@T0704\endcsname{fenshiyu@分式域（fraction field）}
\expandafter\let\csname BookOriginal@trFractionField\endcsname\trFractionField
\renewcommand{\trFractionField}{\BookTerm{T0704}{\csname BookOriginal@trFractionField\endcsname}}
\expandafter\def\csname BookIndex@T0705\endcsname{chengxingziji@乘性子集（multiplicative subset）}
\expandafter\let\csname BookOriginal@trMultiplicativeSubset\endcsname\trMultiplicativeSubset
\renewcommand{\trMultiplicativeSubset}{\BookTerm{T0705}{\csname BookOriginal@trMultiplicativeSubset\endcsname}}
\expandafter\def\csname BookIndex@T0706\endcsname{ziyoumo@自由模（free module）}
\expandafter\let\csname BookOriginal@trFreeModule\endcsname\trFreeModule
\renewcommand{\trFreeModule}{\BookTerm{T0706}{\csname BookOriginal@trFreeModule\endcsname}}
\expandafter\def\csname BookIndex@T0707\endcsname{zimo@子模（submodule）}
\expandafter\let\csname BookOriginal@trSubmodule\endcsname\trSubmodule
\renewcommand{\trSubmodule}{\BookTerm{T0707}{\csname BookOriginal@trSubmodule\endcsname}}
\expandafter\def\csname BookIndex@T0708\endcsname{shuangxianxing@双线性（bilinear）}
\expandafter\let\csname BookOriginal@trBilinear\endcsname\trBilinear
\renewcommand{\trBilinear}{\BookTerm{T0708}{\csname BookOriginal@trBilinear\endcsname}}
\expandafter\def\csname BookIndex@T0709\endcsname{youzhenghe@右正合（right-exact）}
\expandafter\let\csname BookOriginal@trRightExact\endcsname\trRightExact
\renewcommand{\trRightExact}{\BookTerm{T0709}{\csname BookOriginal@trRightExact\endcsname}}
\expandafter\def\csname BookIndex@T0710\endcsname{zuozhenghe@左正合（left-exact）}
\expandafter\let\csname BookOriginal@trLeftExact\endcsname\trLeftExact
\renewcommand{\trLeftExact}{\BookTerm{T0710}{\csname BookOriginal@trLeftExact\endcsname}}
\expandafter\def\csname BookIndex@T0711\endcsname{huanyingshe@环映射（ring morphism）}
\expandafter\let\csname BookOriginal@trRingMorphism\endcsname\trRingMorphism
\renewcommand{\trRingMorphism}{\BookTerm{T0711}{\csname BookOriginal@trRingMorphism\endcsname}}
\expandafter\def\csname BookIndex@T0712\endcsname{daishu@代数（algebraic /\allowbreak{} algebra）}
\expandafter\let\csname BookOriginal@trAlgebra\endcsname\trAlgebra
\renewcommand{\trAlgebra}{\BookTerm{T0712}{\csname BookOriginal@trAlgebra\endcsname}}
\expandafter\def\csname BookIndex@T0713\endcsname{zihuan@子环（subring）}
\expandafter\let\csname BookOriginal@trSubring\endcsname\trSubring
\renewcommand{\trSubring}{\BookTerm{T0713}{\csname BookOriginal@trSubring\endcsname}}
\expandafter\def\csname BookIndex@T0714\endcsname{liangding@良定（well-defined）}
\expandafter\let\csname BookOriginal@trWellDefined\endcsname\trWellDefined
\renewcommand{\trWellDefined}{\BookTerm{T0714}{\csname BookOriginal@trWellDefined\endcsname}}
\expandafter\def\csname BookIndex@T0715\endcsname{dengjiaguanxi@等价关系（equivalence relation）}
\expandafter\let\csname BookOriginal@trEquivalenceRelation\endcsname\trEquivalenceRelation
\renewcommand{\trEquivalenceRelation}{\BookTerm{T0715}{\csname BookOriginal@trEquivalenceRelation\endcsname}}
\expandafter\def\csname BookIndex@T0716\endcsname{dengjialei@等价类（equivalence class）}
\expandafter\let\csname BookOriginal@trEquivalenceClass\endcsname\trEquivalenceClass
\renewcommand{\trEquivalenceClass}{\BookTerm{T0716}{\csname BookOriginal@trEquivalenceClass\endcsname}}
\expandafter\def\csname BookIndex@T0717\endcsname{eryuanyunsuan@二元运算（binary operation）}
\expandafter\let\csname BookOriginal@trBinaryOperation\endcsname\trBinaryOperation
\renewcommand{\trBinaryOperation}{\BookTerm{T0717}{\csname BookOriginal@trBinaryOperation\endcsname}}
\expandafter\def\csname BookIndex@T0718\endcsname{youxianbiaoxian@有限表现（finitely presented）}
\expandafter\let\csname BookOriginal@trFinitelyPresented\endcsname\trFinitelyPresented
\renewcommand{\trFinitelyPresented}{\BookTerm{T0718}{\csname BookOriginal@trFinitelyPresented\endcsname}}
\expandafter\def\csname BookIndex@T0719\endcsname{shang@商（quotient）}
\expandafter\let\csname BookOriginal@trQuotient\endcsname\trQuotient
\renewcommand{\trQuotient}{\BookTerm{T0719}{\csname BookOriginal@trQuotient\endcsname}}
\expandafter\def\csname BookIndex@T0720\endcsname{duijiaotaishe@对角态射（diagonal morphism）}
\expandafter\let\csname BookOriginal@trDiagonalMorphism\endcsname\trDiagonalMorphism
\renewcommand{\trDiagonalMorphism}{\BookTerm{T0720}{\csname BookOriginal@trDiagonalMorphism\endcsname}}
\expandafter\def\csname BookIndex@T0721\endcsname{dikaerfangtu@笛卡尔方图（fibered/\allowbreak{}pullback/\allowbreak{}Cartesian diagram）}
\expandafter\let\csname BookOriginal@trCartesianDiagram\endcsname\trCartesianDiagram
\renewcommand{\trCartesianDiagram}{\BookTerm{T0721}{\csname BookOriginal@trCartesianDiagram\endcsname}}
\expandafter\def\csname BookIndex@T0722\endcsname{yuji@余积（coproduct）}
\expandafter\let\csname BookOriginal@trCoproduct\endcsname\trCoproduct
\renewcommand{\trCoproduct}{\BookTerm{T0722}{\csname BookOriginal@trCoproduct\endcsname}}
\expandafter\def\csname BookIndex@T0723\endcsname{xianweiyuji@纤维余积（fibered coproduct）}
\expandafter\let\csname BookOriginal@trFiberedCoproduct\endcsname\trFiberedCoproduct
\renewcommand{\trFiberedCoproduct}{\BookTerm{T0723}{\csname BookOriginal@trFiberedCoproduct\endcsname}}
\expandafter\def\csname BookIndex@T0724\endcsname{bujiaobing@不交并（disjoint union）}
\expandafter\let\csname BookOriginal@trDisjointUnion\endcsname\trDisjointUnion
\renewcommand{\trDisjointUnion}{\BookTerm{T0724}{\csname BookOriginal@trDisjointUnion\endcsname}}
\expandafter\def\csname BookIndex@T0725\endcsname{dantaishe@单态射（monomorphism）}
\expandafter\let\csname BookOriginal@trMonomorphism\endcsname\trMonomorphism
\renewcommand{\trMonomorphism}{\BookTerm{T0725}{\csname BookOriginal@trMonomorphism\endcsname}}
\expandafter\def\csname BookIndex@T0726\endcsname{mantaishe@满态射（epimorphism）}
\expandafter\let\csname BookOriginal@trEpimorphism\endcsname\trEpimorphism
\renewcommand{\trEpimorphism}{\BookTerm{T0726}{\csname BookOriginal@trEpimorphism\endcsname}}
\expandafter\def\csname BookIndex@T0727\endcsname{kechuqun@可除群（divisible group）}
\expandafter\let\csname BookOriginal@trDivisibleGroup\endcsname\trDivisibleGroup
\renewcommand{\trDivisibleGroup}{\BookTerm{T0727}{\csname BookOriginal@trDivisibleGroup\endcsname}}
\expandafter\def\csname BookIndex@T0728\endcsname{kebiao@可表（representable）}
\expandafter\let\csname BookOriginal@trRepresentable\endcsname\trRepresentable
\renewcommand{\trRepresentable}{\BookTerm{T0728}{\csname BookOriginal@trRepresentable\endcsname}}
\expandafter\def\csname BookIndex@T0729\endcsname{hanzifanchou@函子范畴（functor category）}
\expandafter\let\csname BookOriginal@trFunctorCategory\endcsname\trFunctorCategory
\renewcommand{\trFunctorCategory}{\BookTerm{T0729}{\csname BookOriginal@trFunctorCategory\endcsname}}
\expandafter\def\csname BookIndex@T0730\endcsname{mitianqianru@米田嵌入（Yoneda embedding）}
\expandafter\let\csname BookOriginal@trYonedaEmbedding\endcsname\trYonedaEmbedding
\renewcommand{\trYonedaEmbedding}{\BookTerm{T0730}{\csname BookOriginal@trYonedaEmbedding\endcsname}}
\expandafter\def\csname BookIndex@T0731\endcsname{fanduixiang@泛对象（universal object）}
\expandafter\let\csname BookOriginal@trUniversalObject\endcsname\trUniversalObject
\renewcommand{\trUniversalObject}{\BookTerm{T0731}{\csname BookOriginal@trUniversalObject\endcsname}}
\expandafter\def\csname BookIndex@T0732\endcsname{tu@图（diagram）}
\expandafter\let\csname BookOriginal@trDiagram\endcsname\trDiagram
\renewcommand{\trDiagram}{\BookTerm{T0732}{\csname BookOriginal@trDiagram\endcsname}}
\expandafter\def\csname BookIndex@T0733\endcsname{zhibiaofanchou@指标范畴（index category）}
\expandafter\let\csname BookOriginal@trIndexCategory\endcsname\trIndexCategory
\renewcommand{\trIndexCategory}{\BookTerm{T0733}{\csname BookOriginal@trIndexCategory\endcsname}}
\expandafter\def\csname BookIndex@T0734\endcsname{zhibiaoji@指标集（index set）}
\expandafter\let\csname BookOriginal@trIndexSet\endcsname\trIndexSet
\renewcommand{\trIndexSet}{\BookTerm{T0734}{\csname BookOriginal@trIndexSet\endcsname}}
\expandafter\def\csname BookIndex@T0735\endcsname{nijixian@逆极限（inverse limit）}
\expandafter\let\csname BookOriginal@trInverseLimit\endcsname\trInverseLimit
\renewcommand{\trInverseLimit}{\BookTerm{T0735}{\csname BookOriginal@trInverseLimit\endcsname}}
\expandafter\def\csname BookIndex@T0736\endcsname{toushejixian@投射极限（projective limit）}
\expandafter\let\csname BookOriginal@trProjectiveLimit\endcsname\trProjectiveLimit
\renewcommand{\trProjectiveLimit}{\BookTerm{T0736}{\csname BookOriginal@trProjectiveLimit\endcsname}}
\expandafter\def\csname BookIndex@T0737\endcsname{zhengxiangjixian@正向极限（direct limit）}
\expandafter\let\csname BookOriginal@trDirectLimit\endcsname\trDirectLimit
\renewcommand{\trDirectLimit}{\BookTerm{T0737}{\csname BookOriginal@trDirectLimit\endcsname}}
\expandafter\def\csname BookIndex@T0738\endcsname{guinajixian@归纳极限（inductive limit）}
\expandafter\let\csname BookOriginal@trInductiveLimit\endcsname\trInductiveLimit
\renewcommand{\trInductiveLimit}{\BookTerm{T0738}{\csname BookOriginal@trInductiveLimit\endcsname}}
\expandafter\def\csname BookIndex@T0739\endcsname{neishejixian@内射极限（injective limit）}
\expandafter\let\csname BookOriginal@trInjectiveLimit\endcsname\trInjectiveLimit
\renewcommand{\trInjectiveLimit}{\BookTerm{T0739}{\csname BookOriginal@trInjectiveLimit\endcsname}}
\expandafter\def\csname BookIndex@T0740\endcsname{luguo@滤过（filtration /\allowbreak{} filtered）}
\expandafter\let\csname BookOriginal@trFiltered\endcsname\trFiltered
\renewcommand{\trFiltered}{\BookTerm{T0740}{\csname BookOriginal@trFiltered\endcsname}}
\expandafter\def\csname BookIndex@T0741\endcsname{youxiang@有向（directed）}
\expandafter\let\csname BookOriginal@trDirected\endcsname\trDirected
\renewcommand{\trDirected}{\BookTerm{T0741}{\csname BookOriginal@trDirected\endcsname}}
\expandafter\def\csname BookIndex@T0742\endcsname{wanbeihua@完备化（completion）}
\expandafter\let\csname BookOriginal@trCompletion\endcsname\trCompletion
\renewcommand{\trCompletion}{\BookTerm{T0742}{\csname BookOriginal@trCompletion\endcsname}}
\expandafter\def\csname BookIndex@T0743\endcsname{xingshimijishu@形式幂级数（）}
\expandafter\let\csname BookOriginal@trFormalPowerSeries\endcsname\trFormalPowerSeries
\renewcommand{\trFormalPowerSeries}{\BookTerm{T0743}{\csname BookOriginal@trFormalPowerSeries\endcsname}}
\expandafter\def\csname BookIndex@T0744\endcsname{shoulianmijishu@收敛幂级数（convergent power series）}
\expandafter\let\csname BookOriginal@trConvergentPowerSeries\endcsname\trConvergentPowerSeries
\renewcommand{\trConvergentPowerSeries}{\BookTerm{T0744}{\csname BookOriginal@trConvergentPowerSeries\endcsname}}
\expandafter\def\csname BookIndex@T0745\endcsname{pjin@$p$~进（p-adic）}
\expandafter\let\csname BookOriginal@trPadic\endcsname\trPadic
\renewcommand{\trPadic}{\BookTerm{T0745}{\csname BookOriginal@trPadic\endcsname}}
\expandafter\def\csname BookIndex@T0746\endcsname{pjinzhengshu@$p$~进整数（padic integers）}
\expandafter\let\csname BookOriginal@trPadicIntegers\endcsname\trPadicIntegers
\renewcommand{\trPadicIntegers}{\BookTerm{T0746}{\csname BookOriginal@trPadicIntegers\endcsname}}
\expandafter\def\csname BookIndex@T0747\endcsname{pjinshu@$p$~进数（p-adics）}
\expandafter\let\csname BookOriginal@trPadics\endcsname\trPadics
\renewcommand{\trPadics}{\BookTerm{T0747}{\csname BookOriginal@trPadics\endcsname}}
\expandafter\def\csname BookIndex@T0748\endcsname{yushuxuejia@余数学家（comathematician）}
\expandafter\let\csname BookOriginal@trComathematician\endcsname\trComathematician
\renewcommand{\trComathematician}{\BookTerm{T0748}{\csname BookOriginal@trComathematician\endcsname}}
\expandafter\def\csname BookIndex@T0749\endcsname{yuzuozhe@余作者（coauthor）}
\expandafter\let\csname BookOriginal@trCoauthor\endcsname\trCoauthor
\renewcommand{\trCoauthor}{\BookTerm{T0749}{\csname BookOriginal@trCoauthor\endcsname}}
\expandafter\def\csname BookIndex@T0750\endcsname{bansuidui@伴随对（adjoint pair）}
\expandafter\let\csname BookOriginal@trAdjointPair\endcsname\trAdjointPair
\renewcommand{\trAdjointPair}{\BookTerm{T0750}{\csname BookOriginal@trAdjointPair\endcsname}}
\expandafter\def\csname BookIndex@T0751\endcsname{zuobansui@左伴随（left adjoint）}
\expandafter\let\csname BookOriginal@trLeftAdjoint\endcsname\trLeftAdjoint
\renewcommand{\trLeftAdjoint}{\BookTerm{T0751}{\csname BookOriginal@trLeftAdjoint\endcsname}}
\expandafter\def\csname BookIndex@T0752\endcsname{youbansui@右伴随（right adjoint）}
\expandafter\let\csname BookOriginal@trRightAdjoint\endcsname\trRightAdjoint
\renewcommand{\trRightAdjoint}{\BookTerm{T0752}{\csname BookOriginal@trRightAdjoint\endcsname}}
\expandafter\def\csname BookIndex@T0753\endcsname{ziran@自然（natural）}
\expandafter\let\csname BookOriginal@trNatural\endcsname\trNatural
\renewcommand{\trNatural}{\BookTerm{T0753}{\csname BookOriginal@trNatural\endcsname}}
\expandafter\def\csname BookIndex@T0754\endcsname{danwei@单位（unit）}
\expandafter\let\csname BookOriginal@trAdjunctionUnit\endcsname\trAdjunctionUnit
\renewcommand{\trAdjunctionUnit}{\BookTerm{T0754}{\csname BookOriginal@trAdjunctionUnit\endcsname}}
\expandafter\def\csname BookIndex@T0755\endcsname{yudanwei@余单位（counit）}
\expandafter\let\csname BookOriginal@trAdjunctionCounit\endcsname\trAdjunctionCounit
\renewcommand{\trAdjunctionCounit}{\BookTerm{T0755}{\csname BookOriginal@trAdjunctionCounit\endcsname}}
\expandafter\def\csname BookIndex@T0756\endcsname{biaoshi@表示（representation）}
\expandafter\let\csname BookOriginal@trRepresentation\endcsname\trRepresentation
\renewcommand{\trRepresentation}{\BookTerm{T0756}{\csname BookOriginal@trRepresentation\endcsname}}
\expandafter\def\csname BookIndex@T0757\endcsname{biaoshilun@表示论（representation theory）}
\expandafter\let\csname BookOriginal@trRepresentationTheory\endcsname\trRepresentationTheory
\renewcommand{\trRepresentationTheory}{\BookTerm{T0757}{\csname BookOriginal@trRepresentationTheory\endcsname}}
\expandafter\def\csname BookIndex@T0758\endcsname{youxianqun@有限群（finite group）}
\expandafter\let\csname BookOriginal@trFiniteGroup\endcsname\trFiniteGroup
\renewcommand{\trFiniteGroup}{\BookTerm{T0758}{\csname BookOriginal@trFiniteGroup\endcsname}}
\expandafter\def\csname BookIndex@T0759\endcsname{ziqun@子群（subgroup）}
\expandafter\let\csname BookOriginal@trSubgroup\endcsname\trSubgroup
\renewcommand{\trSubgroup}{\BookTerm{T0759}{\csname BookOriginal@trSubgroup\endcsname}}
\expandafter\def\csname BookIndex@T0760\endcsname{youdao@诱导（induction）}
\expandafter\let\csname BookOriginal@trInduction\endcsname\trInduction
\renewcommand{\trInduction}{\BookTerm{T0760}{\csname BookOriginal@trInduction\endcsname}}
\expandafter\def\csname BookIndex@T0761\endcsname{xianzhi@限制（restriction）}
\expandafter\let\csname BookOriginal@trRestriction\endcsname\trRestriction
\renewcommand{\trRestriction}{\BookTerm{T0761}{\csname BookOriginal@trRestriction\endcsname}}
\expandafter\def\csname BookIndex@T0762\endcsname{fuluobeiniwusihufanlu@弗罗贝尼乌斯互反律（Frobenius reciprocity）}
\expandafter\let\csname BookOriginal@trFrobeniusReciprocity\endcsname\trFrobeniusReciprocity
\renewcommand{\trFrobeniusReciprocity}{\BookTerm{T0762}{\csname BookOriginal@trFrobeniusReciprocity\endcsname}}
\expandafter\def\csname BookIndex@T0763\endcsname{weichuihanzi@纬垂函子（suspension functor）}
\expandafter\let\csname BookOriginal@trSuspensionFunctor\endcsname\trSuspensionFunctor
\renewcommand{\trSuspensionFunctor}{\BookTerm{T0763}{\csname BookOriginal@trSuspensionFunctor\endcsname}}
\expandafter\def\csname BookIndex@T0764\endcsname{huanlukongjian@环路空间（loop space）}
\expandafter\let\csname BookOriginal@trLoopSpace\endcsname\trLoopSpace
\renewcommand{\trLoopSpace}{\BookTerm{T0764}{\csname BookOriginal@trLoopSpace\endcsname}}
\expandafter\def\csname BookIndex@T0765\endcsname{banqun@半群（semigroup）}
\expandafter\let\csname BookOriginal@trSemigroup\endcsname\trSemigroup
\renewcommand{\trSemigroup}{\BookTerm{T0765}{\csname BookOriginal@trSemigroup\endcsname}}
\expandafter\def\csname BookIndex@T0766\endcsname{abeier@阿贝尔（Abelian）}
\expandafter\let\csname BookOriginal@trAbelian\endcsname\trAbelian
\renewcommand{\trAbelian}{\BookTerm{T0766}{\csname BookOriginal@trAbelian\endcsname}}
\expandafter\def\csname BookIndex@T0767\endcsname{abeierbanqun@阿贝尔半群（abelian semigroup）}
\expandafter\let\csname BookOriginal@trAbelianSemigroup\endcsname\trAbelianSemigroup
\renewcommand{\trAbelianSemigroup}{\BookTerm{T0767}{\csname BookOriginal@trAbelianSemigroup\endcsname}}
\expandafter\def\csname BookIndex@T0768\endcsname{qunhua@群化（groupification）}
\expandafter\let\csname BookOriginal@trGroupification\endcsname\trGroupification
\renewcommand{\trGroupification}{\BookTerm{T0768}{\csname BookOriginal@trGroupification\endcsname}}
\expandafter\def\csname BookIndex@T0769\endcsname{niyuan@逆元（inverse）}
\expandafter\let\csname BookOriginal@trInverseElement\endcsname\trInverseElement
\renewcommand{\trInverseElement}{\BookTerm{T0769}{\csname BookOriginal@trInverseElement\endcsname}}
\expandafter\def\csname BookIndex@T0770\endcsname{biaoliangkuozhang@标量扩张（extension of scalars）}
\expandafter\let\csname BookOriginal@trExtensionOfScalars\endcsname\trExtensionOfScalars
\renewcommand{\trExtensionOfScalars}{\BookTerm{T0770}{\csname BookOriginal@trExtensionOfScalars\endcsname}}
\expandafter\def\csname BookIndex@T0771\endcsname{biaoliangxianzhi@标量限制（restriction of scalars）}
\expandafter\let\csname BookOriginal@trRestrictionOfScalars\endcsname\trRestrictionOfScalars
\renewcommand{\trRestrictionOfScalars}{\BookTerm{T0771}{\csname BookOriginal@trRestrictionOfScalars\endcsname}}
\expandafter\def\csname BookIndex@T0772\endcsname{yiwang@遗忘（forgetful）}
\expandafter\let\csname BookOriginal@trForgetful\endcsname\trForgetful
\renewcommand{\trForgetful}{\BookTerm{T0772}{\csname BookOriginal@trForgetful\endcsname}}
\expandafter\def\csname BookIndex@T0773\endcsname{jiaxing@加性（additive）}
\expandafter\let\csname BookOriginal@trAdditive\endcsname\trAdditive
\renewcommand{\trAdditive}{\BookTerm{T0773}{\csname BookOriginal@trAdditive\endcsname}}
\expandafter\def\csname BookIndex@T0774\endcsname{jiaxingfanchou@加性范畴（additive category）}
\expandafter\let\csname BookOriginal@trAdditiveCategory\endcsname\trAdditiveCategory
\renewcommand{\trAdditiveCategory}{\BookTerm{T0774}{\csname BookOriginal@trAdditiveCategory\endcsname}}
\expandafter\def\csname BookIndex@T0775\endcsname{jiaxinghanzi@加性函子（additive functor）}
\expandafter\let\csname BookOriginal@trAdditiveFunctor\endcsname\trAdditiveFunctor
\renewcommand{\trAdditiveFunctor}{\BookTerm{T0775}{\csname BookOriginal@trAdditiveFunctor\endcsname}}
\expandafter\def\csname BookIndex@T0776\endcsname{zhenghe@正合（exact）}
\expandafter\let\csname BookOriginal@trExact\endcsname\trExact
\renewcommand{\trExact}{\BookTerm{T0776}{\csname BookOriginal@trExact\endcsname}}
\expandafter\def\csname BookIndex@T0777\endcsname{lingtaishe@零态射（0-morphism）}
\expandafter\let\csname BookOriginal@trZeroMorphism\endcsname\trZeroMorphism
\renewcommand{\trZeroMorphism}{\BookTerm{T0777}{\csname BookOriginal@trZeroMorphism\endcsname}}
\expandafter\def\csname BookIndex@T0778\endcsname{ziduixiang@子对象（subobject）}
\expandafter\let\csname BookOriginal@trSubobject\endcsname\trSubobject
\renewcommand{\trSubobject}{\BookTerm{T0778}{\csname BookOriginal@trSubobject\endcsname}}
\expandafter\def\csname BookIndex@T0779\endcsname{shangduixiang@商对象（quotient object）}
\expandafter\let\csname BookOriginal@trQuotientObject\endcsname\trQuotientObject
\renewcommand{\trQuotientObject}{\BookTerm{T0779}{\csname BookOriginal@trQuotientObject\endcsname}}
\expandafter\def\csname BookIndex@T0780\endcsname{xiang@像（image）}
\expandafter\let\csname BookOriginal@trImage\endcsname\trImage
\renewcommand{\trImage}{\BookTerm{T0780}{\csname BookOriginal@trImage\endcsname}}
\expandafter\def\csname BookIndex@T0781\endcsname{fuxing@复形（complex）}
\expandafter\let\csname BookOriginal@trComplexN\endcsname\trComplexN
\renewcommand{\trComplexN}{\BookTerm{T0781}{\csname BookOriginal@trComplexN\endcsname}}
\expandafter\def\csname BookIndex@T0782\endcsname{tonglunde@同伦的（homotopic）}
\expandafter\let\csname BookOriginal@trHomotopic\endcsname\trHomotopic
\renewcommand{\trHomotopic}{\BookTerm{T0782}{\csname BookOriginal@trHomotopic\endcsname}}
\expandafter\def\csname BookIndex@T0783\endcsname{lianjietongtai@连接同态（connecting homomorphism）}
\expandafter\let\csname BookOriginal@trConnectingHomomorphism\endcsname\trConnectingHomomorphism
\renewcommand{\trConnectingHomomorphism}{\BookTerm{T0783}{\csname BookOriginal@trConnectingHomomorphism\endcsname}}
\expandafter\def\csname BookIndex@T0784\endcsname{bilian@闭链（cycle）}
\expandafter\let\csname BookOriginal@trCycle\endcsname\trCycle
\renewcommand{\trCycle}{\BookTerm{T0784}{\csname BookOriginal@trCycle\endcsname}}
\expandafter\def\csname BookIndex@T0785\endcsname{bianyuanlian@边缘链（boundary）}
\expandafter\let\csname BookOriginal@trBoundary\endcsname\trBoundary
\renewcommand{\trBoundary}{\BookTerm{T0785}{\csname BookOriginal@trBoundary\endcsname}}
\expandafter\def\csname BookIndex@T0786\endcsname{zhangzhenghelie@长正合列（long exact sequence）}
\expandafter\let\csname BookOriginal@trLongExactSequence\endcsname\trLongExactSequence
\renewcommand{\trLongExactSequence}{\BookTerm{T0786}{\csname BookOriginal@trLongExactSequence\endcsname}}
\expandafter\def\csname BookIndex@T0787\endcsname{duanzhenghelie@短正合列（short exact sequence）}
\expandafter\let\csname BookOriginal@trShortExactSequence\endcsname\trShortExactSequence
\renewcommand{\trShortExactSequence}{\BookTerm{T0787}{\csname BookOriginal@trShortExactSequence\endcsname}}
\expandafter\def\csname BookIndex@T0788\endcsname{pingtanxing@平坦性（flatness）}
\expandafter\let\csname BookOriginal@trFlatness\endcsname\trFlatness
\renewcommand{\trFlatness}{\BookTerm{T0788}{\csname BookOriginal@trFlatness\endcsname}}
\expandafter\def\csname BookIndex@T0789\endcsname{jiaohuanyaobanqun@交换幺半群（commutative monoid）}
\expandafter\let\csname BookOriginal@trCommutativeMonoid\endcsname\trCommutativeMonoid
\renewcommand{\trCommutativeMonoid}{\BookTerm{T0789}{\csname BookOriginal@trCommutativeMonoid\endcsname}}
\expandafter\def\csname BookIndex@T0790\endcsname{yuduijiao@余对角（codiagonal）}
\expandafter\let\csname BookOriginal@trCodiagonal\endcsname\trCodiagonal
\renewcommand{\trCodiagonal}{\BookTerm{T0790}{\csname BookOriginal@trCodiagonal\endcsname}}
\expandafter\def\csname BookIndex@T0791\endcsname{shengchengyuan@生成元（generator）}
\expandafter\let\csname BookOriginal@trGenerator\endcsname\trGenerator
\renewcommand{\trGenerator}{\BookTerm{T0791}{\csname BookOriginal@trGenerator\endcsname}}
\expandafter\def\csname BookIndex@T0792\endcsname{guanxi@关系（relation）}
\expandafter\let\csname BookOriginal@trRelation\endcsname\trRelation
\renewcommand{\trRelation}{\BookTerm{T0792}{\csname BookOriginal@trRelation\endcsname}}
\expandafter\def\csname BookIndex@T0793\endcsname{banahekongjian@巴拿赫空间（Banach space）}
\expandafter\let\csname BookOriginal@trBanachSpace\endcsname\trBanachSpace
\renewcommand{\trBanachSpace}{\BookTerm{T0793}{\csname BookOriginal@trBanachSpace\endcsname}}
\expandafter\def\csname BookIndex@T0794\endcsname{fuleidemiqieerqianrudingli@弗雷德--米切尔嵌入定理（freyd mitchell embedding theorem）}
\expandafter\let\csname BookOriginal@trFreydMitchellEmbeddingTheorem\endcsname\trFreydMitchellEmbeddingTheorem
\renewcommand{\trFreydMitchellEmbeddingTheorem}{\BookTerm{T0794}{\csname BookOriginal@trFreydMitchellEmbeddingTheorem\endcsname}}
\expandafter\def\csname BookIndex@T0795\endcsname{nixiangxitong@逆向系统（inverse system）}
\expandafter\let\csname BookOriginal@trInverseSystem\endcsname\trInverseSystem
\renewcommand{\trInverseSystem}{\BookTerm{T0795}{\csname BookOriginal@trInverseSystem\endcsname}}
\expandafter\def\csname BookIndex@T0796\endcsname{yuluguo@余滤过（cofiltered）}
\expandafter\let\csname BookOriginal@trCofiltered\endcsname\trCofiltered
\renewcommand{\trCofiltered}{\BookTerm{T0796}{\csname BookOriginal@trCofiltered\endcsname}}
\expandafter\def\csname BookIndex@T0797\endcsname{guoduyingshe@过渡映射（transition map）}
\expandafter\let\csname BookOriginal@trTransitionMap\endcsname\trTransitionMap
\renewcommand{\trTransitionMap}{\BookTerm{T0797}{\csname BookOriginal@trTransitionMap\endcsname}}
\expandafter\def\csname BookIndex@T0798\endcsname{jinghua@茎化（stalkification）}
\expandafter\let\csname BookOriginal@trStalkification\endcsname\trStalkification
\renewcommand{\trStalkification}{\BookTerm{T0798}{\csname BookOriginal@trStalkification\endcsname}}
\expandafter\def\csname BookIndex@T0799\endcsname{shuangfuxing@双复形（double complex）}
\expandafter\let\csname BookOriginal@trDoubleComplex\endcsname\trDoubleComplex
\renewcommand{\trDoubleComplex}{\BookTerm{T0799}{\csname BookOriginal@trDoubleComplex\endcsname}}
\expandafter\def\csname BookIndex@T0800\endcsname{quanfuxing@全复形（total complex）}
\expandafter\let\csname BookOriginal@trTotalComplex\endcsname\trTotalComplex
\renewcommand{\trTotalComplex}{\BookTerm{T0800}{\csname BookOriginal@trTotalComplex\endcsname}}
\expandafter\def\csname BookIndex@T0801\endcsname{fanduijiaoxian@反对角线（antidiagonal）}
\expandafter\let\csname BookOriginal@trAntidiagonal\endcsname\trAntidiagonal
\renewcommand{\trAntidiagonal}{\BookTerm{T0801}{\csname BookOriginal@trAntidiagonal\endcsname}}
\expandafter\def\csname BookIndex@T0802\endcsname{chaoshangtongdiao@超上同调（hypercohomology）}
\expandafter\let\csname BookOriginal@trHypercohomology\endcsname\trHypercohomology
\renewcommand{\trHypercohomology}{\BookTerm{T0802}{\csname BookOriginal@trHypercohomology\endcsname}}
\expandafter\def\csname BookIndex@T0803\endcsname{ye@页（page）}
\expandafter\let\csname BookOriginal@trPage\endcsname\trPage
\renewcommand{\trPage}{\BookTerm{T0803}{\csname BookOriginal@trPage\endcsname}}
\expandafter\def\csname BookIndex@T0804\endcsname{diyixiangxian@第一象限（first quadrant）}
\expandafter\let\csname BookOriginal@trFirstQuadrant\endcsname\trFirstQuadrant
\renewcommand{\trFirstQuadrant}{\BookTerm{T0804}{\csname BookOriginal@trFirstQuadrant\endcsname}}
\expandafter\def\csname BookIndex@T0805\endcsname{fanjiaohuan@反交换（anticommute）}
\expandafter\let\csname BookOriginal@trAnticommute\endcsname\trAnticommute
\renewcommand{\trAnticommute}{\BookTerm{T0805}{\csname BookOriginal@trAnticommute\endcsname}}
\expandafter\def\csname BookIndex@T0806\endcsname{dingxiang@定向（orientation）}
\expandafter\let\csname BookOriginal@trOrientation\endcsname\trOrientation
\renewcommand{\trOrientation}{\BookTerm{T0806}{\csname BookOriginal@trOrientation\endcsname}}
\expandafter\def\csname BookIndex@T0807\endcsname{zishang@子商（subquotient）}
\expandafter\let\csname BookOriginal@trSubquotient\endcsname\trSubquotient
\renewcommand{\trSubquotient}{\BookTerm{T0807}{\csname BookOriginal@trSubquotient\endcsname}}
\expandafter\def\csname BookIndex@T0808\endcsname{shoulian@收敛（converge）}
\expandafter\let\csname BookOriginal@trConverge\endcsname\trConverge
\renewcommand{\trConverge}{\BookTerm{T0808}{\csname BookOriginal@trConverge\endcsname}}
\expandafter\def\csname BookIndex@T0809\endcsname{dida@抵达（abut）}
\expandafter\let\csname BookOriginal@trAbut\endcsname\trAbut
\renewcommand{\trAbut}{\BookTerm{T0809}{\csname BookOriginal@trAbut\endcsname}}
\expandafter\def\csname BookIndex@T0810\endcsname{tasuo@塌缩（collapse）}
\expandafter\let\csname BookOriginal@trCollapse\endcsname\trCollapse
\renewcommand{\trCollapse}{\BookTerm{T0810}{\csname BookOriginal@trCollapse\endcsname}}
\expandafter\def\csname BookIndex@T0811\endcsname{yingshezhui@映射锥（mapping cone）}
\expandafter\let\csname BookOriginal@trMappingCone\endcsname\trMappingCone
\renewcommand{\trMappingCone}{\BookTerm{T0811}{\csname BookOriginal@trMappingCone\endcsname}}
\expandafter\def\csname BookIndex@T0812\endcsname{puhanzi@谱函子（spectral functor）}
\expandafter\let\csname BookOriginal@trSpectralFunctor\endcsname\trSpectralFunctor
\renewcommand{\trSpectralFunctor}{\BookTerm{T0812}{\csname BookOriginal@trSpectralFunctor\endcsname}}
\expandafter\def\csname BookIndex@T0813\endcsname{qie@切（tangent）}
\expandafter\let\csname BookOriginal@trTangent\endcsname\trTangent
\renewcommand{\trTangent}{\BookTerm{T0813}{\csname BookOriginal@trTangent\endcsname}}
\expandafter\def\csname BookIndex@T0814\endcsname{quqidian@去奇点（desingularization）}
\expandafter\let\csname BookOriginal@trDesingularization\endcsname\trDesingularization
\renewcommand{\trDesingularization}{\BookTerm{T0814}{\csname BookOriginal@trDesingularization\endcsname}}
\expandafter\def\csname BookIndex@T0815\endcsname{qidianjiexiao@奇点解消（resolution of singularities）}
\expandafter\let\csname BookOriginal@trResolutionOfSingularities\endcsname\trResolutionOfSingularities
\renewcommand{\trResolutionOfSingularities}{\BookTerm{T0815}{\csname BookOriginal@trResolutionOfSingularities\endcsname}}
\expandafter\def\csname BookIndex@T0816\endcsname{quanbianhuan@全变换（total transform）}
\expandafter\let\csname BookOriginal@trTotalTransform\endcsname\trTotalTransform
\renewcommand{\trTotalTransform}{\BookTerm{T0816}{\csname BookOriginal@trTotalTransform\endcsname}}
\expandafter\def\csname BookIndex@T0817\endcsname{zhongxin@中心（center）}
\expandafter\let\csname BookOriginal@trCenter\endcsname\trCenter
\renewcommand{\trCenter}{\BookTerm{T0817}{\csname BookOriginal@trCenter\endcsname}}
\expandafter\def\csname BookIndex@T0818\endcsname{xianweilifangti@纤维立方体（fibered cube）}
\expandafter\let\csname BookOriginal@trFiberedCube\endcsname\trFiberedCube
\renewcommand{\trFiberedCube}{\BookTerm{T0818}{\csname BookOriginal@trFiberedCube\endcsname}}
\expandafter\def\csname BookIndex@T0819\endcsname{jixiaomoxinggangling@极小模型纲领（minimal model program）}
\expandafter\let\csname BookOriginal@trMinimalModelProgram\endcsname\trMinimalModelProgram
\renewcommand{\trMinimalModelProgram}{\BookTerm{T0819}{\csname BookOriginal@trMinimalModelProgram\endcsname}}
\expandafter\def\csname BookIndex@T0820\endcsname{zhangkaibibaoyinli@胀开闭包引理（Blow-up Closure Lemma）}
\expandafter\let\csname BookOriginal@trBlowupClosureLemma\endcsname\trBlowupClosureLemma
\renewcommand{\trBlowupClosureLemma}{\BookTerm{T0820}{\csname BookOriginal@trBlowupClosureLemma\endcsname}}
\expandafter\def\csname BookIndex@T0821\endcsname{yushidingli@余式定理（Remainder Theorem）}
\expandafter\let\csname BookOriginal@trRemainderTheorem\endcsname\trRemainderTheorem
\renewcommand{\trRemainderTheorem}{\BookTerm{T0821}{\csname BookOriginal@trRemainderTheorem\endcsname}}
\expandafter\def\csname BookIndex@T0822\endcsname{fencidaishu@分次代数（graded algebra）}
\expandafter\let\csname BookOriginal@trGradedAlgebra\endcsname\trGradedAlgebra
\renewcommand{\trGradedAlgebra}{\BookTerm{T0822}{\csname BookOriginal@trGradedAlgebra\endcsname}}
\expandafter\def\csname BookIndex@T0823\endcsname{dengjintu@邓金图（Dynkin diagram）}
\expandafter\let\csname BookOriginal@trDynkinDiagram\endcsname\trDynkinDiagram
\renewcommand{\trDynkinDiagram}{\BookTerm{T0823}{\csname BookOriginal@trDynkinDiagram\endcsname}}
\expandafter\def\csname BookIndex@T0824\endcsname{biangengdingli@变更定理（Alteration Theorem）}
\expandafter\let\csname BookOriginal@trAlterationTheorem\endcsname\trAlterationTheorem
\renewcommand{\trAlterationTheorem}{\BookTerm{T0824}{\csname BookOriginal@trAlterationTheorem\endcsname}}
\expandafter\def\csname BookIndex@T0825\endcsname{ruofenjiedingli@弱分解定理（Weak Factorization Theorem）}
\expandafter\let\csname BookOriginal@trWeakFactorizationTheorem\endcsname\trWeakFactorizationTheorem
\renewcommand{\trWeakFactorizationTheorem}{\BookTerm{T0825}{\csname BookOriginal@trWeakFactorizationTheorem\endcsname}}
\expandafter\def\csname BookIndex@T0826\endcsname{adiyafanzhuan@阿蒂亚翻转（atiyah flop）}
\expandafter\let\csname BookOriginal@trAtiyahFlop\endcsname\trAtiyahFlop
\renewcommand{\trAtiyahFlop}{\BookTerm{T0826}{\csname BookOriginal@trAtiyahFlop\endcsname}}
\expandafter\def\csname BookIndex@T0827\endcsname{tuoyuanxianweihua@椭圆纤维化（elliptic fibration）}
\expandafter\let\csname BookOriginal@trEllipticFibration\endcsname\trEllipticFibration
\renewcommand{\trEllipticFibration}{\BookTerm{T0827}{\csname BookOriginal@trEllipticFibration\endcsname}}
\expandafter\def\csname BookIndex@T0828\endcsname{xiangjiaotongdiao@相交同调（intersection homology）}
\expandafter\let\csname BookOriginal@trIntersectionHomology\endcsname\trIntersectionHomology
\renewcommand{\trIntersectionHomology}{\BookTerm{T0828}{\csname BookOriginal@trIntersectionHomology\endcsname}}
\expandafter\def\csname BookIndex@T0829\endcsname{daozhengguizhuidexingbian@到正规锥的形变（deformation to the normal cone）}
\expandafter\let\csname BookOriginal@trDeformationToTheNormalCone\endcsname\trDeformationToTheNormalCone
\renewcommand{\trDeformationToTheNormalCone}{\BookTerm{T0829}{\csname BookOriginal@trDeformationToTheNormalCone\endcsname}}
\expandafter\def\csname BookIndex@T0830\endcsname{jiexitonggou@解析同构（analytically isomorphic）}
\expandafter\let\csname BookOriginal@trAnalyticallyIsomorphic\endcsname\trAnalyticallyIsomorphic
\renewcommand{\trAnalyticallyIsomorphic}{\BookTerm{T0830}{\csname BookOriginal@trAnalyticallyIsomorphic\endcsname}}
\expandafter\def\csname BookIndex@T0831\endcsname{bubianlianglun@不变量论（invariant theory）}
\expandafter\let\csname BookOriginal@trInvariantTheory\endcsname\trInvariantTheory
\renewcommand{\trInvariantTheory}{\BookTerm{T0831}{\csname BookOriginal@trInvariantTheory\endcsname}}
\expandafter\def\csname BookIndex@T0832\endcsname{qiyiguiji@奇异轨迹（singular locus）}
\expandafter\let\csname BookOriginal@trSingularLocus\endcsname\trSingularLocus
\renewcommand{\trSingularLocus}{\BookTerm{T0832}{\csname BookOriginal@trSingularLocus\endcsname}}
\expandafter\def\csname BookIndex@T0833\endcsname{deltahanzi@$\delta$-函子（delta-functor）}
\expandafter\let\csname BookOriginal@trDeltaFunctor\endcsname\trDeltaFunctor
\renewcommand{\trDeltaFunctor}{\BookTerm{T0833}{\csname BookOriginal@trDeltaFunctor\endcsname}}
\expandafter\def\csname BookIndex@T0834\endcsname{kexiaoqu@可消去（effaceable）}
\expandafter\let\csname BookOriginal@trEffaceable\endcsname\trEffaceable
\renewcommand{\trEffaceable}{\BookTerm{T0834}{\csname BookOriginal@trEffaceable\endcsname}}
\expandafter\def\csname BookIndex@T0835\endcsname{lingdiao@零调（acyclic）}
\expandafter\let\csname BookOriginal@trAcyclic\endcsname\trAcyclic
\renewcommand{\trAcyclic}{\BookTerm{T0835}{\csname BookOriginal@trAcyclic\endcsname}}
\expandafter\def\csname BookIndex@T0836\endcsname{songruan@松软（flasque）}
\expandafter\let\csname BookOriginal@trFlasque\endcsname\trFlasque
\renewcommand{\trFlasque}{\BookTerm{T0836}{\csname BookOriginal@trFlasque\endcsname}}
\expandafter\def\csname BookIndex@T0837\endcsname{xingbian@形变（deformation）}
\expandafter\let\csname BookOriginal@trDeformation\endcsname\trDeformation
\renewcommand{\trDeformation}{\BookTerm{T0837}{\csname BookOriginal@trDeformation\endcsname}}
\expandafter\def\csname BookIndex@T0838\endcsname{xingbianlilun@形变理论（deformation theory）}
\expandafter\let\csname BookOriginal@trDeformationTheory\endcsname\trDeformationTheory
\renewcommand{\trDeformationTheory}{\BookTerm{T0838}{\csname BookOriginal@trDeformationTheory\endcsname}}
\expandafter\def\csname BookIndex@T0839\endcsname{kaiyingshe@开映射（open map）}
\expandafter\let\csname BookOriginal@trOpenMap\endcsname\trOpenMap
\renewcommand{\trOpenMap}{\BookTerm{T0839}{\csname BookOriginal@trOpenMap\endcsname}}
\expandafter\def\csname BookIndex@T0840\endcsname{xierbotegaixing@希尔伯特概形（Hilbert scheme）}
\expandafter\let\csname BookOriginal@trHilbertScheme\endcsname\trHilbertScheme
\renewcommand{\trHilbertScheme}{\BookTerm{T0840}{\csname BookOriginal@trHilbertScheme\endcsname}}
\expandafter\def\csname BookIndex@T0841\endcsname{xierboteduoxiangshi@希尔伯特多项式（Hilbert polynomial）}
\expandafter\let\csname BookOriginal@trHilbertPolynomial\endcsname\trHilbertPolynomial
\renewcommand{\trHilbertPolynomial}{\BookTerm{T0841}{\csname BookOriginal@trHilbertPolynomial\endcsname}}
\expandafter\def\csname BookIndex@T0842\endcsname{banzhenghe@半正合（half-exact）}
\expandafter\let\csname BookOriginal@trHalfExact\endcsname\trHalfExact
\renewcommand{\trHalfExact}{\BookTerm{T0842}{\csname BookOriginal@trHalfExact\endcsname}}
\expandafter\def\csname BookIndex@T0843\endcsname{qijipingtan@奇迹平坦（Miracle flatness）}
\expandafter\let\csname BookOriginal@trMiracleFlat\endcsname\trMiracleFlat
\renewcommand{\trMiracleFlat}{\BookTerm{T0843}{\csname BookOriginal@trMiracleFlat\endcsname}}
\expandafter\def\csname BookIndex@T0844\endcsname{sanzhongdian@三重点（triple point）}
\expandafter\let\csname BookOriginal@trTriplePoint\endcsname\trTriplePoint
\renewcommand{\trTriplePoint}{\BookTerm{T0844}{\csname BookOriginal@trTriplePoint\endcsname}}
\expandafter\def\csname BookIndex@T0845\endcsname{fenliejiedian@分裂结点（split node）}
\expandafter\let\csname BookOriginal@trSplitNode\endcsname\trSplitNode
\renewcommand{\trSplitNode}{\BookTerm{T0845}{\csname BookOriginal@trSplitNode\endcsname}}
\expandafter\def\csname BookIndex@T0846\endcsname{feifenliejiedian@非分裂结点（non-split node）}
\expandafter\let\csname BookOriginal@trNonSplitNode\endcsname\trNonSplitNode
\renewcommand{\trNonSplitNode}{\BookTerm{T0846}{\csname BookOriginal@trNonSplitNode\endcsname}}
\expandafter\def\csname BookIndex@T0847\endcsname{jiahou@加厚（thickening）}
\expandafter\let\csname BookOriginal@trThickening\endcsname\trThickening
\renewcommand{\trThickening}{\BookTerm{T0847}{\csname BookOriginal@trThickening\endcsname}}
\expandafter\def\csname BookIndex@T0848\endcsname{deltabubianliang@$\delta$-不变量（delta-invariant）}
\expandafter\let\csname BookOriginal@trDeltaInvariant\endcsname\trDeltaInvariant
\renewcommand{\trDeltaInvariant}{\BookTerm{T0848}{\csname BookOriginal@trDeltaInvariant\endcsname}}
\expandafter\def\csname BookIndex@T0849\endcsname{taishe@态射（morphism）}
\expandafter\let\csname BookOriginal@trMorphisms\endcsname\trMorphisms
\renewcommand{\trMorphisms}{\BookTerm{T0849}{\csname BookOriginal@trMorphisms\endcsname}}
\expandafter\def\csname BookIndex@T0850\endcsname{jihe@集合（set）}
\expandafter\let\csname BookOriginal@trSetsWord\endcsname\trSetsWord
\renewcommand{\trSetsWord}{\BookTerm{T0850}{\csname BookOriginal@trSetsWord\endcsname}}
\expandafter\def\csname BookIndex@T0851\endcsname{ceng@层（sheaf）}
\expandafter\let\csname BookOriginal@trSheaves\endcsname\trSheaves
\renewcommand{\trSheaves}{\BookTerm{T0851}{\csname BookOriginal@trSheaves\endcsname}}
\expandafter\def\csname BookIndex@T0852\endcsname{nuote@诺特（noetherian）}
\expandafter\let\csname BookOriginal@trNoetherian\endcsname\trNoetherian
\renewcommand{\trNoetherian}{\BookTerm{T0852}{\csname BookOriginal@trNoetherian\endcsname}}
\expandafter\def\csname BookIndex@T0853\endcsname{ating@阿廷（artinian）}
\expandafter\let\csname BookOriginal@trArtinian\endcsname\trArtinian
\renewcommand{\trArtinian}{\BookTerm{T0853}{\csname BookOriginal@trArtinian\endcsname}}
\expandafter\def\csname BookIndex@T0854\endcsname{danshe@单射（injective）}
\expandafter\let\csname BookOriginal@trInjection\endcsname\trInjection
\renewcommand{\trInjection}{\BookTerm{T0854}{\csname BookOriginal@trInjection\endcsname}}
\expandafter\def\csname BookIndex@T0855\endcsname{fuxing@复形（complex）}
\expandafter\let\csname BookOriginal@trComplexes\endcsname\trComplexes
\renewcommand{\trComplexes}{\BookTerm{T0855}{\csname BookOriginal@trComplexes\endcsname}}
\expandafter\def\csname BookIndex@T0856\endcsname{bilian@闭链（cycle）}
\expandafter\let\csname BookOriginal@trCycles\endcsname\trCycles
\renewcommand{\trCycles}{\BookTerm{T0856}{\csname BookOriginal@trCycles\endcsname}}
\expandafter\def\csname BookIndex@T0857\endcsname{bianyuanlian@边缘链（boundary）}
\expandafter\let\csname BookOriginal@trBoundaries\endcsname\trBoundaries
\renewcommand{\trBoundaries}{\BookTerm{T0857}{\csname BookOriginal@trBoundaries\endcsname}}
\expandafter\def\csname BookIndex@T0858\endcsname{shengchengyuan@生成元（generator）}
\expandafter\let\csname BookOriginal@trGenerators\endcsname\trGenerators
\renewcommand{\trGenerators}{\BookTerm{T0858}{\csname BookOriginal@trGenerators\endcsname}}
\expandafter\def\csname BookIndex@T0859\endcsname{dingli@定理（Theorem）}
\expandafter\def\csname BookIndex@T0860\endcsname{mingti@命题（Proposition）}
\expandafter\def\csname BookIndex@T0861\endcsname{yinli@引理（Lemma）}
\expandafter\def\csname BookIndex@T0862\endcsname{caixiang@猜想（Conjecture）}
\expandafter\def\csname BookIndex@T0863\endcsname{dingyi@定义（Definition）}
\expandafter\def\csname BookIndex@T0864\endcsname{li@例（Example）}
\expandafter\def\csname BookIndex@T0865\endcsname{lianxi@练习（Exercise）}
\expandafter\def\csname BookIndex@T0866\endcsname{jihao@记号（Notation）}
\expandafter\def\csname BookIndex@T0867\endcsname{yueding@约定（Convention）}
\expandafter\def\csname BookIndex@T0868\endcsname{zhuji@注记（Remark）}
\expandafter\def\csname BookIndex@T0869\endcsname{jinggao@警告（Warning）}
\expandafter\def\csname BookIndex@T0870\endcsname{jingshi@警示（Caveat）}
\expandafter\def\csname BookIndex@T0871\endcsname{zhengzehanshu@正则函数（regular function）}
\expandafter\def\csname BookIndex@T0872\endcsname{trveroneseqianru@\trVeronese{}嵌入（Veronese embedding）}
\expandafter\def\csname BookIndex@T0873\endcsname{qungaixingzuoyong@群概形作用（group scheme action）}
\expandafter\def\csname BookIndex@T0874\endcsname{youxiaotrcartierchuzi@有效\trCartier{}除子（effective Cartier divisor）}
\expandafter\def\csname BookIndex@T0875\endcsname{genxing@根性（radicial）}
\expandafter\def\csname BookIndex@T0876\endcsname{shidian@实点（real point）}
\expandafter\def\csname BookIndex@T0877\endcsname{tezhengbiao@特征标（character）}
\expandafter\def\csname BookIndex@T0878\endcsname{zhengguibi@正规闭（normally closed）}
\expandafter\def\csname BookIndex@T0879\endcsname{duijiaozigaixing@对角子概形（diagonal）}
\expandafter\def\csname BookIndex@T0880\endcsname{trzariskitapu@\trZariski{}拓扑（Zariski topology）}
\expandafter\def\csname BookIndex@T0881\endcsname{trdedekindzhenghuan@\trDedekind{}整环（Dedekind domain）}
\expandafter\def\csname BookIndex@T0882\endcsname{duozhongcishu@多重次数（multidegree）}
\expandafter\def\csname BookIndex@T0883\endcsname{yuehuatongdiao@约化同调（reduced homology）}
\expandafter\def\csname BookIndex@T0884\endcsname{danxing@单形（simplex）}
\expandafter\def\csname BookIndex@T0885\endcsname{jihezhengze@几何正则（geometrically regular）}
\expandafter\def\csname BookIndex@T0886\endcsname{lunduan@论断（Claim）}
\expandafter\def\csname BookIndex@T0887\endcsname{trcartierchuzi@\trCartier{}除子（Cartier divisor）}
\expandafter\def\csname BookIndex@T0888\endcsname{trweilchuzi@\trWeil{}除子（Weil divisor）}
\expandafter\def\csname BookIndex@T0889\endcsname{trcechshangtongdiao@\trCech{}上同调（Cech cohomology）}
\expandafter\def\csname BookIndex@T0890\endcsname{trserreduiou@\trSerre{}对偶（Serre duality）}
\expandafter\def\csname BookIndex@T0891\endcsname{trriemannqumian@\trRiemann{}曲面（Riemann surface）}
\expandafter\def\csname BookIndex@T0892\endcsname{tryonedaqianru@\trYoneda{}嵌入（Yoneda embedding）}
\expandafter\def\csname BookIndex@T0893\endcsname{trhilbertduoxiangshi@\trHilbert{}多项式（Hilbert polynomial）}
\makeatother

\let\BookOriginalIndex\index
\renewcommand{\index}[1]{\ifBookIndexing\BookOriginalIndex{#1}\fi}
